% Bachelorarbeit – KOMA-Script Vorlage (DE)
\documentclass[12pt,a4paper,oneside,bibliography=totoc,listof=totoc,shortlabels]{scrreprt}

% Sprache, Kodierung, Typografie
\usepackage[T1]{fontenc}
% Sicherstellen, dass standardmä\ ssig Serifenschrift genutzt wird
\renewcommand{\familydefault}{\rmdefault}
\usepackage[ngerman]{babel}
\usepackage{csquotes}
\usepackage{microtype}
\usepackage[document]{ragged2e}  % Bessere Hyphenation-Kontrolle für problematische Absätze

% Automatic formatting improvements
% Hyphenation hints für lange deutsche/englische Wörter
\hyphenation{
  Vir-tu-a-li-sie-rung
  Kom-mu-ni-ka-tions-sys-te-me
  Sprach-ver-ar-bei-tung
  Netz-werk-ver-bin-dung
  Im-ple-men-tie-rung
  Tele-fo-nie-sys-tem
  Spra-cher-ken-nung
  Text-zu-Spra-che
}
% End of automatic improvements

% Layout & Links
\usepackage{geometry}
\geometry{a4paper, left=30mm, right=25mm, top=25mm, bottom=30mm}
\usepackage{graphicx}
\usepackage{xcolor}
\usepackage{booktabs}
\usepackage{siunitx}
\usepackage{caption}
\usepackage{subcaption}
\usepackage{listings}
\usepackage{longtable}
\usepackage{float}
\usepackage{placeins}
% Zusatz für Tabellenlayouts in Kapitel 4
\usepackage{array}
\usepackage{tabularx}

% Globale Caption-Position: Tabellen oben, Abbildungen unten
\captionsetup[table]{position=top}
\captionsetup[figure]{position=bottom}

% KOMA: Serifenschrift für Überschriften und standardisierte Captions
\setkomafont{disposition}{\normalfont\bfseries\sffamily}
\setkomafont{caption}{\normalfont\small}
\setkomafont{captionlabel}{\bfseries}

% Zusätzliche KOMA-Optimierungen
\KOMAoptions{
  toc=listof,
  listof=totoc,
  captions=tableheading
}

% Code-Snippets
\lstset{
  basicstyle=\ttfamily\small,
  columns=fullflexible,
  breaklines=true,
  keepspaces=true,
  showstringspaces=false,
  upquote=true,
  frame=lines,
  aboveskip=12pt,
  belowskip=12pt,
  xleftmargin=10pt,
  numbers=none
}

% Definition von YAML Syntax-Highlighting
\lstdefinelanguage{yaml}{
  keywords={true,false,null,y,n},
  keywordstyle=\color{blue}\bfseries,
  basicstyle=\ttfamily\small,
  sensitive=false,
  comment=[l]{\#},
  commentstyle=\color{gray}\ttfamily,
  stringstyle=\color{red}\ttfamily,
  morestring=[b]',
  morestring=[b]",
}

% Definition von INI Syntax
\lstdefinelanguage{ini}{
  basicstyle=\ttfamily\small,
  columns=fullflexible,
  morecomment=[s][\color{gray}]{[}{]},
  morecomment=[l]{\#},
  morecomment=[l]{;},
  commentstyle=\color{gray}\ttfamily,
  stringstyle=\color{red}\ttfamily,
  keywordstyle=\color{blue}\bfseries,
}

\lstalias{bash}{sh}

% Kopf- und Fu\ sszeilen
\usepackage{scrlayer-scrpage}
\clearpairofpagestyles
\pagestyle{scrheadings}
\automark{chapter}
\ihead{\small\leftmark}
\ohead{\small\thepage}
\KOMAoptions{headsepline=0.4pt}

% Abstände und Formatierung
\usepackage{setspace}
\onehalfspacing

\setlength{\parskip}{6pt plus2pt minus1pt}
\setlength{\parindent}{0pt}

% Listen
\usepackage[shortlabels]{enumitem}
\setlist{
  before=\vspace{6pt},
  after=\vspace{6pt},
  nosep,
  leftmargin=25pt,
  labelsep=8pt
}
\setlist[itemize,1]{itemsep=4pt}
\setlist[enumerate,1]{itemsep=4pt}

% Hyperref
\usepackage{hyperref}
\hypersetup{
  colorlinks=true,
  linkcolor=black,
  citecolor=teal,
  urlcolor=teal,
  bookmarksnumbered=true
}

% Glossar und Abkürzungen
\usepackage[toc]{glossaries}

% Definiere Glossary-Typ für alle Einträge (Akronyme + Fachwörter zusammen)
\newglossary[glg]{fachwort}{gls}{glo}{Glossar}

% Verwende list-Stil für bessere Trennung der Einträge mit Zeilenumbrüchen
\setglossarystyle{altlist}

\makeglossaries

% Glossareinträge
\newglossaryentry{gls-AI}{
  type=fachwort,
  name=AI,
  first={Artificial Intelligence (AI)},
  text={AI},
  description={Artificial Intelligence - Englische Bezeichnung für KI}
}

\newglossaryentry{gls-ALSA}{
  type=fachwort,
  name=ALSA,
  first={Advanced Linux Sound Architecture (ALSA)},
  text={ALSA},
  description={Advanced Linux Sound Architecture, ein Framework im Linux-Kernel zur Verwaltung und Steuerung von Audio-Hardware}
}

\newglossaryentry{gls-API}{
  type=fachwort,
  name=API,
  first={Application Programming Interface (API)},
  text={API},
  description={Application Programming Interface, eine Schnittstelle für die Kommunikation zwischen Softwarekomponenten}
}

\newglossaryentry{gls-ARI}{
  type=fachwort,
  name=ARI,
  first={Asterisk REST Interface (ARI)},
  text={ARI},
  description={Asterisk REST Interface, eine REST-basierte Programmierschnittstelle zur Steuerung von Asterisk-Telefonsystemen und zur Implementierung eigener Kommunikationslogik}
}

\newglossaryentry{gls-ASR}{
  type=fachwort,
  name=ASR,
  first={Automatic Speech Recognition (ASR)},
  text={ASR},
  description={Automatic Speech Recognition, die automatische Erkennung gesprochener Sprache durch Maschinen}
}

\newglossaryentry{gls-AVV}{
  type=fachwort,
  name=AVV,
  first={Auftragsverarbeitungsvertrag (AVV)},
  text={AVV},
  description={Auftragsverarbeitungsvertrag, ein nach Art.28 DSGVO vorgeschriebener Vertrag zwischen Verantwortlichem und Auftragsverarbeiter, der Zweck, Art und Umfang der Datenverarbeitung sowie technische und organisatorische Ma\ss nahmen zum Schutz personenbezogener Daten regelt}
}

\newglossaryentry{gls-BPO}{
  type=fachwort,
  name=BPO,
  first={Business Process Outsourcing (BPO)},
  text={BPO},
  description={Auslagerung von Geschäftsprozessen an externe Dienstleister, um Kosten zu senken, Effizienz zu steigern oder sich auf Kernkompetenzen zu konzentrieren}
}

\newglossaryentry{gls-CDR}{
  type=fachwort,
  name=CDR,
  first={Call Detail Record (CDR)},
  text={CDR},
  description={Call Detail Record, ein Datensatz mit Informationen über einen durchgeführten Telefonanruf (Dauer, Teilnehmer, Uhrzeit, etc.)}
}

\newglossaryentry{gls-CLI}{
  type=fachwort,
  name=CLI,
  first={Command-Line Interface (CLI)},
  text={CLI},
  description={Command-Line Interface, eine textbasierte Benutzeroberfläche zur Bedienung von Programmen}
}

\newglossaryentry{gls-CPU}{
  type=fachwort,
  name=CPU,
  first={Central Processing Unit (CPU)},
  text={CPU},
  description={Central Processing Unit, der Hauptprozessor eines Computers, der Befehle ausführt}
}

\newglossaryentry{gls-DACH}{
  type=fachwort,
  name=DACH,
  first={Deutschland, Österreich und Schweiz (DACH)},
  text={DACH},
  description={Abkürzung für Deutschland, Österreich und die Schweiz im wirtschaftlichen und geografischen Kontext}
}

\newglossaryentry{gls-DID}{
  type=fachwort,
  name=DID,
  first={Direct Inward Dialing (DID)},
  text={DID},
  description={Direct Inward Dialing, ein Telekommunikationsverfahren, bei dem externe Anrufe direkt an eine interne Nebenstelle einer Telefonanlage weitergeleitet werden, ohne dass eine Vermittlung über eine zentrale Rufannahme erforderlich ist}
}

\newglossaryentry{gls-DPA}{
  type=fachwort,
  name=DPA,
  first={Data Processing Agreement (DPA)},
  text={DPA},
  description={Data Processing Agreement, ein vertragliches Dokument, das die Bedingungen der Datenverarbeitung zwischen Auftraggeber und Auftragsverarbeiter regelt}
}

\newglossaryentry{gls-DPIA}{
  type=fachwort,
  name=DPIA,
  first={Data Protection Impact Assessment (DPIA)},
  text={DPIA},
  description={Datenschutz-Folgenabschätzung anhand Art.~35 DSGVO zur systematischen Bewertung der Risiken einer Datenverarbeitung für die Rechte und Freiheiten betroffener Personen}
}

\newglossaryentry{gls-DSGVO}{
  type=fachwort,
  name=DSGVO,
  first={Datenschutz-Grundverordnung (DSGVO)},
  text={DSGVO},
  description={Datenschutz-Grundverordnung, eine EU-Verordnung zur Regelung des Schutzes personenbezogener Daten und zur Gewährleistung der informationellen Selbstbestimmung}
}

\newglossaryentry{gls-DTLS}{
  type=fachwort,
  name=DTLS,
  first={Datagram Transport Layer Security (DTLS)},
  text={DTLS},
  description={Datagram Transport Layer Security, ein Sicherheitsprotokoll für UDP-basierte Kommunikation}
}

\newglossaryentry{gls-DTMF}{
  type=fachwort,
  name=DTMF,
  first={Dual-Tone Multi-Frequency (DTMF)},
  text={DTMF},
  description={Dual-Tone Multi-Frequency, ein Verfahren zur Signalübertragung von Tastendrücken bei Telefonen, bei dem jede Taste durch eine Kombination aus zwei Frequenzen codiert wird}
}

\newglossaryentry{gls-E2E}{
  type=fachwort,
  name=E2E,
  first={End-to-End (E2E)},
  text={E2E},
  description={End-to-End, eine Bezeichnung für die Gesamtlatenz von Anfang bis Ende eines Prozesses oder einer Kommunikation}
}

\newglossaryentry{gls-EIA}{
  type=fachwort,
  name=EIA,
  first={Ethical Impact Assessment (EIA)},
  text={EIA},
  description={Systematische Bewertung der ethischen Auswirkungen von Technologien oder Systemen, insbesondere im Hinblick auf gesellschaftliche, rechtliche und moralische Konsequenzen}
}

\newglossaryentry{gls-ETSI}{
  type=fachwort,
  name=ETSI,
  first={European Telecommunications Standards Institute (ETSI)},
  text={ETSI},
  description={European Telecommunications Standards Institute, ein Standardisierungsgremium für Telekommunikation und Netzwerke}
}

\newglossaryentry{gls-GDPR}{
  type=fachwort,
  name=GDPR,
  first={General Data Protection Regulation (GDPR)},
  text={GDPR},
  description={Englische Bezeichnung der Datenschutz-Grundverordnung (DSGVO), die den rechtlichen Rahmen für die Verarbeitung personenbezogener Daten innerhalb der Europäischen Union festlegt}
}

\newglossaryentry{gls-HTTP}{
  type=fachwort,
  name=HTTP,
  first={Hypertext Transfer Protocol (HTTP)},
  text={HTTP},
  description={Hypertext Transfer Protocol, ein zustandsloses Anwendungsprotokoll zur Übertragung von Daten im World Wide Web}
}

\newglossaryentry{gls-HTTPS}{
  type=fachwort,
  name=HTTPS,
  first={Hypertext Transfer Protocol Secure (HTTPS)},
  text={HTTPS},
  description={Hypertext Transfer Protocol Secure, ein Protokoll zur sicheren Übertragung von Daten im Internet, das Verschlüsselung (SSL/TLS) und Authentifizierung von Webseiten bietet}
}

\newglossaryentry{gls-ICE}{
  type=fachwort,
  name=ICE,
  first={Interactive Connectivity Establishment (ICE)},
  text={ICE},
  description={Interactive Connectivity Establishment, ein Protokoll zur Überwindung von Firewall- und NAT-Hindernissen}
}

\newglossaryentry{gls-IETF}{
  type=fachwort,
  name=IETF,
  first={Internet Engineering Task Force (IETF)},
  text={IETF},
  description={Internet Engineering Task Force, eine internationale Organisation, die offene Standards für das Internet entwickelt, insbesondere Protokolle und Architekturen}
}

\newglossaryentry{gls-ISO}{
  type=fachwort,
  name=ISO,
  first={International Organization for Standardization (ISO)},
  text={ISO},
  description={International Organization for Standardization, eine internationale Organisation zur Entwicklung und Veröffentlichung von Normen für verschiedene Technologien, Prozesse und Qualitätsstandards}
}

\newglossaryentry{gls-ITU-T}{
  type=fachwort,
  name=ITU-T,
  first={International Telecommunication Union (ITU-T)},
  text={ITU-T},
  description={International Telecommunication Union - Standardisierungsgremium für Telekommunikation}
}

\newglossaryentry{gls-IVR}{
  type=fachwort,
  name=IVR,
  first={Interactive Voice Response (IVR)},
  text={IVR},
  description={Interactive Voice Response, ein System für automatisierte Sprachmenüs und Sprachverarbeitung}
}

\newglossaryentry{gls-JSON}{
  type=fachwort,
  name=JSON,
  first={JavaScript Object Notation (JSON)},
  text={JSON},
  description={JavaScript Object Notation, ein text-basiertes Datenformat für den Austausch strukturierter Daten}
}

\newglossaryentry{gls-JWT}{
  type=fachwort,
  name=JWT,
  first={JSON Web Token (JWT)},
  text={JWT},
  description={JSON Web Token, ein offener Standard (RFC 7519) zur sicheren Übertragung von Claims zwischen Parteien, häufig verwendet für Authentifizierung und Autorisierung in Webanwendungen}
}

\newglossaryentry{gls-KI}{
  type=fachwort,
  name=KI,
  first={Künstliche Intelligenz (KI)},
  text={KI},
  description={Künstliche Intelligenz bezeichnet die Fähigkeit von Maschinen, menschenähnliche kognitive Leistungen wie Lernen, Problemlösen und Sprachverstehen auszuführen}
}

\newglossaryentry{gls-KMS}{
  type=fachwort,
  name=KMS,
  first={Key Management Service (KMS)},
  text={KMS},
  description={Key Management Service, ein Dienst zur sicheren Verwaltung, Speicherung und Rotation kryptografischer Schlüssel, der häufig in Cloud- und Container-Umgebungen eingesetzt wird}
}

\newglossaryentry{gls-KMU}{
  type=fachwort,
  name=KMU,
  first={kleine und mittlere Unternehmen (KMU)},
  text={KMU},
  description={kleine und mittelständische Unternehmen}
}

\newglossaryentry{gls-LAN}{
  type=fachwort,
  name=LAN,
  first={Local Area Network (LAN)},
  text={LAN},
  description={Local Area Network, ein räumlich begrenztes Rechnernetzwerk, das Geräte innerhalb eines lokalen Bereichs wie eines Gebäudes oder Campus verbindet}
}

\newglossaryentry{gls-LLM}{
  type=fachwort,
  name=LLM,
  first={Large Language Model (LLM)},
  text={LLM},
  description={Large Language Model, ein auf gro\ss en Textkorpora trainiertes KI-Sprachmodell, das natürliche Sprache generieren oder verstehen kann}
}

\newglossaryentry{gls-MANO}{
  type=fachwort,
  name=MANO,
  first={Management and Orchestration (MANO)},
  text={MANO},
  description={Eine Architektur- und Funktionsschicht im NFV-Umfeld zur Verwaltung, Orchestrierung und Überwachung von virtualisierten Netzwerkfunktionen (VNFs) sowie der zugrunde liegenden Infrastruktur}
}

\newglossaryentry{gls-MEC}{
  type=fachwort,
  name=MEC,
  first={Multi-access Edge Computing (MEC)},
  text={MEC},
  description={Multi-access Edge Computing, ein Architekturkonzept, bei dem Rechen- und Speicherkapazitäten nahe am Netzwerkrand bereitgestellt werden, um Latenzen zu reduzieren und zeitkritische Anwendungen zu unterstützen}
}

\newglossaryentry{gls-MOS}{
  type=fachwort,
  name=MOS,
  first={Mean Opinion Score (MOS)},
  text={MOS},
  description={Mean Opinion Score, eine Metrik zur Bewertung der Audioqualität auf einer Skala}
}

\newglossaryentry{gls-NAT}{
  type=fachwort,
  name=NAT,
  first={Network Address Translation (NAT)},
  text={NAT},
  description={Network Address Translation, ein Verfahren zur Umwandlung von IP-Adressen in Netzwerkpaketen}
}

\newglossaryentry{gls-NDJSON}{
  type=fachwort,
  name=NDJSON,
  first={Newline Delimited JSON (NDJSON)},
  text={NDJSON},
  description={Newline Delimited JSON, ein Datenformat zur sequentiellen Darstellung von JSON-Objekten, bei dem jedes Objekt in einer eigenen Zeile steht, wodurch sich das Format besonders für Streaming, Log-Dateien und die Verarbeitung gro\ss er Datenmengen eignet}
}

\newglossaryentry{gls-NFV}{
  type=fachwort,
  name=NFV,
  first={Network Function Virtualization (NFV)},
  text={NFV},
  description={Network Function Virtualization, ein Paradigma zur Ablösung dedizierter Hardware durch softwarebasierte, containerisierte Netzwerkdienste}
}

\newglossaryentry{gls-NFVI}{
  type=fachwort,
  name=NFVI,
  first={Network Function Virtualization Infrastructure (NFVI)},
  text={NFVI},
  description={Network Function Virtualization Infrastructure, die physische Hardware- und Software-Ressourcenschicht, auf der VNFs laufen}
}

\newglossaryentry{gls-NFVO}{
  type=fachwort,
  name=NFVO,
  first={Network Function Virtualization Orchestration (NFVO)},
  text={NFVO},
  description={Network Function Virtualization Orchestration, die Verwaltungs- und Orchestrierungsschicht für VNFs und deren Lifecycle-Management}
}

\newglossaryentry{gls-NLP}{
  type=fachwort,
  name=NLP,
  first={Natural Language Processing (NLP)},
  text={NLP},
  description={Natural Language Processing bezeichnet die automatische Verarbeitung, Analyse und Generierung menschlicher Sprache durch Computer, z.B. für Textverständnis, Sprachmodellierung oder Dialogsysteme}
}

\newglossaryentry{gls-OSS}{
  type=fachwort,
  name=OSS,
  first={Open Source Software (OSS)},
  text={OSS},
  description={Open Source Software, Software mit öffentlich verfügbarem Quellcode, die frei verändert und verteilt werden darf}
}

\newglossaryentry{gls-PABX}{
  type=fachwort,
  name=PABX,
  first={Private Automatic Branch Exchange (PABX)},
  text={PABX},
  description={Private Automatic Branch Exchange, eine private, automatische Nebenstellenanlage zur Vermittlung interner und externer Sprachverbindungen in Unternehmensnetzwerken. Auch als klassische Telefonanlage oder Private Branch Exchange (PBX) bekannt}
}

\newglossaryentry{gls-PBX}{
  type=fachwort,
  name=PBX,
  first={Private Branch Exchange (PBX)},
  text={PBX},
  description={Private Branch Exchange, eine Telefonanlage für interne und externe Sprachkommunikation}
}

\newglossaryentry{gls-PCM}{
  type=fachwort,
  name=PCM,
  first={Pulse Code Modulation (PCM)},
  text={PCM},
  description={Pulse Code Modulation, ein Verfahren zur digitalen Darstellung analoger Signale durch Abtastung, Quantisierung und Codierung, häufig verwendet in der digitalen Sprach- und Audioübertragung}
}

\newglossaryentry{gls-PCM16}{
  type=fachwort,
  name=PCM16,
  first={Pulse Code Modulation 16 Bit (PCM16)},
  text={PCM16},
  description={Pulse Code Modulation mit 16 Bit Auflösung, ein unkomprimiertes Audioformat zur digitalen Darstellung von Audiosignalen, das häufig in der Sprachverarbeitung und Echtzeitkommunikation verwendet wird}
}

\newglossaryentry{gls-PESQ}{
  type=fachwort,
  name=PESQ,
  first={Perceptual Evaluation of Speech Quality (PESQ)},
  text={PESQ},
  description={Perceptual Evaluation of Speech Quality, ein objektives Verfahren zur Bewertung der Sprachqualität durch Vergleich eines Referenzsignals mit einem gestörten Signal, standardisiert in ITU-T P.862 und weit verbreitet zur Bewertung von Sprachübertragungssystemen}
}

\newglossaryentry{gls-PII}{
  type=fachwort,
  name=PII,
  first={Personally Identifiable Information (PII)},
  text={PII},
  description={Personally Identifiable Information, persönliche Daten, die zur Identifikation einer Person verwendet werden können}
}

\newglossaryentry{gls-PJSIP}{
  type=fachwort,
  name=PJSIP,
  first={PJSIP (PJSIP)},
  text={PJSIP},
  description={Open-Source-Bibliothek für SIP-basierte Sprach- und Videoübertragung, die Funktionen für VoIP, Instant Messaging und NAT-Traversierung bereitstellt}
}

\newglossaryentry{gls-POLQA}{
  type=fachwort,
  name=POLQA,
  first={Perceptual Objective Listening Quality Assessment (POLQA)},
  text={POLQA},
  description={Perceptual Objective Listening Quality Assessment, ein objektives Messverfahren zur Bewertung der wahrgenommenen Sprachqualität in Telekommunikationssystemen, standardisiert in ITU-T P.863 und ausgelegt für moderne Sprachdienste wie VoIP, HD-Voice und Mobilfunknetze}
}

\newglossaryentry{gls-PSTN}{
  type=fachwort,
  name=PSTN,
  first={Public Switched Telephone Network (PSTN)},
  text={PSTN},
  description={Public Switched Telephone Network, das klassische öffentliche Telefonnetz mit Sprachübertragung}
}

\newglossaryentry{gls-RAG}{
  type=fachwort,
  name=RAG,
  first={Retrieval-Augmented Generation (RAG)},
  text={RAG},
  description={Retrieval-Augmented Generation, ein Ansatz in der natürlichen Sprachverarbeitung, bei dem ein generatives Sprachmodell mit einer externen Wissensquelle kombiniert wird, um Antworten auf Basis abgerufener, kontextrelevanter Informationen zu erzeugen}
}

\newglossaryentry{gls-RAM}{
  type=fachwort,
  name=RAM,
  first={Random Access Memory (RAM)},
  text={RAM},
  description={Random Access Memory, der Arbeitsspeicher eines Computers für temporäre Datenablage}
}

\newglossaryentry{gls-RBAC}{
  type=fachwort,
  name=RBAC,
  first={Role-Based Access Control (RBAC)},
  text={RBAC},
  description={Role-Based Access Control, ein Zugriffskontrollmodell, das Berechtigungen auf Basis von Benutzerrollen vergibt}
}

\newglossaryentry{gls-REST}{
  type=fachwort,
  name=REST,
  first={Representational State Transfer (REST)},
  text={REST},
  description={Representational State Transfer, ein Architekturstil für die Gestaltung von Web-APIs}
}

\newglossaryentry{gls-RFC}{
  type=fachwort,
  name=RFC,
  first={Request for Comments (RFC)},
  text={RFC},
  description={Request for Comments, eine von der Internet Engineering Task Force veröffentlichte Dokumentenreihe zur Spezifikation, Beschreibung und Diskussion von Internetstandards und -protokollen}
}

\newglossaryentry{gls-RMS}{
  type=fachwort,
  name=RMS,
  first={Root Mean Square (RMS)},
  text={RMS},
  description={Root Mean Square, ein mathematisches Ma\ss zur Bestimmung des quadratischen Mittelwerts eines Signals, häufig verwendet zur Beschreibung der effektiven Stärke von Wechselspannungen oder Strömen}
}

\newglossaryentry{gls-RTC}{
  type=fachwort,
  name=RTC,
  first={Real-Time Communication (RTC)},
  text={RTC},
  description={Real-Time Communication, Sammelbegriff für Technologien zur zeitkritischen Übertragung von Audio-, Video- und Datenströmen in Echtzeit über Netzwerke}
}

\newglossaryentry{gls-RTCP}{
  type=fachwort,
  name=RTCP,
  first={Real-Time Control Protocol (RTCP)},
  text={RTCP},
  description={Ein Begleitprotokoll zu RTP, das für die Überwachung und Kontrolle von Echtzeit-Medienströmen zuständig ist}
}

\newglossaryentry{gls-RTP}{
  type=fachwort,
  name=RTP,
  first={Real-time Transport Protocol (RTP)},
  text={RTP},
  description={Real-time Transport Protocol, das Protokoll zur Übertragung von Audio und Videodaten in Echtzeit}
}

\newglossaryentry{gls-SCC}{
  type=fachwort,
  name=SCC,
  first={Standard Contractual Clauses (SCC)},
  text={SCC},
  description={Standard Contractual Clauses, von der EU bereitgestellte Standardvertragsmuster für die sichere Datenübermittlung in Drittstaaten}
}

\newglossaryentry{gls-SDK}{
  type=fachwort,
  name=SDK,
  first={Software Development Kit (SDK)},
  text=SDK,
  description={Software Development Kit, eine Sammlung von Werkzeugen und Bibliotheken zur Entwicklung von Anwendungen}
}

\newglossaryentry{gls-SFU}{
  type=fachwort,
  name=SFU,
  first={Selective Forwarding Unit (SFU)},
  text=SFU,
  description={Selective Forwarding Unit, ein Media-Server für WebRTC, der Audio- und Videodatenströme zwischen Teilnehmern weiterleitet}
}

\newglossaryentry{gls-SIP}{
  type=fachwort,
  name=SIP,
  first={Session Initiation Protocol (SIP)},
  text={SIP},
  description={Session Initiation Protocol, ein Standardprotokoll zur Signalisierung und Verwaltung von VoIP-Anrufen}
}

\newglossaryentry{gls-SLA}{
  type=fachwort,
  name=SLA,
  first={Service Level Agreement (SLA)},
  text=SLA,
  description={Service Level Agreement, eine vertragliche Vereinbarung, die die erwartete Servicequalität (z.B. Verfügbarkeit, Latenz) definiert}
}

\newglossaryentry{gls-SPA}{
  type=fachwort,
  name=SPA,
  first={Single-Page Application (SPA)},
  text=SPA,
  description={Single-Page Application, eine Webanwendung, die in einem Single Document geladen wird und dynamisch Inhalte aktualisiert}
}

\newglossaryentry{gls-SRTP}{
  type=fachwort,
  name=SRTP,
  first={Secure Real-time Transport Protocol (SRTP)},
  text=SRTP,
  description={Secure Real-time Transport Protocol, eine sichere Variante von RTP mit Verschlüsselung}
}

\newglossaryentry{gls-STT}{
  type=fachwort,
  name=STT,
  first={Speech-to-Text (STT)},
  text=STT,
  description={Speech-to-Text, die Umwandlung von Sprache in geschriebenen Text}
}

\newglossaryentry{gls-STUN}{
  type=fachwort,
  name=STUN,
  first={Session Traversal Utilities for NAT (STUN)},
  text=STUN,
  description={Session Traversal Utilities for NAT, ein Protokoll zur Ermittlung öffentlicher IP-Adressen und Ports zur Unterstützung der Kommunikation durch NAT-Gateways}
}

\newglossaryentry{gls-TCP}{
  type=fachwort,
  name=TCP,
  first={Transmission Control Protocol (TCP)},
  text=TCP,
  description={Transmission Control Protocol, ein verbindungsorientiertes Transportprotokoll der Internetprotokollfamilie, das eine zuverlässige, geordnete und fehlerfreie Datenübertragung gewährleistet}
}

\newglossaryentry{gls-TLS}{
  type=fachwort,
  name=TLS,
  first={Transport Layer Security (TLS)},
  text=TLS,
  description={Transport Layer Security, ein Kryptografieprotokoll für sichere Kommunikation über Netzwerke}
}

\newglossaryentry{gls-TTL}{
  type=fachwort,
  name=TTL,
  first={Time-To-Live (TTL)},
  text=TTL,
  description={Time-To-Live, ein Parameter in Netzwerkprotokollen (insbesondere IP), der die maximale Anzahl von Hops (Sprüngen) begrenzt, die ein Paket durchlaufen kann, bevor es verworfen wird}
}

\newglossaryentry{gls-TTS}{
  type=fachwort,
  name=TTS,
  first={Text-to-Speech (TTS)},
  text=TTS,
  description={Text-to-Speech, die Umwandlung von geschriebenem Text in gesprochene Sprache}
}

\newglossaryentry{gls-TURN}{
  type=fachwort,
  name=TURN,
  first={Traversal Using Relays around NAT (TURN)},
  text=TURN,
  description={Traversal Using Relays around NAT, ein Protokoll zur relay-basierten Weiterleitung von Medienströmen, wenn direkte Verbindungen durch NAT oder Firewalls nicht möglich sind}
}

\newglossaryentry{gls-UDP}{
  type=fachwort,
  name=UDP,
  first={User Datagram Protocol (UDP)},
  text=UDP,
  description={User Datagram Protocol, ein verbindungsloses Transportprotokoll der Internetprotokollfamilie, das eine schnelle, aber nicht garantierte Übertragung von Daten ohne Reihenfolge oder Zustellsicherung ermöglicht}
}

\newglossaryentry{gls-UI}{
  type=fachwort,
  name=UI,
  first={User Interface (UI)},
  text=UI,
  description={User Interface, die Benutzeroberfläche einer Anwendung, über die Benutzer mit dem System interagieren}
}

\newglossaryentry{gls-VAD}{
  type=fachwort,
  name=VAD,
  first={Voice Activity Detection (VAD)},
  text=VAD,
  description={Voice Activity Detection, die Erkennung von Sprachaktivität in Audiodaten}
}

\newglossaryentry{gls-VLAN}{
  type=fachwort,
  name=VLAN,
  first={Virtual Local Area Network (VLAN)},
  text=VLAN,
  description={Virtual Local Area Network, ein logisches Teilnetz zur Segmentierung von Netzwerken auf Layer 2, das zur Trennung von Datenverkehr und zur Erhöhung von Sicherheit und Übersichtlichkeit eingesetzt wird}
}

\newglossaryentry{gls-VM}{
  type=fachwort,
  name=VM,
  first={Virtuelle Maschine (VM)},
  text=VM,
  description={Virtuelle Maschine, eine softwarebasierte Emulation eines Computersystems, die es ermöglicht, Betriebssysteme und Anwendungen isoliert auf derselben Hardware auszuführen}
}

\newglossaryentry{gls-VNF}{
  type=fachwort,
  name=VNF,
  first={Virtualized Network Function (VNF)},
  text=VNF,
  description={Virtualized Network Function, eine softwarebasierte Implementierung einer Netzwerkfunktion, die auf Standard-Hardware läuft}
}

\newglossaryentry{gls-YAML}{
  type=fachwort,
  name=YAML,
  first={YAML Ain't Markup Language (YAML)},
  text=YAML,
  description={YAML Ain’t Markup Language, ein menschenlesbares Datenserialisierungsformat zur strukturierten Darstellung von Konfigurationsdaten, das durch Einrückungen anstelle von Klammern oder Tags arbeitet}
}

\newglossaryentry{gls-Asterisk}{
  type=fachwort,
  name=Asterisk,
  description={Open-Source-Software-PBX, die Signalisierung, Rufaufbau und Medienverarbeitung für VoIP- und klassische Telefonsysteme bereitstellt}
}
\glsadd{gls-Asterisk}

\newglossaryentry{gls-AudioDec}{
  type=fachwort,
  name=AudioDec,
  description={Neuraler Audio-Codec-Ansatz mit sehr niedriger Latenz, der für Echtzeit-Anwendungen optimiert ist}
}
\glsadd{gls-AudioDec}

\newglossaryentry{gls-Aufzeichnungs-Opt-in}{
  type=fachwort,
  name=Aufzeichnungs-Opt-in,
  description={Die ausdrückliche Zustimmung eines Benutzers oder Gesprächspartners zur Aufzeichnung von Audio-, Video- oder Transaktionsdaten, die für Compliance-, Qualitäts- oder Analysezwecke erhoben werden}
}
\glsadd{gls-Aufzeichnungs-Opt-in}

\newglossaryentry{gls-Bridge}{
  type=fachwort,
  name=Bridge,
  description={Netzwerk- oder Softwarekomponente, die zwei unterschiedliche Netzwerke, Protokolle oder Systeme verbindet und Datenaustausch ermöglicht, z.B. SIP-WebRTC-Bridge}
}
\glsadd{gls-Bridge}

\newglossaryentry{gls-Codec}{
  type=fachwort,
  name=Codec,
  description={Encoder-Decoder-Verfahren zur Digitalisierung, Kompression und Dekomprimierung von Audio- oder Videodaten. Beispiele: G.711 (klassische Telefonie), Opus (WebRTC), PCM16 (Spracherkennung)}
}
\glsadd{gls-Codec}

\newglossaryentry{gls-Container}{
  type=fachwort,
  name=Container,
  description={Eine leichtgewichtige virtualisierte Laufzeitumgebung, in der Software und ihre Abhängigkeiten isoliert ausgeführt werden}
}
\glsadd{gls-Container}

\newglossaryentry{gls-Whisper}{
  type=fachwort,
  name=Whisper,
  description={OpenAI-Spracherkennungsmodell für mehrsprachige Speech-to-Text-Verarbeitung, verfügbar als Cloud-API und als lokal ausführbares Modell}
}
\glsadd{gls-Whisper}

\newglossaryentry{gls-Whitelist}{
  type=fachwort,
  name=Whitelist,
  description={Liste von erlaubten IP-Adressen, Benutzern oder anderen Entities. Nur Einträge auf der Whitelist werden akzeptiert, alle anderen werden blockiert}
}
\glsadd{gls-Whitelist}

\newglossaryentry{gls-Baresip}{
  type=fachwort,
  name=Baresip,
  description={Open-Source SIP User Agent für experimentelle Tests und grundlegende SIP-Registrierung, der für Protokollkompatibilitätsprüfungen eingesetzt wird}
}
\glsadd{gls-Baresip}

\newglossaryentry{gls-LiveKit}{
  type=fachwort,
  name=LiveKit,
  description={Open-Source WebRTC-Framework und Media-Server für Echtzeitkommunikation, das Raumerstellung, Audiostream-Verarbeitung und KI-Integration unterstützt}
}
\glsadd{gls-LiveKit}

\newglossaryentry{gls-Media-Proxy}{
  type=fachwort,
  name=Media-Proxy,
  description={Netzwerkkomponente zur Weiterleitung, Anpassung und Transkodierung von Audio- und Videoströmen zwischen unterschiedlichen Protokollen und Codecs}
}
\glsadd{gls-Media-Proxy}

\newglossaryentry{gls-RTP-Bridge}{
  type=fachwort,
  name=RTP-Bridge,
  description={Komponente zur Weiterleitung, Anpassung und Transkodierung von RTP-Medienströmen zwischen verschiedenen Diensten und Protokollen}
}
\glsadd{gls-RTP-Bridge}

\newglossaryentry{gls-Hypervisor}{
  type=fachwort,
  name=Hypervisor,
  description={Virtualisierungssoftware, die es ermöglicht, mehrere virtuelle Maschinen auf einem physischen Host-System auszuführen und deren Ressourcenzugriff zu verwalten}
}
\glsadd{gls-Hypervisor}

\newglossaryentry{gls-Docker}{
  type=fachwort,
  name=Docker,
  description={Containerisierungsplattform zur Erstellung, Verwaltung und Ausführung von Anwendungen in isolierten Container-Umgebungen}
}
\glsadd{gls-Docker}

\newglossaryentry{gls-Docker-Image}{
  type=fachwort,
  name=Docker-Image,
  description={Schichtenbasierte, unveränderliche Vorlage zur Erstellung von Docker-Containern, die alle benötigten Abhängigkeiten, Bibliotheken und Konfigurationen enthält}
}
\glsadd{gls-Docker-Image}

\newglossaryentry{gls-Kubernetes}{
  type=fachwort,
  name=Kubernetes,
  description={Open-Source-Orchestrierungsplattform für die Automatisierung von Bereitstellung, Skalierung und Verwaltung containerisierter Anwendungen}
}
\glsadd{gls-Kubernetes}

\newglossaryentry{gls-Kamailio}{
  type=fachwort,
  name=Kamailio,
  description={Open-Source SIP-Server und Proxy für hochperformante VoIP-Telefonie, der für Lastverteilung und SIP-Routing eingesetzt wird}
}
\glsadd{gls-Kamailio}

\newglossaryentry{gls-Dialplan}{
  type=fachwort,
  name=Dialplan,
  description={Regelwerk in einer PBX-Software zur Definition von Anrufverarbeitungslogik, Routing-Entscheidungen und Aktionsabläufen}
}
\glsadd{gls-Dialplan}

\newglossaryentry{gls-Jitter-Puffer}{
  type=fachwort,
  name=Jitter-Puffer,
  description={Zwischenspeicher zur Kompensation von Verzögerungsschwankungen (Jitter) bei der Übertragung von Medienpaketen in IP-Netzwerken, um eine gleichmäßige Wiedergabe zu gewährleisten}
}
\glsadd{gls-Jitter-Puffer}

\newglossaryentry{gls-Streaming}{
  type=fachwort,
  name=Streaming,
  description={Übertragungstechnik, bei der Daten kontinuierlich in Echtzeit gesendet und verarbeitet werden, ohne das vollständige Ergebnis abzuwarten}
}
\glsadd{gls-Streaming}

\newglossaryentry{gls-Token-LLM}{
  type=fachwort,
  name=Token,
  description={Kleinste Verarbeitungseinheit in Sprachmodellen, die Wörter, Wortteile oder Zeichen repräsentiert und bei der Textgenerierung verwendet wird}
}
\glsadd{gls-Token-LLM}

\newglossaryentry{gls-ConfigMaps}{
  type=fachwort,
  name=ConfigMaps,
  description={Kubernetes-Ressource zur Verwaltung von Konfigurationsdaten als Key-Value-Paare, die von Containern geladen werden können, ohne dass diese in Container-Images eingebettet werden müssen}
}
\glsadd{gls-ConfigMaps}

\newglossaryentry{gls-Secrets}{
  type=fachwort,
  name=Secrets,
  description={Kubernetes-Objekt zur sicheren Speicherung vertraulicher Daten wie Passwörter, API-Schlüssel und Zertifikate}
}
\glsadd{gls-Secrets}

\newglossaryentry{gls-Persistent-Volume}{
  type=fachwort,
  name=Persistent Volume,
  description={Kubernetes-Ressource zur Bereitstellung dauerhaften Speicherplatzes für Container, der über den Lebenszyklus einzelner Pods hinaus erhalten bleibt}
}
\glsadd{gls-Persistent-Volume}

\newglossaryentry{gls-ReplicaSet}{
  type=fachwort,
  name=ReplicaSet,
  description={Kubernetes-Controller zur Sicherstellung, dass eine definierte Anzahl identischer Pod-Replikas jederzeit ausgeführt wird}
}
\glsadd{gls-ReplicaSet}

\newglossaryentry{gls-Pod}{
  type=fachwort,
  name=Pod,
  description={Kleinste ausführbare Einheit in Kubernetes, die einen oder mehrere Container enthält und gemeinsam bereitgestellt wird}
}
\glsadd{gls-Pod}

\newglossaryentry{gls-Namespace}{
  type=fachwort,
  name=Namespace,
  description={Logische Trennung von Ressourcen innerhalb eines Kubernetes-Clusters zur Isolation verschiedener Anwendungen oder Teams}
}
\glsadd{gls-Namespace}

\newglossaryentry{gls-NetworkPolicy}{
  type=fachwort,
  name=NetworkPolicy,
  description={Kubernetes-Ressource zur Definition netzwerkbasierter Zugriffskontrollregeln zwischen Pods}
}
\glsadd{gls-NetworkPolicy}

\newglossaryentry{gls-Ingress-Controller}{
  type=fachwort,
  name=Ingress-Controller,
  description={Kubernetes-Komponente zur Verwaltung des externen HTTP/HTTPS-Zugriffs auf Services innerhalb eines Clusters}
}
\glsadd{gls-Ingress-Controller}

\newglossaryentry{gls-Least-Privilege-Prinzip}{
  type=fachwort,
  name=Least-Privilege-Prinzip,
  description={Sicherheitsprinzip, bei dem Benutzern und Prozessen nur die minimal notwendigen Berechtigungen erteilt werden, um ihre Aufgaben zu erfüllen}
}
\glsadd{gls-Least-Privilege-Prinzip}

\newglossaryentry{gls-Digest-Authentifizierung}{
  type=fachwort,
  name=Digest-Authentifizierung,
  description={HTTP-Authentifizierungsverfahren, bei dem Passwörter gehasht übertragen werden, um Abhören zu erschweren}
}
\glsadd{gls-Digest-Authentifizierung}

\newglossaryentry{gls-Horizontal-Pod-Autoscaler}{
  type=fachwort,
  name=Horizontal Pod Autoscaler,
  description={Kubernetes-Mechanismus zur automatischen Anpassung der Anzahl von Pod-Replicas basierend auf CPU-, Speicher- oder benutzerdefinierten Metriken}
}
\glsadd{gls-Horizontal-Pod-Autoscaler}

\newglossaryentry{gls-SIP-Trunk}{
  type=fachwort,
  name=SIP-Trunk,
  description={Virtuelle Verbindungsleitung zwischen einer PBX und einem VoIP-Provider zur Übertragung mehrerer gleichzeitiger Sprachkanäle über IP-Netzwerke}
}
\glsadd{gls-SIP-Trunk}

\newglossaryentry{gls-Intent-Erkennung}{
  type=fachwort,
  name=Intent-Erkennung,
  description={Verfahren zur automatischen Identifikation der Absicht oder des Ziels einer Nutzeräußerung in einem Dialog, häufig eingesetzt in Chatbots und Sprachassistenten}
}
\glsadd{gls-Intent-Erkennung}

\newglossaryentry{gls-Dialog-Management}{
  type=fachwort,
  name=Dialog-Management,
  description={Steuerungslogik für die Verwaltung des Gesprächsflusses in einem KI-basierten Dialog, die Kontexte verfolgt und Antworten koordiniert}
}
\glsadd{gls-Dialog-Management}

\newglossaryentry{gls-Task-Orchestrierung}{
  type=fachwort,
  name=Task-Orchestrierung,
  description={Koordination und Steuerung mehrerer zusammenhängender Aufgaben oder Prozesse in einem System}
}
\glsadd{gls-Task-Orchestrierung}

\newglossaryentry{gls-Geocodierung}{
  type=fachwort,
  name=Geocodierung,
  description={Prozess zur Umwandlung von Adressinformationen in geografische Koordinaten (Längen- und Breitengrad)}
}
\glsadd{gls-Geocodierung}

\newglossaryentry{gls-Prometheus}{
  type=fachwort,
  name=Prometheus,
  description={Open-Source-Monitoring-System zur Erfassung und Speicherung von Zeitreihendaten (Metriken) aus IT-Systemen}
}
\glsadd{gls-Prometheus}

\newglossaryentry{gls-Grafana}{
  type=fachwort,
  name=Grafana,
  description={Open-Source-Plattform zur Visualisierung und Analyse von Metriken, Logs und Traces aus verschiedenen Datenquellen}
}
\glsadd{gls-Grafana}

\newglossaryentry{gls-PostgreSQL}{
  type=fachwort,
  name=PostgreSQL,
  description={Open-Source relationales Datenbankmanagementsystem mit umfangreicher SQL-Unterstützung}
}
\glsadd{gls-PostgreSQL}

\newglossaryentry{gls-MySQL}{
  type=fachwort,
  name=MySQL,
  description={Weit verbreitetes Open-Source relationales Datenbankmanagementsystem}
}
\glsadd{gls-MySQL}

\newglossaryentry{gls-Monolithisches-Modell}{
  type=fachwort,
  name=Monolithisches Modell,
  description={Softwarearchitektur, bei der alle Funktionen und Komponenten in einer einzigen, eng gekoppelten Anwendung integriert sind}
}
\glsadd{gls-Monolithisches-Modell}

\newglossaryentry{gls-Mikroservice}{
  type=fachwort,
  name=Mikroservice,
  description={Architekturmuster, bei dem eine Anwendung in kleine, unabhängig deploybare Services aufgeteilt wird, die über definierte Schnittstellen kommunizieren}
}
\glsadd{gls-Mikroservice}

\newglossaryentry{gls-Hybrid-Ansatz}{
  type=fachwort,
  name=Hybrid-Ansatz,
  description={Architekturmodell, das Elemente monolithischer und mikroservice-basierter Ansätze kombiniert}
}
\glsadd{gls-Hybrid-Ansatz}

\newglossaryentry{gls-Barge-in}{
  type=fachwort,
  name=Barge-in,
  description={Funktion in Sprachdialogsystemen, die es Nutzern ermöglicht, die Systemausgabe zu unterbrechen und sofort zu antworten}
}
\glsadd{gls-Barge-in}

\newglossaryentry{gls-Edge-Deployment}{
  type=fachwort,
  name=Edge-Deployment,
  description={Bereitstellungsmodell, bei dem Anwendungen nahe am Netzwerkrand (z.B. auf lokalen Servern, Edge-Knoten oder IoT-Geräten) statt in zentralen Rechenzentren betrieben werden}
}
\glsadd{gls-Edge-Deployment}

\newglossaryentry{gls-Vendor-Lock-in}{
  type=fachwort,
  name=Vendor-Lock-in,
  description={Abhängigkeit von einem spezifischen Anbieter, bei der ein Wechsel zu alternativen Lösungen technisch oder wirtschaftlich erschwert wird}
}
\glsadd{gls-Vendor-Lock-in}

\newglossaryentry{gls-On-Premise}{
  type=fachwort,
  name=On-Premise,
  description={Betriebsmodell, bei dem Software und Infrastruktur auf eigenen Servern im Unternehmen oder lokalen Rechenzentren statt in der Cloud betrieben werden}
}
\glsadd{gls-On-Premise}

\newglossaryentry{gls-GitHub-Copilot}{
  type=fachwort,
  name=GitHub Copilot,
  description={KI-gestütztes Programmierwerkzeug zur automatischen Code-Vervollständigung und Generierung von Code-Vorschlägen}
}
\glsadd{gls-GitHub-Copilot}

\newglossaryentry{gls-Next-js}{
  type=fachwort,
  name=Next.js,
  description={React-Framework für serverseitiges Rendering und statische Website-Generierung}
}
\glsadd{gls-Next-js}

\newglossaryentry{gls-React}{
  type=fachwort,
  name=React,
  description={JavaScript-Bibliothek zur Entwicklung von Benutzeroberflächen mit komponentenbasierter Architektur}
}
\glsadd{gls-React}

\newglossaryentry{gls-Node-js}{
  type=fachwort,
  name=Node.js,
  description={JavaScript-Laufzeitumgebung für serverseitige Anwendungen}
}
\glsadd{gls-Node-js}

\newglossaryentry{gls-npm}{
  type=fachwort,
  name=npm,
  description={Package Manager für JavaScript und Node.js zur Verwaltung von Abhängigkeiten}
}
\glsadd{gls-npm}

\newglossaryentry{gls-pip}{
  type=fachwort,
  name=pip,
  description={Package Installer für Python zur Installation und Verwaltung von Python-Paketen}
}
\glsadd{gls-pip}

\newglossaryentry{gls-Docker-Desktop}{
  type=fachwort,
  name=Docker Desktop,
  description={Desktop-Anwendung zur lokalen Entwicklung und Verwaltung von Docker-Containern unter Windows und macOS}
}
\glsadd{gls-Docker-Desktop}

\newglossaryentry{gls-Docker-Compose}{
  type=fachwort,
  name=Docker-Compose,
  description={Werkzeug zur Definition und Verwaltung von Multi-Container-Docker-Anwendungen mittels YAML-Konfigurationsdateien}
}
\glsadd{gls-Docker-Compose}

\newglossaryentry{gls-Bridge-Netzwerk}{
  type=fachwort,
  name=Bridge-Netzwerk,
  description={Docker-Netzwerkmodus, bei dem Container über ein virtuelles Netzwerk miteinander kommunizieren}
}
\glsadd{gls-Bridge-Netzwerk}

\newglossaryentry{gls-Redis}{
  type=fachwort,
  name=Redis,
  description={Open-Source In-Memory-Datenbank für schnelle Datenzugriffe, häufig für Caching und Session-Verwaltung eingesetzt}
}
\glsadd{gls-Redis}

\newglossaryentry{gls-Observer-Pattern}{
  type=fachwort,
  name=Observer-Pattern,
  description={Entwurfsmuster, bei dem Objekte automatisch über Zustandsänderungen anderer Objekte benachrichtigt werden}
}
\glsadd{gls-Observer-Pattern}

\newglossaryentry{gls-Hook}{
  type=fachwort,
  name=Hook,
  description={React-Funktion zur Nutzung von Zustandsverwaltung und Seiteneffekten in funktionalen Komponenten}
}
\glsadd{gls-Hook}

\newglossaryentry{gls-Server-Side-Token-Pattern}{
  type=fachwort,
  name=Server-Side Token Pattern,
  description={Architekturmuster, bei dem Authentifizierungs-Tokens serverseitig generiert werden, um Geheimnisse vor dem Client zu schützen}
}
\glsadd{gls-Server-Side-Token-Pattern}

\newglossaryentry{gls-Exponentielles-Backoff}{
  type=fachwort,
  name=Exponentielles Backoff,
  description={Strategie zur automatischen Wiederholung fehlgeschlagener Vorgänge mit exponentiell steigenden Wartezeiten}
}
\glsadd{gls-Exponentielles-Backoff}

\newglossaryentry{gls-Python-Package-Index}{
  type=fachwort,
  name=Python Package Index,
  description={Zentrales Repository für Python-Softwarepakete (auch PyPI genannt)}
}
\glsadd{gls-Python-Package-Index}

\newglossaryentry{gls-Zustandsmaschine}{
  type=fachwort,
  name=Zustandsmaschine,
  description={Modell zur Beschreibung des Verhaltens eines Systems durch definierte Zustände und Übergänge}
}
\glsadd{gls-Zustandsmaschine}

\newglossaryentry{gls-WebSocket}{
  type=fachwort,
  name=WebSocket,
  description={Protokoll für bidirektionale, persistente Kommunikation zwischen Client und Server}
}
\glsadd{gls-WebSocket}

\newglossaryentry{gls-Silero-VAD}{
  type=fachwort,
  name=Silero VAD,
  description={Lokal ausführbares Voice Activity Detection Modell zur Erkennung von Sprachaktivität in Audio}
}
\glsadd{gls-Silero-VAD}

\newglossaryentry{gls-GPT-4o-mini}{
  type=fachwort,
  name=GPT-4o-mini,
  description={Kompaktes Large Language Model von OpenAI mit reduziertem Ressourcenbedarf}
}
\glsadd{gls-GPT-4o-mini}

\newglossaryentry{gls-Docker-Container}{
  type=fachwort,
  name=Docker-Container,
  description={Laufende Instanz eines Docker-Images, die eine isolierte Anwendungsumgebung bereitstellt}
}
\glsadd{gls-Docker-Container}

\newglossaryentry{gls-Port-Forwarding}{
  type=fachwort,
  name=Port-Forwarding,
  description={Netzwerktechnik zur Weiterleitung eingehender Verbindungen von einem Port auf einen anderen Port oder ein anderes Gerät, auch als Portweiterleitung bezeichnet}
}
\glsadd{gls-Port-Forwarding}

\newglossaryentry{gls-Portfreigabe}{
  type=fachwort,
  name=Portfreigabe,
  description={Konfiguration an Router oder Firewall zur Freigabe eines Ports für eingehende Verbindungen aus dem Internet an ein internes Gerät}
}
\glsadd{gls-Portfreigabe}

\newglossaryentry{gls-Firewall}{
  type=fachwort,
  name=Firewall,
  description={Sicherheitssystem zur Kontrolle des ein- und ausgehenden Netzwerkverkehrs basierend auf definierten Regeln}
}
\glsadd{gls-Firewall}

\newglossaryentry{gls-ulaw}{
  type=fachwort,
  name=ulaw,
  description={Audio-Codec für Sprachkompression in der Telefonie (µ-law PCM), standardisiert für nordamerikanische Telefonnetze}
}
\glsadd{gls-ulaw}

\newglossaryentry{gls-alaw}{
  type=fachwort,
  name=alaw,
  description={Audio-Codec für Sprachkompression in der Telefonie (A-law PCM), standardisiert für europäische und internationale Telefonnetze}
}
\glsadd{gls-alaw}

\newglossaryentry{gls-Opus-Codec}{
  type=fachwort,
  name=Opus,
  description={Hocheffizienter Audio-Codec für Sprach- und Musikübertragung mit niedriger Latenz, standardisiert in RFC 6716 und weit verbreitet in WebRTC}
}
\glsadd{gls-Opus-Codec}

\newglossaryentry{gls-G711}{
  type=fachwort,
  name=G.711,
  description={ITU-T Standard für Sprachcodec mit PCM-Kodierung, weit verbreitet in klassischer Telefonie mit 64 kbit/s Datenrate}
}
\glsadd{gls-G711}

\newglossaryentry{gls-Catch-all-Fallback}{
  type=fachwort,
  name=Catch-all-Fallback,
  description={Regelkonfiguration zur Behandlung aller nicht explizit definierten Fälle in einem Routing- oder Verarbeitungssystem}
}
\glsadd{gls-Catch-all-Fallback}

\newglossaryentry{gls-SIP-INVITE}{
  type=fachwort,
  name=SIP INVITE,
  description={SIP-Nachrichtentyp zum Initiieren einer Session oder Einladung zu einem Anruf}
}
\glsadd{gls-SIP-INVITE}

\newglossaryentry{gls-Dispatch-Rule}{
  type=fachwort,
  name=Dispatch Rule,
  description={Regelwerk zur Weiterleitung eingehender SIP-Anrufe an bestimmte Ziele oder Räume}
}
\glsadd{gls-Dispatch-Rule}

\newglossaryentry{gls-Cleanup-Funktion}{
  type=fachwort,
  name=Cleanup-Funktion,
  description={Programmiermuster oder Funktion zur Freigabe von Ressourcen und Bereinigung von Systemzuständen}
}
\glsadd{gls-Cleanup-Funktion}

\newglossaryentry{gls-ISO-8601}{
  type=fachwort,
  name=ISO 8601,
  description={Internationaler Standard für die Darstellung von Datum und Uhrzeit in einheitlichem Format}
}
\glsadd{gls-ISO-8601}

\newglossaryentry{gls-Turn}{
  type=fachwort,
  name=Turn,
  description={Einzelner Dialog-Schritt in einer Konversation, bestehend aus Nutzeräußerung und Systemantwort}
}
\glsadd{gls-Turn}

\newglossaryentry{gls-Word-Error-Rate}{
  type=fachwort,
  name=Word-Error-Rate,
  description={Metrik zur Bewertung der Genauigkeit von Spracherkennung, die den Anteil falsch erkannter Wörter misst}
}
\glsadd{gls-Word-Error-Rate}

\newglossaryentry{gls-Client-Initiated}{
  type=fachwort,
  name=Client-Initiated,
  description={Verbindungsabbruch oder Aktion, die vom Client ausgelöst wurde}
}
\glsadd{gls-Client-Initiated}

\newglossaryentry{gls-Paketverlust}{
  type=fachwort,
  name=Paketverlust,
  description={Verlust von Datenpaketen während der Netzwerkübertragung, der zu Qualitätsminderung führen kann}
}
\glsadd{gls-Paketverlust}

\newglossaryentry{gls-Jitter}{
  type=fachwort,
  name=Jitter,
  description={Variabilität in der Verzögerung von Netzwerkpaketen, die zu unregelmäßiger Wiedergabe führen kann}
}
\glsadd{gls-Jitter}

\newglossaryentry{gls-Echo}{
  type=fachwort,
  name=Echo,
  description={Akustisches Phänomen, bei dem der eigene Ton mit Verzögerung zurückgehört wird}
}
\glsadd{gls-Echo}

\newglossaryentry{gls-Payload}{
  type=fachwort,
  name=Payload,
  description={Nutzdaten einer Netzwerknachricht oder API-Anfrage, die die eigentlichen Informationen enthalten}
}
\glsadd{gls-Payload}

\newglossaryentry{gls-Snapshot}{
  type=fachwort,
  name=Snapshot,
  description={Momentaufnahme eines Systemzustands zu einem bestimmten Zeitpunkt, die für Backup, Wiederherstellung oder Versionierung verwendet wird}
}
\glsadd{gls-Snapshot}

\newglossaryentry{gls-Incident-Response-Plan}{
  type=fachwort,
  name=Incident-Response-Plan,
  description={Dokumentierter Prozess zur Reaktion auf Sicherheitsvorfälle und Systemausfälle}
}
\glsadd{gls-Incident-Response-Plan}

\newglossaryentry{gls-Auto-Scaling}{
  type=fachwort,
  name=Auto-Scaling,
  description={Automatische Anpassung von Ressourcen (z.B. Anzahl Container) basierend auf Systemlast}
}
\glsadd{gls-Auto-Scaling}

\newglossaryentry{gls-Load-Balancer}{
  type=fachwort,
  name=Load-Balancer,
  description={System zur Verteilung von Netzwerkverkehr auf mehrere Server zur Optimierung von Auslastung und Verfügbarkeit}
}
\glsadd{gls-Load-Balancer}

\newglossaryentry{gls-Edge-Computing}{
  type=fachwort,
  name=Edge-Computing,
  description={Bereitstellung von Rechen- und Speicherressourcen nahe am Ort der Datenentstehung zur Reduzierung von Latenzen}
}
\glsadd{gls-Edge-Computing}

\newglossaryentry{gls-Multimodal}{
  type=fachwort,
  name=Multimodal,
  description={Verarbeitung und Integration mehrerer Modalitäten (z.B. Text, Audio, Video) in einem System}
}
\glsadd{gls-Multimodal}

\newglossaryentry{gls-On-Device-Inferenz}{
  type=fachwort,
  name=On-Device-Inferenz,
  description={Ausführung von maschinellen Lernmodellen direkt auf einem Endgerät statt in der Cloud}
}
\glsadd{gls-On-Device-Inferenz}

\newglossaryentry{gls-Trace}{
  type=fachwort,
  name=Trace,
  description={Aufzeichnung des Ausführungspfads einer Anfrage durch ein verteiltes System zur Performance-Analyse}
}
\glsadd{gls-Trace}

\newglossaryentry{gls-Timeout}{
  type=fachwort,
  name=Timeout,
  description={Zeitüberschreitung, nach der eine Verbindung oder Operation abgebrochen wird}
}
\glsadd{gls-Timeout}

\newglossaryentry{gls-Latenz}{
  type=fachwort,
  name=Latenz,
  description={Zeitverzögerung zwischen Senden und Empfangen von Daten in einem System}
}
\glsadd{gls-Latenz}

\newglossaryentry{gls-Throughput}{
  type=fachwort,
  name=Throughput,
  description={Datenmenge, die pro Zeiteinheit durch ein System übertragen werden kann}
}
\glsadd{gls-Throughput}

\newglossaryentry{gls-Bitrate}{
  type=fachwort,
  name=Bitrate,
  description={Datenmenge pro Zeiteinheit bei der Übertragung digitaler Medien, gemessen in Bits pro Sekunde}
}
\glsadd{gls-Bitrate}

\newglossaryentry{gls-Abtastrate}{
  type=fachwort,
  name=Abtastrate,
  description={Anzahl der Abtastwerte pro Sekunde bei der Digitalisierung analoger Signale, gemessen in Hertz (z.B. 8 kHz, 16 kHz, 48 kHz)}
}
\glsadd{gls-Abtastrate}

\newglossaryentry{gls-Best-Effort}{
  type=fachwort,
  name=Best-Effort,
  description={Dienstmodell ohne Leistungsgarantien, bei dem das System sein Bestes versucht}
}
\glsadd{gls-Best-Effort}

\newglossaryentry{gls-Uptime}{
  type=fachwort,
  name=Uptime,
  description={Zeitdauer, während der ein System betriebsbereit und verfügbar ist}
}
\glsadd{gls-Uptime}

\newglossaryentry{gls-Failover}{
  type=fachwort,
  name=Failover,
  description={Automatisches Umschalten auf ein Backup-System bei Ausfall des primären Systems}
}
\glsadd{gls-Failover}

\newglossaryentry{gls-Single-Host-Deployment}{
  type=fachwort,
  name=Single-Host-Deployment,
  description={Bereitstellung eines Systems auf einem einzelnen physischen oder virtuellen Server}
}
\glsadd{gls-Single-Host-Deployment}

\newglossaryentry{gls-Anrufaufkommen}{
  type=fachwort,
  name=Anrufaufkommen,
  description={Anzahl der Anrufe in einem bestimmten Zeitraum}
}
\glsadd{gls-Anrufaufkommen}

\newglossaryentry{gls-Transkript}{
  type=fachwort,
  name=Transkript,
  description={Textuelle Darstellung gesprochener Sprache}
}
\glsadd{gls-Transkript}

\newglossaryentry{gls-Inferenz}{
  type=fachwort,
  name=Inferenz,
  description={Anwendung eines trainierten maschinellen Lernmodells zur Vorhersage oder Generierung von Ausgaben}
}
\glsadd{gls-Inferenz}

\newglossaryentry{gls-Quantisierung}{
  type=fachwort,
  name=Quantisierung,
  description={Reduzierung der Präzision von Modellparametern zur Verringerung von Speicher- und Rechenbedarf bei maschinellen Lernmodellen}
}
\glsadd{gls-Quantisierung}

\newglossaryentry{gls-CPaaS}{
  type=fachwort,
  name=CPaaS,
  first={Communications Platform as a Service (CPaaS)},
  text=CPaaS,
  description={Communications Platform as a Service, eine Cloud-basierte Plattform zur Bereitstellung von Kommunikationsdiensten}
}

\newglossaryentry{gls-mTLS}{
  type=fachwort,
  name=mTLS,
  first={Mutual TLS (mTLS)},
  text=mTLS,
  description={Mutual TLS, eine Erweiterung von TLS, bei der beide Seiten (Client und Server) sich gegenseitig mit Zertifikaten authentifizieren}
}

\newglossaryentry{gls-QoS}{
  type=fachwort,
  name=QoS,
  first={Quality of Service (QoS)},
  text=QoS,
  description={Quality of Service, ein Konzept zur Verwaltung von Netzwerk-Performance-Parametern wie Latenz, Bandbreite und Paketverlust}
}

\newglossaryentry{gls-WebRTC}{
  type=fachwort,
  name=WebRTC,
  first={Web Real-Time Communication (WebRTC)},
  text=WebRTC,
  description={Web Real-Time Communication, ein Standard für Echtzeitkommunikation über Browser und Apps}
}

\newglossaryentry{gls-VoIP}{
  type=fachwort,
  name=VoIP,
  first={Voice over IP (VoIP)},
  text=VoIP,
  description={Voice over IP, Technologie zur Übertragung von Sprachkommunikation über IP-basierte Netzwerke, bei der Audiosignale digitalisiert, paketvermittelt übertragen und am Ziel wieder in Sprache umgewandelt werden}
}

\newglossaryentry{gls-WAV}{
  type=fachwort,
  name=WAV,
  first={Waveform Audio File Format (WAV)},
  text=WAV,
  description={Waveform Audio File Format, ein Audio-Dateiformat zur Speicherung von Audiosignalen in digitaler Form, das häufig für die Speicherung von Musik und Sprache verwendet wird}
}


% Literatur
\usepackage[
  backend=biber,
  style=alphabetic,
  sorting=nyt,
  maxcitenames=2,
  mincitenames=1,
  maxbibnames=99
]{biblatex}
\addbibresource{references.bib}

% Metadaten
\newcommand{\ThesisTitle}{Machbarkeitsstudie zur Integration von Large Language Models in Container-basierte SIP-Telefonsysteme}
\newcommand{\ThesisAuthor}{Jan Liebeck}
\newcommand{\ThesisMatrikel}{75502}
\newcommand{\ThesisSupervisorCompany}{B.Sc. Felix Piasta (stack1 GmbH)}
\newcommand{\ThesisSupervisorUniversity}{Prof. Dr.-Ing. Robert Müller (HTWK Leipzig)}
\newcommand{\ThesisDate}{26.01.2026}

% Widow/Orphan-Kontrolle
\widowpenalty=10000
\clubpenalty=10000
\hyphenpenalty=5000

\begin{document}

% Titelblatt nach HTWK Leipzig Vorgaben
% Hochschulgerechte Formatierung mit erforderlichen Informationen

\begin{titlepage}
\centering

% Logo und Hochschule (optional)
\vspace*{1cm}

{\normalsize Hochschule für Technik, Wirtschaft und Kultur (HTWK) Leipzig\par}
{\normalsize Fakultät Informatik und Medien\par}

\vspace{2cm}

% Studiengangbezeichnung und Arbeitstyp
{\large Bachelorarbeit\par}
{\normalsize im Studiengang Medieninformatik (B.Sc.)\par}
\vspace{1cm}

{\normalsize zur Erlangung des akademischen Grades\par}
{\normalsize \textit{Bachelor of Science} (B.Sc.)\par}

\vspace{2.5cm}

% Titel der Arbeit
{\Large \bfseries \ThesisTitle\par}

\vspace{2.5cm}

% Autor
{\large Verfasst von: \ThesisAuthor\par}

\vspace{2cm}

% Betreuer und Gutachter
{\normalsize Betreuer:\par}
{\normalsize \ThesisSupervisorUniversity\par}

\vspace{0.8cm}

{\normalsize Gutachter:\par}
{\normalsize \ThesisSupervisorCompany\par}

\vfill

% Abgabedatum
{\normalsize Abgabedatum: \ThesisDate\par}

\end{titlepage}

% Neue Seite nach Titelblatt
\cleardoublepage


\pagenumbering{Roman}
\pagestyle{scrheadings}

\clearpage
\manualmark
\markboth{Abstrakt}{Abstrakt}
% Abstract für Bachelorarbeit
% Nach Schreibwerkstatt Universität Stuttgart: Prägnant, objektiv, ohne Ausblick, ohne Vollformen

\chapter*{Abstrakt}
\markboth{Abstrakt}{Abstrakt}
\thispagestyle{scrheadings}
\addcontentsline{toc}{chapter}{Abstrakt}

Klassische \gls{gls-SIP}-Telefonsysteme ermöglichen Echtzeitkommunikation, jedoch bieten sie keine intelligente Sprachverarbeitung. Diese Arbeit untersucht die technische Machbarkeit der Integration von Large Language Models in containerbasierte \gls{gls-SIP}-Telefonsysteme. Die zentrale Forschungsfrage lautet: Unter welchen Bedingungen ermöglichen containerbasierte Ansätze eine stabile Verbindung von \gls{gls-SIP}-basierter Telefonie mit \gls{gls-KI}-Sprachmodellen?

Der Prototyp verbindet drei Kernkomponenten: Asterisk übernimmt die \gls{gls-SIP}-\\Signalisierung, LiveKit verarbeitet \gls{gls-WebRTC}-Medienströme, und ein Python-Agent steuert die \gls{gls-KI}-Logik. Die Sprachverarbeitung erfolgt über OpenAI Whisper für die Spracherkennung, GPT-4o-mini für die Dialogverwaltung und OpenAI Text-to-Speech für die Sprachsynthese. Die Docker-basierte Implementierung garantiert Reproduzierbarkeit. Die iterative Entwicklung begann mit einem leichtgewichtigen \gls{gls-SIP}-Agenten für erste Schritte und das Erlernen grundlegender \gls{gls-SIP}-Mechanismen und führte schrittweise zur finalen Architektur.

Die Evaluation konzentriert sich auf die \gls{gls-SIP}-Integration. Fünf repräsentative Testanrufe wurden detailliert ausgewertet. Die Median-Latenz beträgt 0,78 s für Spracherkennung und 1,41 s für Sprachsynthese. Das System zeigt stabile bidirektionale Audioübertragung ohne Timeouts. Eine \gls{gls-STT}-Fehlinterpretation wurde dokumentiert. Die CPU-Auslastung liegt bei 11,2\%, der RAM-Bedarf zwischen 11,2 und 12,1 GB. Die Audioqualität erreicht \gls{gls-MOS}-Werte von 4,0 bis 4,5.

Die Arbeit präsentiert eine \gls{gls-ETSI}-konforme Architektur, demonstriert eine funktionsfähige Implementierung, liefert Evaluationsdaten aus \gls{gls-SIP}-Tests und formuliert Handlungsempfehlungen für kleine und mittlere Unternehmen. Open-Source-Lösungen ermöglichen Datenschutz-konforme Deployments ohne Vendor-Lock-in. Die Ergebnisse zeigen: \gls{gls-KI}-Sprachmodelle lassen sich ohne spezialisierte Hardware mit klassischen \gls{gls-VoIP}-Systemen kombinieren, um dialogorientierte Sprachverarbeitung zu ermöglichen.\\

\noindent\textbf{Schlagwörter:} Large Language Model, \gls{gls-SIP}, Telefonie, Docker-basierte Architektur, Asterisk, LiveKit, \gls{gls-KI}-Integration


\clearpage
\manualmark
\markboth{Abstract}{Abstract}
\chapter*{Abstract}
\markboth{Abstract}{Abstract}
\thispagestyle{scrheadings}
\addcontentsline{toc}{chapter}{Abstract}

Classic \gls{gls-SIP} telephone systems enable real-time communication, but they do not offer intelligent speech processing. This thesis examines the technical feasibility of integrating large language models into container-based \gls{gls-SIP} telephone systems. The central research question is: Under what conditions do container-based approaches enable a stable connection between \gls{gls-SIP}-based telephony and \gls{gls-AI} speech models?

The prototype connects three core components: Asterisk handles \gls{gls-SIP} signaling, LiveKit processes \gls{gls-WebRTC} media streams, and a Python agent controls the \gls{gls-AI} logic. Speech processing is handled by OpenAI Whisper for speech recognition, GPT-4o-mini for dialogue management, and OpenAI Text-to-Speech for speech synthesis. The Docker-based implementation guarantees reproducibility. Iterative development began with a lightweight \gls{gls-SIP} agent called baresip to evaluate basic \gls{gls-SIP} mechanisms and gradually led to the final architecture.

The evaluation focuses on \gls{gls-SIP} integration. Five representative test calls were evaluated in detail. The median latency is 0.78 s for speech recognition and 1.41 s for speech synthesis. The system shows stable bidirectional audio transmission without timeouts. One \gls{gls-STT} misinterpretation was documented. \gls{gls-CPU} utilization is 11.2\%, and \gls{gls-RAM} requirements are between 11.2 and 12.1 GB. Audio quality achieves \gls{gls-MOS} values of 4.0 to 4.5.

The paper presents an \gls{gls-ETSI}-compliant architecture, demonstrates a functional implementation, provides evaluation data from \gls{gls-SIP} tests, and formulates recommendations for action for small and medium-sized enterprises. Open source solutions enable data protection-compliant deployments without vendor lock-in. The results show that \gls{gls-AI} language models can be combined with standard \gls{gls-VoIP} systems without specialized hardware to enable dialogue-oriented speech processing.

\noindent\textbf{Keywords:} Large Language Model, \gls{gls-SIP}, telephony, Docker-based architecture, Asterisk, LiveKit, \gls{gls-AI} integration


\clearpage
\automark{chapter}

\ihead{\small\leftmark}
\tableofcontents
\addtocontents{toc}{\protect\thispagestyle{scrheadings}}
\listoftables
\addtocontents{lot}{\protect\thispagestyle{scrheadings}}
\printglossary[type=fachwort, title={Glossar}, nonumberlist]

\clearpage
\pagenumbering{arabic}

\chapter{Einleitung}
\label{chap:einleitung}

Moderne Technologien wie \gls{gls-KI} und die Integration in Telefonsysteme führen dazu, dass sich sowohl die Art der Interaktion als auch die Qualität und Komplexität dieser Systeme weiterentwickeln. Sprachassistenten ermöglichen die automatische Erkennung menschlicher Sprache und können beispielsweise individuelle Funktionen zur Steuerung von Smart-Home-Komponenten übernehmen. Die Integration von \gls{gls-LLM} erweitert diese Funktionen um dialogorientierte Steuerung und Kontextverarbeitung über mehrere Gesprächsbeiträge hinweg. Dadurch werden natürlichere und flexiblere Dialoge möglich, die klassische Telefonie sowie moderne Sprachplattformen ergänzen und neue Formen der Mensch-Maschine-Kommunikation schaffen \parencite{Schmiedchen2026KI}. Hierbei stellt sich die Frage, wie \gls{gls-KMU} mit begrenzter Mitarbeiter- und Umsatzgrö\ss e von solchen Technologien profitieren können.

\section{Problemstellung und Zielsetzung}
\gls{gls-VoIP}-Systeme, die auf dem \gls{gls-SIP}-Protokoll basieren, sind häufig stark an spezielle Telefonie-Hardware gebunden und bringen kaum intelligente Sprachverarbeitung mit. Die Verbindung solcher Systeme mit einem \gls{gls-LLM} eröffnet neue Möglichkeiten, etwa für automatisierte Sprachdialoge, die Kundensysteme effizient unterstützen und
gleichzeitig natürlich wirken. Forschung zur Integration von Sprachverarbeitung mit \gls{gls-LLM} zeigt, dass die Kombination von Audio-Encodern mit Sprachmodellen robuste multilinguale Spracherkennung ermöglicht \parencite{whisper_quant_2025}. Au\ss erdem wird deutlich, dass sich klassische \gls{gls-SIP}-basierte Telefonsysteme mit modernen \gls{gls-KI}-Dialogkomponenten und containerbasierten Architekturen kombinieren lassen \parencite{Antonova2025}. In der Praxis zeigt sich jedoch, dass die Übertragung und Verarbeitung von Audio zwischen virtualisierten Systemkomponenten eine gro\ss e Herausforderung darstellt. In containerbasierten Umgebungen, in denen verschiedene Dienste getrennt voneinander laufen, müssen Sprachdaten zuverlässig und mit möglichst niedriger Latenz zwischen diesen Diensten übertragen werden \parencite{ReliabilityVirtualizedNetworks}.

Neuere Studien weisen darauf hin, dass \gls{gls-SIP}-Instanzen und die dazugehörigen Media-Proxy in Container-Umgebungen eingebunden werden sollten. Ein Media-Proxy leitet Audio- und Videoströme weiter, passt sie bei Bedarf an oder transkodiert sie, damit unterschiedliche Protokolle miteinander kompatibel sind. Dabei ist es wichtig, dass die Verwaltung der Dienste automatisiert und orchestriert erfolgt, um Skalierbarkeit und Fehlertoleranz sicherzustellen \parencite{Antonova2025}. Die Realisierung solcher Systeme auf Basis von Open-Source-Softwarekomponenten und Container-Orchestrierung mit Docker ermöglicht eine reproduzierbare, kosteneffiziente und flexible Bereitstellung. Im Rahmen dieser Arbeit wurde durch eine iterative Evaluierung verschiedener \gls{gls-SIP}-Implementierungen eine praxisnahe Architektur entwickelt. Ausgehend von baresip als experimentellem Tool zur Überprüfung der SIP-Registrierung und Protokollkompatibilität wurde die Architektur schrittweise mit Asterisk als stabilem \gls{gls-SIP}-Backend und LiveKit als Integrations-Framework für \gls{gls-KI}-Sprachverarbeitungsdienste realisiert.

Das \gls{gls-RTP} übernimmt die Echtzeitübertragung von Audio und Video, während \gls{gls-WebRTC} die Kommunikation im Browser ermöglicht. Für die Medienverarbeitung ist eine RTP-Bridge erforderlich, die Mediendaten zwischen den Diensten weiterleitet, anpasst oder transkodiert. Sicherheitsfunktionen wie das \gls{gls-SRTP}, zur Verschlüsselung der Mediendaten und \gls{gls-DTLS} für die sichere Schlüsselverteilung, schützen die übertragenen Daten. Die Interoperabilität zwischen \gls{gls-SIP}/\gls{gls-RTP} und \gls{gls-WebRTC} stellt sicher, dass unterschiedliche Clients zuverlässig miteinander kommunizieren können. Ohne diese Komponenten lassen sich stabile, bidirektionale Audio-Pipelines mit kontrollierter Latenz kaum realisieren\\ \parencite{Antonova2025, ReliabilityVirtualizedNetworks}.

\section{Aufbau der Arbeit}
 Kapitel \ref{chap:grundlagen} beschreibt die technischen Grundlagen zu allen genutzten Technologien und Protokollen. Anschlie\ss end enthält Kapitel \ref{chap:anforderungen} Anforderungen für die Implementierung. Das Kapitel \ref{chap:architektur} behandelt Architekturthematiken und Kapitel \ref{chap:marktanalyse} gibt Einblicke in den Markt von \gls{gls-KI}-basierten Systemen. Im nächsten Kapitel \ref{chap:implementierung} werden die prototypische Umsetzung und die Integration der zentralen Komponenten dokumentiert. Danach präsentiert Kapitel \ref{chap:evaluation} die Messergebnisse und deren Auswertung.
Das abschlie\ss ende Kapitel \ref{chap:diskussion} vereint die Interpretation der Ergebnisse mit den zentralen Erkenntnissen.

\section{Beitrag der Arbeit}
Der Beitrag der Arbeit umfasst: (1) eine Architektur für die Kombination von \gls{gls-SIP}-/\gls{gls-VoIP}-Systemen mit einem \gls{gls-LLM}-Dienst in virtualisierten Containern, (2) eine prototypische Implementierung der genutzten Komponenten mit reproduzierbarem Setup, (3) eine Evaluation wesentlicher Qualitätsmerkmale und (4) einen praxisorientierten Leitfaden mit technischen Handlungsempfehlungen zur Nutzung des prototypischen Systems. \\

\section{Methodik}

Bei der Realisierung dieser Arbeit wurde ein iterativer und experimenteller Ansatz verfolgt. Der Fokus lag auf dem Nachweis der technischen Machbarkeit der Integration von \gls{gls-SIP}-Systemen mit \gls{gls-KI}-Sprachverarbeitung. Baresip\footnote{Offizielle Webseite: \url{https://github.com/baresip/baresip}} diente als experimentelles Tool zur Überprüfung grundlegender \gls{gls-SIP}-Registrierungs- und Protokoll-Kompatibilitätsfragen. Aufgrund der begrenzten Dialplan-Logik wurde zügig auf Asterisk gewechselt, um die grundlegenden Mechanismen einer \gls{gls-SIP}-Trunk-Registrierung in einer produktionsnahen Umgebung abzubilden. Diese Phase führte zur Auswahl von Asterisk als stabiles und produktiv einsetzbares \gls{gls-SIP}-Backend. Das LiveKit-Framework\footnote{Offizielle Webseite: \url{https://livekit.io/}} wurde anschlie\ss end als geeignete Lösung zur Integration mit \gls{gls-KI}-Sprachverarbeitungs-\\diensten implementiert. Da der Schwerpunkt auf der finalen, funktionsfähigen Architektur liegt, werden die Evaluierungsphasen der einzelnen Komponenten nur so weit dokumentiert, wie sie für das Verständnis des Gesamtsystems notwendig sind. Außerdem wurde eine \gls{gls-UI} als praktisches Steuerungstool für reine \gls{gls-WebRTC} Telefonie ohne \gls{gls-SIP} Integration entwickelt, um die grundlegenden \gls{gls-KI}-Komponenten zu prüfen. Die wissenschaftliche Evaluation konzentriert sich jedoch auf die \gls{gls-SIP}-basierte Telefonie-Integration mit \gls{gls-KI}-Sprachverarbeitung, nicht auf die \gls{gls-UI}.

Folgende Hypothesen werden überprüft und im Evaluationskapitel \ref{chap:evaluation} anhand systematischer Tests bestätigt oder verworfen.

\paragraph{H1 (Funktionalität):} Es wird angenommen, dass die bidirektionale Audioübertragung zwischen einem \gls{gls-SIP}-Telefonsystem und einem extern gehosteten \gls{gls-KI}-\\Sprachverarbeitungsdienst in einer containerbasierten Umgebung stabil aufrechterhalten werden kann, ohne dabei systematische Verbindungsabbrüche, Codec-\\Inkompatibilitäten oder Medienpfadfehler aufzuweisen.\\

\textit{Kontext:} Nach RFC~3261 regelt \gls{gls-SIP} die Signalisierung und Codec-Verhandlung in \gls{gls-VoIP}-Systemen. Die Standards RFC~5245 (Interactive Connectivity Establishment), RFC~5389 (Session Traversal Utilities for NAT) und RFC~5766 (Traversal Using Relays around NAT) beschreiben Mechanismen zur Überwindung von Netzwerk-Adressumsetzung und Firewall-Hindernissen, die für die Herstellung stabiler Audio-Kommunikationspfade essentiell sind \parencite{rfc3261,rfc5245,rfc5389,rfc5766}.

\paragraph{H2 (Latenz):} Es wird angenommen, dass die in der Evaluation gemessenen Latenzbeiträge der Sprachverarbeitung (\gls{gls-STT}‑Latenz, Transkriptverzögerung, \gls{gls-TTS}‑Latenz) im praxistauglichen Bereich liegen und interaktive Kommunikation unterstützen.

\textit{Kontext:} In der klassischen Telefonie nach ITU‑T G.114 werden One‑Way‑Latenzen bis 150 ms als akzeptabel, bis 400 ms als noch tolerierbar und oberhalb von 400 ms als kritisch eingestuft \parencite{ITU-T-G114}. In dieser Arbeit werden die Latenzbeiträge der Sprachverarbeitung separat betrachtet, da sie die zentralen, systemseitig beeinflussbaren Schritte abbilden (\gls{gls-STT}, Transkriptverzögerung in der \gls{gls-SIP}-Pipeline, \gls{gls-TTS}). Externe Anteile (Routing, Netz, \gls{gls-SIP}‑Signalisierung) wirken zusätzlich, werden hier aber nicht gemessen.

\paragraph{H3 (Ressourceneffizienz):} Es wird angenommen, dass die eingesetzte containerbasierte Architektur auf Basis von Standard-Serverhardware betrieben werden kann und dabei Systemressourcen (\gls{gls-CPU}, \gls{gls-RAM}, Speicherbandbreite) effizient ausnutzt, sodass keine übermä\ss ige Überprovisionierung erforderlich ist.

\textit{Kontext:} Nach den \gls{gls-ETSI}-\gls{gls-NFV}-Richtlinien besteht das Ziel der Netzwerk-\\Funktionsvirtualisierung darin, spezialisierte Hardware durch software- und containerbasierte Dienste zu ersetzen, die auf Standard-Servern betrieben werden können. Praktische Studien zu solchen Kommunikationssystemen zeigen, dass mit optimierter Dimensionierung stabile und performante Deployments auch auf ressourcenlimitierter Hardware möglich sind \parencite{ETSI-NFV,ETSI-REL,Antonova2025}.

\paragraph{H4 (Skalierbarkeit):} Es wird angenommen, dass die containerbasierte Architektur die Replikation und horizontale Skalierung von Komponenteninstanzen unterstützt, sodass das System bei wachsendem Anrufaufkommen durch Hinzufügen weiterer Container-Replicas skalierbar bleibt.

\textit{Kontext:} Containerorchestrierungsplattformen ermöglichen eine automatisierte Bereitstellung, Scaling und Loadbalancing. Nach \gls{gls-ETSI}-\gls{gls-NFV}-Standard können Netzwerkfunktionen vertikal und horizontal skaliert werden, um dynamisch auf sich ändernde Lasten zu reagieren. Dies wird durch eine dezentrale Komponentenarchitektur und zustandslose Services realisiert \parencite{ETSI-NFV,ETSI-REL,Antonova2025}.

\paragraph{H5 (STT-Robustheit):} Es wird angenommen, dass die eingesetzte Spracherkennungskomponente deutschsprachige Telefonie-Prompts und typische Anrufervorgaben zuverlässig transkribiert, ohne dabei systematische oder regelmä\ss ig wiederholte Erkennungsfehler zu erzeugen.

\textit{Kontext:} Moderne Spracherkennungsmodelle wie OpenAI Whisper sind auf multilingualen Trainingsdaten trainiert und zeigen besonders für Deutsch eine hohe Erkennungsgenauigkeit. Dies stützt die Annahme, dass für telefonische Anwendungen mit kontrollierten Prompt-Szenarien eine zuverlässige Performance zu erwarten ist, auch wenn Test-Bedingungen und reale Produktionsumgebungen unterschiedliche Fehlerraten aufweisen können \parencite{whisper_quant_2025}.

\chapter{Grundlagen}
\label{chap:grundlagen}

Bevor ein \gls{gls-KI}-basiertes Telefonsystem entwickelt wird, lohnt sich der Blick auf Container-Umgebungen: Eine \gls{gls-VM}, die einen eigenen Gast-Kernel verwendet, nutzen im Gegensatz dazu Container den Host-Kernel direkt und sparen so Hypervisor-Overhead \parencite{Antonova2025}. Nach geltenden Richtlinien ist für die Virtualisierung von Systemkomponenten nicht zwingend eine vollständige \gls{gls-VM} erforderlich. Stattdessen können auch Container oder andere Isolationsmechanismen als leichtgewichtige Alternativen eingesetzt werden, solange sie ausreichende Isolation bieten und effiziente Snapshots unterstützen \parencite{ETSI-REL}. In Cloud-nativen Architekturen bilden Container und zusätzliche Orchestrierung den Standard, weil Bereitstellung, Skalierung und Resilienz damit automatisiert werden \parencite{ReliabilityVirtualizedNetworks}.

\section[Container, Isolation und
Deployment]{Container, Isolation und
Bereitstellung}

Dieser Ansatz macht Container deutlich leichter und schneller als \gls{gls-VM}. Anwendungen starten innerhalb von Millisekunden und
benötigen weniger Speicher. Gerade in heterogenen Umgebungen oder bei begrenzter Hardware lassen sich so viele Dienste effizient parallel betreiben. Dies wird auch in Arbeiten zu containerbasierten Telefonsystemen bestätigt, die zeigen, dass Container auf ressourcenschwacher Hardware meist stabiler und performanter laufen als klassische \gls{gls-VM} \parencite{Antonova2025}.
Für die Entwicklung komplexerer Systeme wie \gls{gls-VoIP}-Umgebungen oder \gls{gls-KI}-basierte Dienste sind zudem die Eigenschaften von Containerplattformen wie Docker relevant. Docker-Images bestehen aus mehreren Schichten, wobei jede Schicht eine Änderung am Dateisystem repräsentiert. Diese Schichten werden inkrementell zu einem vollständigen Image zusammengesetzt, was das Verteilen und Aktualisieren vereinfacht. Durch dieses Schichtmodell können identische Basisschichten von mehreren Images gemeinsam genutzt werden, was Speicher spart und das Caching sowie die Wiederverwendbarkeit fördert. Sicherheitsupdates oder Änderungen an einzelnen Komponenten lassen sich gezielt und effizient ausrollen, ohne das gesamte Image neu erstellen zu müssen. In der Praxis führt dieser modulare Aufbau dazu, dass Systeme flexibel erweitert, skaliert und besser gewartet werden können \parencite{docker-docs}.

Gleichzeitig muss man die Grenzen dieses Ansatzes kennen: Da Container sich den Kernel teilen, hängt ihre Sicherheit stark von dessen Stabilität ab. Fehler oder Schwachstellen im Kernel können sich auf alle laufenden Container auswirken. Diese Einschränkung spielt besonders dann eine Rolle, wenn sensible Daten wie Sprachdaten verarbeitet werden. Um eine vollständig sichere und stabile, bidirektionale Audioübertragung zu gewährleisten, müssen die Dienste in Containern einwandfrei miteinander kommunizieren und die Container fehlerfrei miteinander harmonieren \parencite{Antonova2025,ETSI-REL,ReliabilityVirtualizedNetworks}.

\vspace{1em}

\section[Klassische Telefonieprotokolle]{Klassische Telefonieprotokolle}

Um zu verstehen, wie sich \gls{gls-KI}-Komponenten in ein Telefonsystem integrieren lassen, ist ein Blick auf die technischen Grundlagen der Internet-Telefonie notwendig. \gls{gls-VoIP} stellt dabei die Standardtechnologie zur Übertragung von Sprachdaten über IP-Netzwerke dar. Der zugrunde liegende Protokollstapel basiert auf dem Zusammenspiel mehrerer standardisierter Protokolle \parencite{rfc3261,ITU-T-G114}.

Für den Rufaufbau und die Signalisierung wird hauptsächlich das \gls{gls-SIP}-Protokoll eingesetzt. Es handelt sich dabei um ein Standardprotokoll zur Steuerung von Sprach- und Videokommunikation \parencite{rfc3261}. \gls{gls-SIP} regelt unter anderem die Registrierung von Endgeräten sowie den Aufbau und Abbau von Sprachverbindungen über einen\newline INVITE/ACCEPT/BYE-Handshake sowie Registrierungsmechanismen zur Anmeldung von Clients an Proxy-Servern. Die eigentlichen Audiodaten werden anschlie\ss end über \gls{gls-RTP} übertragen, das für die Echtzeitübertragung von Audio- und Videodaten ausgelegt ist. Ergänzend liefert \gls{gls-RTCP} als Begleitprotokoll zu \gls{gls-RTP} Informationen über die Verbindungsqualität und stellt Feedback zur laufenden Sitzung bereit \parencite{rfc3550}.

Open-Source-Systeme wie Asterisk \footnote{Offizielle Webseite: \url{https://www.asterisk.org/}} oder Kamailio \footnote{Offizielle Webseite: \url{https://www.kamailio.org/w/}} setzen konsequent auf diese Standards. Diese Protokollstandards finden sich beispielsweise in einer \gls{gls-PBX} wieder. Ein \gls{gls-PBX} ist eine Telefonanlage für interne und externe Sprachkommunikation, die Funktionen wie Anrufweiterleitung und Konferenzschaltungen bereitstellt. Als bewährte Open-Source-Lösungen bieten solche Systeme vielfältige Möglichkeiten zur Anrufverarbeitung, die von einfachen Durchwahlszenarien über \gls{gls-IVR} mit automatisierten Sprachmenüs bis hin zu komplexeren Konferenzfunktionen reichen \parencite{Antonova2025}. In modernen Architekturen werden \gls{gls-PBX}-Systeme häufig in containerbasierten Umgebungen und häufig in Verbindung mit Cloud-Diensten. betrieben. Der \gls{gls-PBX}-Kern, der \gls{gls-SIP}-Proxy, eine Datenbank und gegebenenfalls ein \gls{gls-RTP}-Proxy laufen dabei als separate Dienste, wodurch sich die Skalierbarkeit und Fehlertoleranz erhöht \parencite{Antonova2025,ReliabilityVirtualizedNetworks}.

Für containerbasierte Umgebungen ist das \gls{gls-NAT}-Traversierung besonders relevant, da die Echtzeitkommunikation auf bidirektionale und dynamische Verbindungen angewiesen ist. \gls{gls-SIP} wurde ursprünglich für Netzwerke mit öffentlich erreichbaren Endpunkten konzipiert und berücksichtigt keine durch \gls{gls-NAT} verursachte Adressumsetzung \parencite{rfc3261}. Insbesondere bei der Aushandlung von Medientransporten werden IP-Adressen und Ports im Protokollinhalt übertragen, die aus der Sicht des \gls{gls-NAT} häufig nicht erreichbar sind, was den direkten Austausch von Audio- und Videodaten verhindert.

Zur Überwindung dieser Einschränkungen kommen zusätzliche \gls{gls-NAT}-Traversierungsmechanismen zum Einsatz. Das \gls{gls-STUN}-Protokoll ermöglicht Endpunkten die Ermittlung ihrer von au\ss en sichtbaren Transportadressen \parencite{rfc5389}. Das \gls{gls-TURN}-Verfahren stellt als Fallback eine relaybasierte Weiterleitung von Mediaströmen bereit, wenn direkte Verbindungen aufgrund restriktiver \gls{gls-NAT}-Typen nicht möglich sind \parencite{rfc5766}. Aufbauend darauf koordiniert \gls{gls-ICE} diese Verfahren, indem verschiedene potenzielle Verbindungspfade systematisch getestet werden, um eine funktionierende Ende-zu-Ende-Verbindung auszuhandeln \parencite{rfc5245}.

In containerbasierten Umgebungen wie Kubernetes\footnote{Offizielle Webseite: \url{https://kubernetes.io/}} oder Docker\footnote{Offizielle Webseite: \url{https://www.docker.com/}} wird diese Problematik durch zusätzliche virtuelle Netzwerkebenen weiter verschärft, da Container in der Regel über mehrere Adressübersetzungen hinweg kommunizieren müssen, was NAT-Traversierung zu einer zentralen technischen Herausforderung für die Echtzeitkommunikation in Telefonsystemen macht\parencite{Antonova2025,ReliabilityVirtualizedNetworks}.

\vspace{1em}

\section{Audio-Formate und Transkodierungs-Pipeline}
Die Kopplung klassischer \gls{gls-SIP}-Telefonie mit einem \gls{gls-LLM}-\gls{gls-API} bringt abweichende Audio-Formate zusammen: Im \gls{gls-PSTN} wird G.711-\gls{gls-PCM} mit 8 kHz und 8 Bit genutzt, wie im \gls{gls-RTP}-Profil beschrieben \parencite{rfc3551,ITU-T-G114}. Browserbasierte \gls{gls-WebRTC}-Pfade setzen hingegen auf den Opus-Codec mit interner 48 kHz Abtastung und variabler Bitrate \parencite{rfc6716}. Für die anschlie\ss ende Spracherkennung verlangen aktuelle \gls{gls-ASR}-Modelle wie Whisper \gls{gls-PCM16}-Audio bei 16 kHz \parencite{whisper_quant_2025}. Daraus entsteht eine zwingende Transcoding-Kette:

$\text{G.711 (8 kHz)} \xrightarrow{\text{\gls{gls-SIP} Bridge}} \text{Opus (48 kHz)} \xrightarrow{\text{Agent Worker}} \text{PCM16 (16 kHz)}$

Jede Stufe kann Verzögerung und mögliche Codec-Artefakte hinzufügen. Um den Overhead zu begrenzen, wird \gls{gls-VAD} eingesetzt: Sowohl das \gls{gls-RTP}-Profil (z.B. Annex B bei G.729) empfiehlt \gls{gls-VAD} zur Unterdrückung von Stillepaketen \parencite{rfc3551}, als auch der Opus-Pfad nutzt \gls{gls-VAD}-Metriken, um Sprachaktivität zu erkennen \parencite{rfc6716}. Durch das Ausblenden von Stille sinkt der Daten- und Token-Verbrauch spürbar, was sich insbesondere bei Cloud-\gls{gls-API} in geringeren Kosten niederschlägt.

Moderne \gls{gls-LLM}-\gls{gls-API} unterstützen Streaming-Antworten, sodass eine Ausgabe bereits gesendet wird, während das Modell noch restliche Tokens erzeugt. Whisper-Streaming liefert Live-Transkripte mit wenigen Sekunden Latenz und erlaubt Puffer-Feintuning für kontinuierliche Ausgaben \parencite{whisper_quant_2025}.

Damit sinkt die wahrgenommene Wartezeit, aber physische Verzögerungen bleiben: Die Kodierung selbst bringt Rahmen- und Vorlaufzeiten mit sich, und ein Jitter-Puffer in paketbasierten Netzen addiert im Mittel etwa die halbe erwartete Verzögerungsschwankung \parencite{rfc3550,ITU-T-G114}. Netzwerk-Laufzeiten und die Modellinferenz kommen zusätzlich hinzu, sodass trotz Streaming weiterhin spürbare Latenzquellen bestehen.

Es ist dabei wichtig zu unterscheiden: Während ITU-T G.114 für klassische Telefonie Einwegverzögerungen von rund 150 Millisekunden als unkritisch definiert \parencite{ITU-T-G114}, ist Whisper-Streaming für Echtzeitanwendungen konzipiert, obwohl es bei Live-Transkriptionen mit Verzögerungen von etwa 3,3 Sekunden arbeitet, wobei die Spracherkennung fortlaufend während des Sprechens erfolgt \parencite{whisper_quant_2025}. Daraus lässt sich ableiten, dass \gls{gls-KI}-basierte Dialoge Latenzen im unteren Sekundenbereich noch als natürlich erscheinen lassen, obwohl sie deutlich über den Limits der klassischen IP-Telefonie liegen.

Der zentrale Vorteil containerbasierter Architekturen wird deutlich, wenn \gls{gls-KI}-Dienste mit klassischen \gls{gls-SIP}-/\gls{gls-VoIP}-Systemen kombiniert werden: Während \gls{gls-SIP} weiterhin Signalisierung und Gesprächssteuerung übernimmt, können die Medienströme parallel zu \gls{gls-ASR}- und \gls{gls-LLM}-Services geleitet werden. Jede Komponente lässt sich dann unabhängig skalieren. Diese modulare Entkopplung ermöglicht es, Hochlast-Komponenten gezielt zu vergrö\ss ern, ohne das gesamte System zu überprovisionieren \parencite{Antonova2025,ReliabilityVirtualizedNetworks}.

\section{Rechtliche Aspekte und Datenschutz}

Da Sprachdaten personenbezogene Informationen enthalten, müssen vor allem in \gls{gls-KI}-basierten Telefonsystemen wichtige datenschutzrechtliche Grundlagen, insbesondere die \gls{gls-DSGVO}, zwingend berücksichtigt werden. Solche Systeme können nicht nur Identitätsinformationen, sondern auch sensible Zusatzdaten (z.B. Gesundheitsinformationen, Standorte, persönliche Präferenzen) nutzen und erfordern damit besonderen Schutz \parencite{Data2025_AIDataQuality}.

Kurzbefunde aus aktueller Forschung untermauern diese Sachlage und stellen heraus, dass die Verlagerung traditioneller \gls{gls-PBX}-Funktionen in containerbasierte Umgebungen zwar Flexibilität und Skalierbarkeit bringt, zugleich aber spezielle Sicherheits- und Orchestrierungsanforderungen (z.B. Netzwerksegmentierung, Schlüsselmanagement, kontrollierte Ingress-Punkte) notwendig macht, um personenbezogene Medienströme angemessen zu schützen \parencite{Antonova2025}. Darüber hinaus betonen Übersichtsarbeiten zur Zuverlässigkeit virtualisierter Netze, dass Virtualisierung die Interdependenzen zwischen Komponenten erhöht und damit neue Fehler und
Ausfallmodi einführt, die in Sicherheits- und
Datenschutzbewertungen explizit berücksichtigt werden müssen \parencite{ReliabilityVirtualizedNetworks}. Arbeiten zur Data-Governance weisen explizit auf die Bedeutung klarer Regelungen zur Datenlokalität, Retention-Politik und zur Nutzung von Subprozessoren hin, weil diese Punkte die rechtliche Bewertung von cloudbasierten \gls{gls-LLM}-Diensten ma\ss geblich beeinflussen \parencite{Data2025_AIDataQuality}.

Wesentliche rechtliche und ethische Prinzipien, die für\gls{gls-KI}-basierte Systeme gelten, basieren auf der \gls{gls-DSGVO} und
einschlägigen Data-Governance-Standards:

\noindent\textbf{Datenminimierung}
\begin{quote}
Es dürfen nur die unbedingt notwendigen Audiodaten und Metadaten verarbeitet werden.
\end{quote}
\noindent\textbf{Zweckbindung}
\begin{quote}
Aufzeichnung, Analyse und
Weiterverarbeitung sind auf den konkreten Zweck zu beschränken und vertraglich festzulegen.
\end{quote}

\noindent\textbf{Transparenz}
\begin{quote}
Betroffene sind über Art, Umfang und
Zweck der Verarbeitung zu informieren. Bei Telefonie ist dies in der Regel über Ansagen oder Datenschutzhinweise zu realisieren.
\end{quote}
\clearpage
\noindent\textbf{Integrität und Vertraulichkeit}
\begin{quote}
Technische Ma\ss nahmen (Verschlüsselung,
Zugriffskontrollen) müssen sicherstellen, dass Unbefugte keinen Zugriff erhalten.
\end{quote}

\noindent\textbf{Rechtmä\ss igkeit}
\begin{quote}
Verarbeitung muss auf einer Rechtsgrundlage beruhen (Einwilligung, Vertragserfüllung, berechtigtes Interesse, etc.).
\end{quote}

Diese Prinzipien lehnen sich an Artikel 5 der \gls{gls-DSGVO} an und werden durch aktuelle Data-Governance-Arbeiten konkretisiert \parencite{Data2025_AIDataQuality}.

\subsubsection{Konkrete Schutzma\ss nahmen}

Für den praktischen Schutz personenbezogener Mediendaten empfiehlt sich ein mehrschichtiges Vorgehen, das technische Härtung und organisatorische Prozesse kombiniert. Zunächst ist die strikte Trennung von Medien- und
Steuerkanälen zentral: Medienströme \gls{gls-RTP}/\gls{gls-SRTP} sollten in eigenen Netzwerkzonen oder in einem \gls{gls-VLAN} geführt werden. Dies sind logische Netzwerk-Segmente zur Isolation. In containerbasierten Umgebungen wie Kubernetes sind separate Namespaces mit expliziten \texttt{NetworkPolicy}-Regeln von Vorteil, um eine logische Trennung vorzunehmen. Ergänzend reduzieren dedizierte Ingress-Controller und restriktive Firewall-Policies die Angriffsfläche und begrenzen seitliche Übersprünge im Falle einer Kompromittierung. Dies entspricht den Empfehlungen für containerbasierte \gls{gls-PBX}-Bereitstellungen \parencite{Antonova2025}.

Ebenso wichtig ist ein zentrales und automatisiertes Geheimnis- und
Schlüsselmanagement: \gls{gls-API}-Keys, \gls{gls-TLS}-Zertifikate und Provider-Tokens dürfen nicht unverschlüsselt in Container-Images oder Quell-Repositories liegen. Stattdessen sollten Geheimnisse in einem gesicherten Speicher verwaltet, Kubernetes-Schlüssel nur verschlüsselt im Ruhezustand gehalten und Zugriffsrechte nach dem Least-Privilege-Prinzip vergeben werden. Kurzlebige Zugriffstokens und automatische Rotation vermindern das Risiko gestohlener Zugangsdaten. Antonova et al. heben genau diese Praxis für Asterisk-Deployments in Kubernetes hervor \parencite{Antonova2025}.

Observability und Logging müssen datenschutzgerecht gestaltet sein: Logs sind für Debugging, Monitoring und Forensik unverzichtbar, dürfen jedoch keine vollständigen Transkripte personenbezogener Gespräche in allgemeinen Log-Stores enthalten. Stattdessen empfiehlt sich die Ma\ss nahme des Maskierens oder Entfernens von personenbezogenen Informationen, ein separat verschlüsseltes Audit-Repository für sensible Ereignisse sowie klar definierte Aufbewahrungs- und Löschfristen, inklusive restriktiver Zugriffskontrollen. Diese technischen Praktiken sind notwendig, um Compliance-Anforderungen und Datenschutzverpflichtungen einhalten zu können \parencite{Data2025_AIDataQuality}.

Die Gewährleistung von Verfügbarkeit erfordert ebenfalls explizite Strategien: ReplicaSets und Auto-Scaling erhöhen die Ausfallsicherheit, lösen aber nicht automatisch Konsistenz- oder Persistenzfragen für zustandsbehaftete Daten (z.B. den \gls{gls-CDR} oder Ticket-Datenbanken). Deshalb sind regelmä\ss ige, getestete Backups (z.B. Snapshotting von Persistent-Volumes), dokumentierte Restore-Prozeduren und konsistente Sicherungsstrategien unabdingbar. Die Literatur zur Zuverlässigkeit virtualisierter Netze beschreibt gängige Modellierungsansätze und betont Testbarkeit und Wiederherstellbarkeit als zentrale Anforderungen \parencite{ReliabilityVirtualizedNetworks}.

Schlie\ss lich müssen technische Ma\ss nahmen durch betriebliche Verfahren ergänzt werden: Privacy-Impact-Assessments, regelmä\ss ige Sicherheits- und Compliance-Audits, dokumentierte Incident-Response-Pläne sowie zielgerichtete Schulungen für Entwickler und Operatoren schaffen die organisatorische Grundlage, damit die technischen Kontrollen wirksam bleiben. Aktuelle Data-Governance-Publikationen und Arbeiten zu Governance-Fragen bei \gls{gls-LLM}-Pipelines empfehlen genau diese Kombination aus Technik und Prozesssteuerung \parencite{Data2025_AIDataQuality,ssrn5348245}.

Neben Privacy-Impact-Assessments empfiehlt die Literatur die Durchführung von \gls{gls-EIA}, um ethische Risiken und Auswirkungen systematisch zu identifizieren und zu steuern. Zentrale Empfehlungen umfassen zudem die Benennung spezifischer Verantwortlicher, die kontinuierliche Überwachung von Schutzmaßnahmen, die Verbreitung von Best Practices sowie gezielte Schulungen für Mitarbeitende. Insbesondere für Hochrisiko-\gls{gls-KI}-Systeme werden Transparenz und Erklärbarkeit als zentrale Anforderungen hervorgehoben, um Vertraulichkeit und Nachvollziehbarkeit zu gewährleisten \parencite{Schmiedchen2026KI}.

\subsubsection{Datenschutzrisiken bei Einsatz von \gls{gls-LLM}}

Viele \gls{gls-LLM}- und Sprach-\gls{gls-API} werden als Cloud-Services angeboten. Dies vereinfacht eine Integration und den Betrieb, bringt jedoch spezifische Risiken mit sich, die einer systematischen Analyse bedürfen. Nach dem JusGov-Rahmenwerk gelten Aspekte wie Notwendigkeit, Proportionalität und Risikobeurteilung als zentral für eine \gls{gls-DPIA} nach Artikel 35 der \gls{gls-GDPR} \parencite{ssrn5348245}.\\Im Kontext von Cloud-\gls{gls-LLM}-Diensten sind folgende Kernrisiken relevant:
\clearpage
\begin{description}
\item[Datenübermittlung in Drittstaaten]
Es besteht das Risiko, dass Audiodaten und daraus erzeugte Transkripte bei der Nutzung cloudbasierter \gls{gls-KI}-Dienste au\ss erhalb der EU verarbeitet oder durch Subprozessoren in Drittstaaten weitergegeben werden. Fehlen geeignete Garantien, wie etwa Zertifizierungen nach \gls{gls-ISO}~27001 oder vertragliche Absicherungen, ergeben sich erhöhte datenschutzrechtliche Risiken im Sinne der \gls{gls-DSGVO} \parencite{Data2025_AIDataQuality}.

\item[Begrenzte Einsicht und Auditierbarkeit]
Die interne Verarbeitungspipeline eines Anbieters entzieht sich direkter Kontrolle. Risiken entstehen durch mangelnde Transparenz in der Speicherung, Löschung und Nachvollziehbarkeit von personenbezogenen Daten \parencite{ssrn5348245}.

\item[Unkontrollierte Datenspeicherung]
Anbieter können Nutzereingaben (Audio, Transkripte, Metadaten) zu Monitoring- oder Qualitätszwecken protokollieren oder über vereinbarte Zeiträume hinaus speichern. Ohne explizite vertragliche Aufbewahrungsfristen besteht das Risiko unerwünschter Persistenz und potenzieller Zweckentfremdung \parencite{Data2025_AIDataQuality}.

\item[Bias und Diskriminierende Muster]
\gls{gls-LLM} können durch Training auf historisch verzerrten Daten diskriminierende Outputs erzeugen und damit gegen das Trustworthy-\gls{gls-AI}-Prinzip der Fairness versto\ss en \parencite{ssrn5348245}.
\end{description}

In der nachfolgenden Aufstellung sind wesentliche Risiken und deren empfohlene Schutzma\ss nahmen tabellarisch zusammengefasst:

\begin{table}[H]
  \centering
  \caption[Risiken von Cloud-\gls{gls-LLM}-Diensten]{Risiken von Cloud-\gls{gls-LLM}-Diensten und empfohlene Schutzma\ss nahmen (Angelehnt an: \textcite{Data2025_AIDataQuality}, \textcite{ssrn5348245})}
  \label{tab:llm_risks}
  {\small
  \setlength{\extrarowheight}{4pt}
  \begin{tabular}{@{} >{\raggedright\arraybackslash}p{0.48\textwidth} >{\raggedright\arraybackslash}p{0.48\textwidth} @{} }
    \toprule
    \textbf{Risiko} & \textbf{Ma\ss nahmen} \\
    \midrule
    Datenübermittlung in Drittstaaten & \gls{gls-SCC}, Datenminimierung, Verschlüsselung vor Übertragung, EU-Datenorte vertraglich festlegen \\
    \addlinespace[2pt]
    Begrenzte Einsicht/Auditierbarkeit & On-premise/Edge-Optionen für kritische Workloads, Lösch-/Audit-Rechte im \gls{gls-DPA}, transparente Subprozessoren-Listen \\
    \addlinespace[2pt]
    Unkontrollierte Datenspeicherung & Explizite Aufbewahrungsfristen im \gls{gls-DPA}, Log-Redaction (Schwärzung sensibler Daten), Metadaten-Minimierung, regelmä\ss ige Kontrollaudits \\
    \addlinespace[2pt]
    Bias und Diskriminierung & Datenqualitäts-Prüfungen, multidisziplinäre Audits (IT/Recht/Ethik), Überwachung auf diskriminierende Outputs, \gls{gls-DPIA}-Dokumentation \\
  \end{tabular}
  }
  \\[1em]
  \footnotesize{\textit{Quellen:} \textcite{ssrn5348245}, \textcite{Data2025_AIDataQuality}}
\end{table}

Neben der Vereinbarung über ein \gls{gls-DPA} sollte vor einem Produktivbetrieb wesentliche Ma\ss nahmen, beispielsweise in Form einer formalen \gls{gls-DPIA}, durchgeführt werden. Die zentralen Aspekte einer formalen \gls{gls-DPIA} sind:
\begin{itemize}
  \item \textbf{Notwendigkeit prüfen}: Ist Cloud-\gls{gls-LLM} wirklich erforderlich oder gibt es weniger invasive Alternativen?
  \item \textbf{Stakeholder-Konsultation}: Nach Möglichkeit sollten Betroffene (Anrufer, Mitarbeiter) in die Bewertung einbezogen werden. Dies orientiert sich an Artikel 35 Absatz 9 \gls{gls-GDPR}.
  \item \textbf{Fortbestehende Risiken}: Falls nach einer Mitigation hohe Risiken bestehen, ist die zuständige Datenschutzbehörde zu konsultieren. Bezugnehmend auf Artikel 36 \gls{gls-GDPR}.
\end{itemize}

Eine \gls{gls-DPIA} ist als dynamisches Dokument zu verstehen und sollte bei wesentlichen Änderungen an der Verarbeitung oder dem technischen Setup aktualisiert werden, um eine kontinuierliche Compliance zu gewährleisten \parencite{ssrn5348245}.

Die dargestellten Grundlagen zeigen, dass containerbasierte Deployment-Modelle, standardisierte \gls{gls-VoIP}-Protokolle und datenschutzkonforme Verarbeitungsprozesse die technische und rechtliche Basis für ein \gls{gls-KI}-basiertes Telefonsystem bilden.

\chapter{Anforderungen}
\label{chap:anforderungen}


Die Anforderungen dieser Arbeit orientieren sich am \gls{gls-ETSI}-\gls{gls-NFV}-Framework \parencite{ETSI-NFV} und an Standards für Zuverlässigkeit und Verfügbarkeit virtualisierter Netzwerk-Funktionen \parencite{ETSI-REL}. Im Folgenden wird zwischen der prototypischen Implementierung und den in der Literatur empfohlenen produktiven Vorgaben unterschieden.

\section{Funktionale Anforderungen}

\subsubsection{Prototypische Implementierung}

Der Prototyp basiert auf den folgenden funktionalen Anforderungen für eine containerbasierte Architektur und \gls{gls-KI}-basierte Sprachinteraktion:

\begin{itemize}
  \item \textbf{\gls{gls-SIP}-Gateway:} RFC 3261-konformer Rufaufbau und Call-Control über ein lokales Asterisk-\gls{gls-PBX} \parencite{rfc3261}
  \item \textbf{Media-Bridge:} Transcoding zwischen Audio-Codecs (G.711 für \gls{gls-PSTN}, Opus für \gls{gls-WebRTC}, \gls{gls-PCM16} für \gls{gls-ASR}) und asymmetrisches \gls{gls-RTP}-Handling für \gls{gls-SIP} $\leftrightarrow$ \gls{gls-WebRTC} \parencite{rfc3550}
  \item \textbf{\gls{gls-KI}-Pipeline:} \gls{gls-ASR} mit \gls{gls-VAD} für Spracherkennung, \gls{gls-LLM} für Intent-Erkennung und Dialogverwaltung, \gls{gls-TTS} für Antwort-Generierung
  \item \textbf{Sitzungsverwaltung:} Statusverfolgung für aktive Calls mit minimaler Persistenz
\end{itemize}

Der detaillierte Signalfluss wird in Kapitel \ref{chap:architektur} beschrieben.

\subsubsection{Produktive Vorgaben nach Standards}

Produktive Systeme erweitern die Basisfunktionalität um mehrere empfehlenswerte Anforderungen. Basierend auf dem \gls{gls-ETSI}-\gls{gls-NFV}-Rahmenwerk \parencite{ETSI-NFV, ETSI-REL} und bewährten Praktiken für containerbasierte Asterisk-Deployments empfiehlt sich eine produktive Architektur mit folgenden Vorgaben \parencite{Antonova2025}:

\begin{itemize}
  \item \textbf{Verteilte Systemkomponenten:} Trennung von \gls{gls-SIP}-Signalisierung (Kamailio), Medienverarbeitung (Asterisk Core), Datenspeicherung (PostgreSQL) und Lastverteilung als separate, unabhängig skalierbare Container-Services \parencite{Antonova2025}
  \item \textbf{Fehlertoleranz und Verfügbarkeit:} Automatische Neu-Initialisierung fehlerhafter Container, automatische Wiederherstellung von Diensten, internes Cluster-Load-Balancing \parencite{Antonova2025}
  \item \textbf{Horizontale Skalierbarkeit:} Automatische Pod-Erzeugung basierend auf \gls{gls-CPU}- und \gls{gls-RAM}-Metriken mittels Horizontal Pod Autoscaler, verteilte Lastverteilung zwischen Asterisk-Replikas \parencite{Antonova2025,KubernetesDocs}
  \item \textbf{Persistierung und Verwaltung:} Zentralisierte Datenspeicherung (\gls{gls-CDR}, Konfiguration, Authentifizierung) in relationalen Datenbanken, deklarative Konfigurationsverwaltung über Kubernetes-Manifeste \parencite{Antonova2025,KubernetesDocs}
  \item \textbf{Observability und Automatisierung:} Monitoring von System-Metriken und Logs, automatisierte Deployment- und Rollback-Prozesse, DevOps-Integration für kontinuierliche Bereitstellung \parencite{Antonova2025}
\end{itemize}

\section{Nicht-funktionale Anforderungen}

Der Prototyp basiert auf den folgenden nicht-funktionalen Anforderungen für eine containerbasierte Architektur und \gls{gls-KI}-basierte Sprachinteraktion:

\subsubsection{Prototypische Implementierung}

\begin{itemize}
  \item \textbf{Latenz:} End-to-End 2--5\,s für \gls{gls-KI}-Dialog-Interaktionen (akzeptabel für asynchrone Konversation, nicht \gls{gls-ITU-T}-konform für klassische Telefonie \parencite{ITU-T-G114}).
  \item \textbf{Verfügbarkeit:} Best-Effort ohne garantierte Uptime, kein Failover, Single-Host-Deployment.
  \item \textbf{Skalierbarkeit:} Vertikales Scaling begrenzt auf Host-Ressourcen, kein horizontales Scaling.
  \item \textbf{Sicherheit:} Container-Isolation via Docker, Basis-Authentifizierung für Zugang zu Verwaltungsschnittstellen, keine Verschlüsselung für Entwicklungs-/Testzwecke.
  \item \textbf{Datenschutz:} Minimale Session-Persistenz, keine langfristige Datenspeicherung.
\end{itemize}

\subsubsection{Produktive Vorgaben nach Standards}

Die folgenden Anforderungen beschreiben das produktive Zielmodell nach \gls{gls-ETSI}-REL-Standard \parencite{ETSI-REL} und \gls{gls-ITU-T} G.114 \parencite{ITU-T-G114}. Sie folgen bewährten Praktiken anhand containerbasierten Asterisk-Deployments \parencite{Antonova2025}:

\begin{itemize}
  \item \textbf{Latenz:} < 400\,ms für Media-Übertragung (Sprachübertragung), < 2\,s für \gls{gls-KI}-Dialog-Systeme (asynchrone Interaktion), nach \gls{gls-ITU-T} G.114 \parencite{ITU-T-G114}.
  \item \textbf{Verfügbarkeit:} Hohe Verfügbarkeit mit Redundanz kritischer Komponenten und automatischem Failover \parencite{ETSI-REL}.
  \item \textbf{Skalierbarkeit:} Horizontale Skalierung über Kubernetes mit automatischer Last-Verteilung basierend auf Standard-Metriken wie \gls{gls-CPU} und \gls{gls-RAM} oder benutzerdefinierten Metriken wie Anrufvolumen \parencite{KubernetesDocs}. Containerbasierte Telefonsysteme können solche \gls{gls-VoIP}-spezifischen Metriken über benutzerdefinierte Monitoring-Lösungen realisieren \parencite{Antonova2025}.
  \item \textbf{Sicherheit:} Verschlüsselte Datenübertragung für kritische Komponenten.\gls{gls-SRTP} für Medien anhand von RFC 3711 \parencite{rfc3711}, Schlüssel-Management für \gls{gls-API}-Schlüssel und sichere Aufbewahrung von Zugangsdaten außerhalb des Quell-Repositories \parencite{Antonova2025}.
  \item \textbf{Datenschutz:} \gls{gls-DSGVO}-konforme Praktiken einschlie\ss lich Datenminimierung, Pseudonymisierung und dokumentierter Aufbewahrungsrichtlinien \parencite{Data2025_AIDataQuality}.
\end{itemize}



\chapter{Architektur}
\label{chap:architektur}

Die vorliegende Architektur folgt dem Paradigma der \gls{gls-NFV}. Das \gls{gls-ETSI} definiert dieses Paradigma als Standard. Das Ziel davon ist es, spezialisierte Hardware durch software- und containerbasierte Dienste zu ersetzen. Diese Dienste sollen stattdessen auf Standard-Servern betrieben werden können \parencite{ETSI-NFV}. Auch softwarebasierte \gls{gls-PBX}- und Telefonsysteme lassen sich im Sinne des \gls{gls-NFV}-Modells als \gls{gls-VNF} einordnen, da sie zentrale Kommunikationsfunktionen wie Signalisierung, Vermittlung und Medienverarbeitung übernehmen und virtualisiert betrieben werden können. Damit wird die Flexibilität und Skalierbarkeit klassischer Kommunikationssysteme anhand der Intention des \gls{gls-ETSI}-Referenzmodells erhöht.

\section{\gls{gls-ETSI} \glstext{gls-NFV}-Architektur-Standard}
\label{sec:etsi_nfv_architektur}

Die \gls{gls-ETSI}-\gls{gls-NFV}-Architektur bietet ein flexibles Rahmenwerk für die Virtualisierung von Netzwerkfunktionen auf Standardhardware und bildet damit die Grundlage für skalierbare Kommunikationssysteme. Sie unterscheidet drei zentrale Schichten: Die Infrastruktur stellt virtualisierte Ressourcen wie Rechenleistung, Speicher und Netzwerk bereit. Darauf aufbauend werden einzelne Netzwerkfunktionen als Softwaremodule betrieben, die unabhängig voneinander auf dieser Infrastruktur laufen. Die Steuerung und Verwaltung aller Komponenten übernimmt die \gls{gls-MANO}-Eben, die aus spezialisierten Modulen für das Lebenszyklusmanagement der \gls{gls-VNF}, der Ressourcenverwaltung und der übergreifenden Koordination besteht \parencite{ETSI-NFV}. Diese modulare Architektur ermöglicht es, spezialisierte Dienste wie Signalisierung, Medienverarbeitung und \gls{gls-KI}-Funktionen als unabhängig skalierbare Services zu implementieren. Unterstützt wird dies durch Orchestrierung und Failover-Mechanismen \parencite{ETSI-REL}.

Der Prototyp dieser Arbeit implementiert ein minimales Subset dieser Architektur mit Single-Host-Deployment statt verteilter Infrastruktur. Alle produktiven \gls{gls-ETSI}-Standards werden als optionale Upgrade-Strukturen angesehen und sind in hochskalierten Nutzungsszenarien erforderlich.

\subsection{Architektur-Übersicht}

Das konkrete System in dieser Arbeit realisiert einen durchgängigen Kommunikationspfad zwischen klassischer Telefonie und \gls{gls-KI}-basierter Sprachverarbeitung. Ein eingehender Anruf wird vom \gls{gls-SIP}-Gateway entgegengenommen, die Medienschicht transkodiert das Audio zwischen verschiedenen Codecs wie G.711, Opus und PCM16 und die Anwendungsschicht verarbeitet Sprache durch \gls{gls-ASR}, \gls{gls-LLM} und \gls{gls-TTS}. Die generierte Sprachantwort durchläuft den gleichen Pfad in umgekehrter Richtung zurück zum Anrufer, wobei jede Schicht isoliert, skaliert und gesichert werden kann.

Ein aktuelles Beispiel für eine moderne, \gls{gls-KI}-basierte Telefoniearchitektur liefert das System von Todorov, Todor N, das für die Automatisierung von Taxi-Bestellungen per Telefon entwickelt wurde. Die Architektur besteht aus einem offenen \gls{gls-PABX} für das Call-Routing, Echtzeit-\gls{gls-STT}/\gls{gls-TTS}-Modulen, einem domänenspezifisch trainierten \gls{gls-LLM}-Kern und robusten Schnittstellen zu Legacy-Dispatch-Systemen. Diese Struktur ist darauf ausgelegt, hohe Anrufvolumina mit minimaler Latenz zu verarbeiten und bietet Fallback-Mechanismen zu menschlichen Operatoren \parencite{Todorov2025LLMDispatcher}.

\subsection{Einordnung anhand von Architektur-Schichten}

Die in dieser Arbeit verwendete Architektur gliedert sich in drei logische Schichten und wird wie folgt beschrieben:

Die \textbf{Signalisierungsschicht}, die für die \gls{gls-SIP}-Steuerung zuständig ist, die \textbf{Medienschicht}, die \gls{gls-RTP}-Audio verwaltet, und die \textbf{Anwendungsschicht} übernehmen die Steuerung der \gls{gls-KI}-Logik. Diese Schichtung ermöglicht erweiterbare und isolierte Komponenten nach dem \gls{gls-NFV}-Paradigma \parencite{ETSI-NFV}.

\textbf{Signalisierungsschicht:}\\
Das \gls{gls-SIP}-Gateway \parencite{rfc3261} empfängt INVITE-Requests über \gls{gls-UDP})sowie \gls{gls-TCP} Port 5060 und führt \gls{gls-NAT}-Traversierung mit erweiterten Funktionen wie \gls{gls-STUN}/\gls{gls-TURN}/\gls{gls-ICE} durch \parencite{rfc5389, rfc5766, rfc5245}. Ein Applikationsserver verarbeitet \gls{gls-SIP}-Signale und erstellt Room-Sessions als Mikroservice.
\clearpage
\textbf{Medienschicht:}\\
Ein \gls{gls-RTP}-Relay leitet \gls{gls-RTP}-Pakete vom \gls{gls-SIP}-Telefon zum \gls{gls-WebRTC}-Server weiter und transkodiert die notwendigen Audio Codecs (G.711 $\leftrightarrow$ Opus $\leftrightarrow$ PCM16), die in RFC 3550 standardisiert sind \parencite{rfc3550}. Der \gls{gls-WebRTC}-Server terminiert \gls{gls-WebRTC}-Streams und stellt Opus-Audio mit 48 kHz bereit. Produktive Systeme implementieren zusätzlich \gls{gls-SRTP}-Verschlüsselung nach RFC 3711 \parencite{rfc3711}.

\textbf{Anwendungsschicht:}\\
Die Anwendungsschicht integriert \gls{gls-KI}-Dienste für Sprachverarbeitung und Dialog-Management. Sie verbindet Audio-Streams mit \gls{gls-ASR} für Spracherkennung, \gls{gls-LLM} für Intent-Verarbeitung und Antwort-Generierung sowie \gls{gls-TTS} für Sprachsynthese. \gls{gls-VAD} ermöglicht dabei die Erkennung von Sprachaktivität zur Optimierung der Verarbeitungskette. Diese Schicht realisiert den interaktiven \gls{gls-KI}-Dialog zwischen Anrufer und System.\parencite{Antonova2025,rfc3551,rfc3550}

\begin{table}[H]
\centering
\caption[Zielarchitektur-Komponenten]{Zielarchitektur: Komponenten und Funktionen nach Antonova et al., \gls{gls-ETSI}-\gls{gls-NFV}-Richtlinien und Todorov \parencite{Antonova2025, ETSI-REL, Todorov2025LLMDispatcher}}
\label{tab:vnf_components}
\small
\begin{tabular}{|>{\raggedright\arraybackslash}p{3.6cm}|>{\raggedright\arraybackslash}p{5.3cm}|>{\raggedright\arraybackslash}p{4.1cm}|}
\hline
\textbf{Komponente} & \textbf{Funktion} & \textbf{Charakteristik} \\
\hline
\gls{gls-PABX} (Asterisk) & Eingehende Anrufe annehmen, Audio an \gls{gls-STT} leiten, Call-Rerouting & Hochvolumige Anrufverwaltung, geringe Latenz \\
\hline
\gls{gls-SIP}-Proxy (Kamailio) & \gls{gls-SIP}-Nachrichten-Routing, Lastverteilung vor der \gls{gls-PABX} & Entlastung des \gls{gls-PABX}-Kerns \\
\hline
\gls{gls-RTP}-Proxy & \gls{gls-RTP}-Relay, Medienpfad unterstützen & Entkopplung des Medienpfads \\
\hline
\gls{gls-STT}/\gls{gls-TTS} & Sprache $\rightarrow$ Text, Text $\rightarrow$ Sprache & Echtzeit-Verarbeitung, robuste Spracherkennung \\
\hline
\gls{gls-LLM}-Kern mit Task-Orchestrierung & Dialogsteuerung, Kontextverarbeitung, \gls{gls-API}-Aufrufe & Niedrige Inferenzlatenz, domänenspezifisches Training \\
\hline
Externe \gls{gls-API}-Schnittstellen & Geocodierung, Adressvalidierung, Dispatch-Integration & Interoperabilität mit Legacy-Systemen \\
\hline
Fallback-Mechanismen & Rerouting zu menschlichen Operatoren & Servicekontinuität bei Fehlerfällen \\
\hline
Datenbank (PostgreSQL/MySQL) & Konfiguration, \gls{gls-CDR}, Authentifizierung & Zuverlässige Datenhaltung \\
\hline
Konfigurationsmanagement (ConfigMaps/Geheimnisse) & Zentrale Verwaltung von Konfigurationen und Geheimnisse & Trennung von Code und Geheimnissen \\
\hline
Load Balancer/Ingress & Traffic-Verteilung, Cluster-Einstieg & Skalierbarer Zugriffspunkt \\
\hline
Monitoring/Logging (Prometheus/Grafana) & Metriken und Log-Analyse & Betriebsüberwachung und Fehleranalyse \\
\hline
Skalierbare Architektur \gls{gls-ETSI}-REL & Komponenten replizieren, Zustand handhaben, Ausfälle abfedern & Zuverlässigkeit, Skalierung, Failover \\
\hline
\end{tabular}
\\[1em]
\footnotesize{\textit{Quellen:} \textcite{Antonova2025} (Asterisk-\gls{gls-PABX}, Kamailio, \gls{gls-RTP}-Proxy, Datenbank, ConfigMaps/Geheimnisse, Load Balancer, Monitoring/Logging), \textcite{ETSI-REL} (Zuverlässigkeit/Skalierung als Architekturprinzipien), \textcite{Todorov2025LLMDispatcher} \gls{gls-PABX}, \gls{gls-STT}/\gls{gls-TTS}, \gls{gls-LLM}-Kern, externe \gls{gls-API}, Fallback))}
\end{table}
\FloatBarrier

\noindent \textit{Anmerkung:} Tabelle~\ref{tab:vnf_components} fasst die Komponenten und Architekturprinzipien aus Antonova, \gls{gls-ETSI}-REL und Todorov zusammen. Antonova et al. beschreiben eine Kubernetes-basierte Asterisk-Architektur mit Kamailio, PostgreSQL/MySQL, Load Balancer und Monitoring-Tools. Todorov ergänzt eine produktive Dispatcher-Architektur mit \gls{gls-STT}/\gls{gls-TTS}, \gls{gls-LLM}-Kern, externen \gls{gls-API}-Schnittstellen und Fallback auf menschliche Operatoren. \gls{gls-ETSI}-REL liefert die Prinzipien für Zuverlässigkeit, Skalierung und Failover.

\subsubsection{Abgrenzung: Zielarchitektur und Prototyp-Implementierung}

Die beschriebene Zielarchitektur stellt ein Produktiv-Referenzmodell dar, das auf etablierten Standards und praktischen Validierungen basiert. Die zentralen Komponenten und Prinzipien sind in Tabelle~\ref{tab:vnf_components} zusammengefasst. Der Prototyp dieser Arbeit implementiert ein reduziertes Single-Host-Subset dieser Zielarchitektur.

Kapitel~\ref{chap:implementierung} beschreibt die konkrete Implementierung dieses reduzierten Prototyps und deren Abweichungen von der Zielarchitektur.

\section{Sicherheitskonzept}

Der Prototyp orientiert sich an \gls{gls-ETSI}-Sicherheitsstandards \parencite{ETSI-REL}, implementiert jedoch ein reduziertes Sicherheitsmodell für Entwicklungs- und Testzwecke. Daher unterscheidet sich die Darstellung speziell zwischen Prototyp-Anforderungen und produktiven Sicherheitsma\ss nahmen nach \gls{gls-RFC}- und \gls{gls-ETSI}-Standards.

\subsection{Sicherheitsarchitektur}

\textbf{Prototyp-Modell:}\\
Die folgenden Sicherheitsprinzipien orientieren sich an etablierten Standards (wie \gls{gls-ETSI}-REL und relevanten \gls{gls-RFC}, wurden jedoch spezifisch für den Prototyp angepasst und reduziert:

\begin{itemize}
  \item \textbf{Netzwerk-Isolation:} containerbasierte Segmentierung zur Trennung von Komponenten und Minimierung der Angriffsfläche
  \item \textbf{Zugriffskontrolle:} tokenbasierte Authentifizierung mit definierten Berechtigungen und zeitlich begrenzter Gültigkeit
  \item \textbf{Service-Authentifizierung:} Authentifizierung zwischen internen Services über kryptographische Verfahren
  \item \textbf{Input-Validierung:} Prüfung und Validierung aller Eingabeparameter zur Vermeidung von Injection-Angriffen
  \item \textbf{Nachvollziehbarkeit:} Logging-Mechanismen zur Protokollierung sicherheitsrelevanter Ereignisse
  \item \textbf{Konfigurationsmanagement:} Externe Verwaltung sensibler Daten ohne Einbettung in Quellcode
  \item \textbf{Transportverschlüsselung:} Verzicht auf \gls{gls-TLS}/\gls{gls-SRTP} im Prototyp zugunsten reduzierter Komplexität
\end{itemize}
\textbf{Produktiv-Anforderungen nach Standards:}\\
Produktive Systeme müssen nach \gls{gls-ETSI}-REL \parencite{ETSI-REL} und RFC-Standards folgende Ma\ss nahmen implementieren:

\begin{itemize}
  \item \textbf{\gls{gls-SIP}/\gls{gls-TLS}:} Port 5061 mit \gls{gls-TLS} 1.3 (RFC 8446 \parencite{rfc8446}, RFC 5630 \parencite{rfc5630})
  \item \textbf{\gls{gls-SRTP}:} Verschlüsselte \gls{gls-RTP}-Streams (RFC 3711 \parencite{rfc3711})
  \item \textbf{\gls{gls-DTLS}-\gls{gls-SRTP}:} Für \gls{gls-WebRTC}-Clients (RFC 5764 \parencite{rfc5764})
  \item \textbf{\gls{gls-mTLS}:} zertifikatbasierte Authentifizierung zwischen Services \parencite{rfc8446}
  \item \textbf{\gls{gls-RBAC}:} Rollenbasierte Zugriffskontrolle für administrative Funktionen \parencite{KubernetesDocs}
\end{itemize}

\subsection{Netzwerk-Segmentierung}

Die Architektur folgt dem Prinzip der funktionalen Trennung nach \gls{gls-ETSI}-REL \parencite{ETSI-REL}:

\begin{table}[H]
\centering
\caption[Funktionale Netzwerk-Segmentierung]{Funktionale Netzwerk-Segmentierung nach \gls{gls-ETSI}-Standards}
\label{tab:network_segmentation}
\small
\begin{tabular}{|>{\raggedright\arraybackslash}p{2.5cm}|>{\raggedright\arraybackslash}p{4.5cm}|>{\raggedright\arraybackslash}p{5cm}|}
\hline
\textbf{Zone} & \textbf{Komponenten} & \textbf{Schutzma\ss nahmen} \\
\hline
Signalisierung & \gls{gls-SIP}-Gateway, \gls{gls-NAT}-Traversierung & \gls{gls-TLS} 1.3, \gls{gls-SIP} (RFC 5630) \\
\hline
Medien & \gls{gls-RTP}-Relay, \gls{gls-WebRTC} & \gls{gls-SRTP}, \gls{gls-DTLS} (RFC 3711, 5764) \\
\hline
Anwendung & \gls{gls-KI}-Dienste, \gls{gls-LLM}, \gls{gls-ASR}/\gls{gls-TTS} & \gls{gls-TLS} 1.3 (RFC 8446), \gls{gls-API}-Verschlüsselung \\
\hline
Datenhaltung & Konfiguration, State-Management & Netzwerk-Isolation, Zugriffskontrolle angelehnt an \gls{gls-ITU-T} \\
\hline
\end{tabular}
\\[1em]
\footnotesize{\textit{Quelle:} \textcite{ETSI-REL}}
\end{table}
\FloatBarrier

Die folgende Zielarchitektur wurde auf der Grundlage etablierter Architekturprinzipien und unter Einbeziehung relevanter Literatur abgeleitet \parencite{ETSI-NFV,Antonova2025,Todorov2025LLMDispatcher}.

\begin{figure}[H]
\centering
\includegraphics[width=0.85\textwidth,height=0.85\textheight,keepaspectratio]{images/architektur.png}
\caption{Architektur-Übersicht: Eigene Darstellung mit PlantUML
(\url{https://plantuml.com/}) in Anlehnung an relevante Literatur.}
\label{fig:architektur}
\end{figure}

Die konkrete Implementierung dieser Architektur-Prinzipien wird in Kapitel~\ref{chap:implementierung} beschrieben.

\subsection{Bereitstellungsmodelle und Deployment}

Bei der Systemarchitektur stehen drei grundlegende Modelle für eine Bereitstellung zur Verfügung: Das \textbf{monolithische Modell} fasst alle Funktionen in einer einzigen Container-Instanz zusammen, was schnelle Prototypentwicklung ermöglicht, aber Skalierungsprobleme mit sich bringt. Das \textbf{vollständig dezentralisierte Modell} zerlegt die Anwendung in viele einzelne, unabhängige Mikroservices, bietet höhere Flexibilität und Ausfallsicherheit, erfordert aber umfangreiche Orchestrierungs- und Verwaltungsma\ss nahmen. Der \textbf{hybride Ansatz} stellt einen Mittelweg dar: Er isoliert nur die kritischsten Funktionen als eigenständige Services, während unkritische Funktionen weiterhin zusammengefasst bleiben. Der vorliegende Prototyp folgt dieser hybriden Strategie, um eine schrittweise Entkopplung von Komponenten ohne umfassende Neuarchitekturierung zu ermöglichen \parencite{Antonova2025}.

Im folgenden Kapitel~\ref{chap:marktanalyse} wird nun genauer auf einige Anbieter im Bereich der \gls{gls-KI}-basierten Telefonie eingegangen. Es werden exemplarisch sieben Anbieter betrachtet, die unterschiedlichste Architekturen nutzen, diverse Zielgruppen adressieren und verschiedene Kernfeatures verwenden. Au\ss erdem werden ausführlich ihre Preismodelle dargestellt und mögliche anfallende Kosten analysiert. Zusätzlich werden wichtige datenschutzrechtliche Aspekte diskutiert und abgeleitete Ma\ss nahmen empfohlen.

\chapter{Marktanalyse}\label{chap:marktanalyse}

% Helfer-Makros und Spaltentyp (werden hier definiert, damit sie vor Nutzung verfügbar sind)
\newcommand{\vendor}[3]{\noindent\raisebox{-0.1\height}{\includegraphics[height=7mm]{#1}}\hspace{0.6em}\textbf{#2}. #3\par\vspace{3pt}}
\newcommand{\vendorText}[2]{\noindent\textbf{#1}. #2\par\vspace{3pt}}
\newcommand{\providerlogo}[1]{\raisebox{0pt}[0pt][0pt]{\includegraphics[height=6mm]{#1}}}
\newcolumntype{Y}{>{\centering\arraybackslash}X}
\newcolumntype{L}{>{\raggedright\arraybackslash}X}
\newcolumntype{C}[1]{>{\centering\arraybackslash}m{#1}}

Diese Kurzanalyse richtet sich an \gls{gls-KMU} in \gls{gls-DACH}. Vor allem werden Unternehmen aus diesen Regionen adressiert, die \gls{gls-KI}-basierte Telefonsysteme mit \gls{gls-SIP}-Integration einführen möchten und dabei volle Datenschutzkontrolle bewahren müssen.

Zur Realisierung stehen drei grundsätzliche Bereitstellungsmodelle zur Verfügung:(1) Cloudbasierte Lösungen für schnellen Markteintritt, (2) Hybrid-Ansätze mit lokal- und teilweise in Cloud-Systeme ausgelagerten Dienste, (3) Vollständig selbstgehostete Systeme für maximale Kontrolle und \gls{gls-DSGVO}-Konformität \parencite{Antonova2025}. Solche Modelle und deren Unterschiede werden auch international diskutiert, etwa in systematischen Übersichten und aktuell produktiven Umgebungen \parencite{Todorov2025LLMDispatcher,Ngobeni2025AIVoIP}. Die Wahl zwischen diesen Modellen impliziert Trade-offs zwischen Flexibilität, Compliance und Betriebskomplexität.

Um Anbieter von \gls{gls-KI}-basierten Telefonsystemen und deren Funktionsumfänge überhaupt vergleichen zu können, muss zunächst klar sein, welche Anforderungen für einen produktiven Betrieb unverzichtbar sind. In der klassischen IP-Telefonie sind das standardisierte \gls{gls-PBX}-Funktionen. Einige Beispiele dafür sind: \gls{gls-SIP}-Registrierung, bidirektionales In-/Outbound-Routing, \gls{gls-DTMF} oder Verarbeitung für \gls{gls-IVR}-Menüs \parencite{rfc3261}. Im Einsatz von \gls{gls-KI}-basierten Telefonsystemen sind andere Eigenschaften notwendig: \gls{gls-WebRTC}/\gls{gls-RTP}-Kompatibilität, Barge-in-Fähigkeiten und transparente Transkodierung (G.711 \Leftrightarrow\ Opus) \parencite{rfc6716}. Die Integration von \gls{gls-KI} birgt einige Herausforderungen, die entscheiden, ob ein System für viele tausend Anrufe täglich stabil funktioniert, eine gute Nutzererfahrung bietet und somit auch die Produktivität in Unternehmen und deren Prozessen steigert. In Systemen, die klassische Protokolle, Technologien und Funktionen mit modernen Standards verknüpfen, müssen die genannten Aspekte im Zusammenspiel implementiert werden. Internationale Studien zeigen, dass diese Herausforderungen und Lösungsansätze in verschiedenen Märkten ähnlich sind \parencite{Todorov2025LLMDispatcher,Ngobeni2025AIVoIP}.
\clearpage
\section{Integrationsaufwand }

Der Integrationsaufwand für \gls{gls-KI}-basierte Telefonsysteme variiert je nach Anbieter und Bereitstellungsmodell deutlich. Gro\ss e Plattformen, welche \gls{gls-CPaaS} anbieten, stellen umfangreiche Dokumentationen, \gls{gls-API} und oft zentral verwaltete Dienste bereit, wodurch die Integration meist schnell und mit wenig Eigenaufwand möglich ist. \gls{gls-DACH}-spezialisierte Anbieter und \gls{gls-OSS} bieten mehr Kontrolle, erfordern jedoch häufig zusätzliche Konfiguration und technisches Know-how, insbesondere bei vollständig selbst gehosteten Systemen. Hybrid-Modelle kombinieren lokale Kernsysteme mit Cloud-Diensten und ermöglichen so Flexibilität, setzen aber ebenfalls Erfahrung in der Systemintegration voraus. Internationale Erfahrungen und systematische Übersichten belegen, dass diese Integrationsherausforderungen und Vorteile auch in anderen Märkten beobachtet werden können \parencite{Todorov2025LLMDispatcher,Ngobeni2025AIVoIP}. Insgesamt bestimmt die Wahl des Anbieters und des Bereitstellungsmodells, wie viel Eigenleistung bei der Integration in eigene Infrastruktur und beim Betrieb solcher Telefonsysteme notwendig ist.

\section{Datenschutz und Compliance}\label{sec:datenschutz_compliance}

Für \gls{gls-KMU} in \gls{gls-DACH} ist Datenschutz nicht optional, sondern zentral. Das bedeutet konkret: \gls{gls-DSGVO}-Konformität und gültige \gls{gls-AVV} mit dem Anbieter bilden die Basis für konforme Systeme. EU-Datenstandorte und dokumentierte Speicherorte sind nicht nur legal erforderlich, sondern beeinflussen auch Latenzen und Zuverlässigkeit. Besonders kritisch sind Audio-Aufzeichnungen. Bei Multi-Vendor-Setups müssen Subprozessoren für \gls{gls-LLM}-Dienste transparent offengelegt und in \gls{gls-AVV} dokumentiert sein. Dies ist ein Aspekt einer Vielzahl weiterer Aspekte, die in der Evaluationsphase explizit geprüft werden sollten. Datenminimierung und Datenqualität
bedeutet auch: Welche Metadaten speichert der Anbieter wirklich? Telefonnummern, Transkripte oder Intent-Labels müssen minimiert und pseudonymisiert werden \parencite{Data2025_AIDataQuality}.

\section{Limitierungen aus der Praxis}
In der Praxis zeigen sich wiederkehrende Herausforderungen, die bei der Auswahl von Anbietern, die \gls{gls-KI}-basierte Systeme bereitstellen, berücksichtigt werden müssen. Die Rufummernverfügbarkeit ist eine kritische Herausforderung und für \gls{gls-SIP}-Nutzung unverzichtbar. Viele Cloud-Plattformen bieten primär Rufnummern an, die an Standorte in den USA gebunden sind, weil die Infrastruktur historisch dort konzentriert ist. LiveKit-Cloud beispielsweise hat optionale \gls{gls-SIP}-Anbindung und Rufnummernkauf in einem Video-Announcement vorgestellt und auch in der offiziellen LiveKit Dokumentation beschrieben \parencite{livekit_telephony_docs}. Jedoch wird in diesem Zusammenhang ein rein cloudbasiertes Bereitstellungsmodell angeboten. Internationale Verfügbarkeiten von Rufnummern, die in das cloudbasierte System direkt integriert werden können, werden nur schrittweise erweitert. Für Unternehmen und Kunden in \gls{gls-DACH} bedeutet das, dass Portierungen oder Rufnummern separat geklärt werden müssen. Dieser Prozess kann einige Zeit und viele Ressourcen benötigen. Hinzu kommt die geografische Abhängigkeit von Rechenregionen und Media-Path-Routing \parencite{rfc3550}: Je weiter der Datenfluss entfernt ist, desto höher ist die Latenz \parencite{ITU-T-G114} und desto schwieriger wird die Einhaltung von \gls{gls-DSGVO}-Anforderungen zur lokalen Datenverarbeitung.

\section{Marktcharakteristika und Trends}
Der Markt für \gls{gls-KI}-basierte Telefonsysteme mit \gls{gls-SIP}-Integration ist seit 2022-2023 stark fragmentiert und wächst schnell. Beispielsweise profitieren Anbieter wie RetellAI \parencite{retellai_sip_blog} von der Verfügbarkeit hochperformanter \gls{gls-LLM}. Anbieter nutzen diese leistungsstarken Technologien in ihrer Softwareentwicklung und machen sie damit für viele weitere Kunden durch ihre Software und Plattformen zugänglich.
Basierend auf Marktbeobachtungen und öffentlichen Ankündigungen sind diese typischerweise mit Risikokapital finanziert und fokussieren sich auf Skalierungen auf Enterprise-Level, anstatt kurzfristige Profitabilität zu steigern und die Stärkung lokaler IT-Partner oder \gls{gls-KMU} zu ermöglichen. \parencite{retellai_seed_blog} Das birgt Risiken für langfristige Partnerschaften.

Open-Source-Infrastruktur profitiert von dieser Markt-Fragmentierung und deren Investmentstrategien. {Asterisk} ist als produktionsreife PBX/VoIP-Plattform etabliert \parencite{asterisk_project}. Moderne Frameworks wie LiveKit ermöglichen \gls{gls-WebRTC}-Infrastrukturen und können sowohl cloudbasiert als auch lokal eingesetzt werden \parencite{livekit_docs}.

\section{Anbieter-Überblick und Kurzbewertung}
Mithilfe der Anforderungen werden nun konkrete Lösungen betrachtet. Die folgende Tabelle~\ref{tab:anbieter_overview} stellt sieben kommerzielle Anbieter gegenüber. Dabei werden gro\ss e, etablierte \gls{gls-CPaaS}-Plattformen und spezialisierte Anbieter in der \gls{gls-DACH}-Region betrachtet. Das Ziel ist es nicht, eine eins-zu-eins-Vergleichbarkeit zu suggerieren, sondern verschiedene Anbieter und deren Bereitstellungsmodelle zu vergleichen und einzuordnen. Was sind die Kernfunktionen? Für welche Zielgruppe ist der Anbieter konzipiert? Wie steht es um Datenschutz und Compliance? Dies sind wichtige Kriterien, die bei Kaufentscheidungen oft zu spät geprüft werden. Die Bewertungen basieren auf öffentlich zugänglichen Produktdokumentationen und
Webseiten der Anbieter (Stand: Januar 2026) und sind bewusst nicht normativ, sondern deskriptiv formuliert.

\renewcommand{\arraystretch}{1.18} % angenehmer Zeilenabstand
\begin{table}[H]
\caption{Überblick zu Kernfeatures, Zielgruppe und DSGVO/Compliance}\label{tab:anbieter_overview}
\small
\begin{tabularx}{\linewidth}{@{}>{\raggedright\arraybackslash}m{3.2cm} L L >{\raggedright\arraybackslash}m{3.2cm}@{}}
\toprule
\textbf{Anbieter} & \textbf{Kernfeatures} & \textbf{Zielgruppe} & \textbf{DSGVO} \\
\midrule
\parbox[t]{3.2cm}{Cloudonix\\\url{cloudonix.com}} & \gls{gls-SIP} Trunking, No-Code \gls{gls-IVR}, \gls{gls-API}s & Entwickler, Unternehmen & \parbox[t]{3.2cm}{US-Fokus, EU-Hosting überprüfen, mit \gls{gls-AVV} absichern} \\
\addlinespace[8pt]
\parbox[t]{3.2cm}{DialogShift\\\url{dialogshift.com}} & Phone \gls{gls-AI}, Chat \gls{gls-AI}, WhatsApp-Integration & Hotellerie, Service-Teams & \parbox[t]{3.2cm}{\gls{gls-DSGVO}-orientiert, mit \gls{gls-AVV} absichern} \\
\addlinespace[8pt]
\parbox[t]{3.2cm}{RetellAI\\\url{retellai.com}} & \gls{gls-LLM} Voice Agent, Multi-Turn, \gls{gls-API} & Kundenservice, Gesundheitssektor & \parbox[t]{3.2cm}{US-Fokus, EU-Hosting überprüfen, mit \gls{gls-AVV} absichern} \\
\addlinespace[8pt]
\parbox[t]{3.2cm}{VideoSDK\\\url{videosdk.live}} & Voice/Video-\gls{gls-API}, \gls{gls-SIP} & Entwickler, Unternehmen & \parbox[t]{3.2cm}{US-Fokus, EU-Hosting überprüfen, mit \gls{gls-AVV} absichern} \\
\addlinespace[8pt]
\parbox[t]{3.2cm}{Botfriends\\\url{botfriends.de}} & Voicebot/Chatbot, Omnichannel, No-Code & Kundenservice, Mittelstand & \parbox[t]{3.2cm}{EU-Hosting verfügbar, mit \gls{gls-AVV} absichern} \\
\addlinespace[8pt]
\parbox[t]{3.2cm}{Synthflow\\\url{synthflow.ai}} & Voice \gls{gls-AI}, Flow Designer, \gls{gls-API} & \gls{gls-BPO}, Gesundheitssektor & \parbox[t]{3.2cm}{EU-Regionen möglich, Kein Self-Hosting, mit \gls{gls-AVV} absichern} \\
\addlinespace[8pt]
\parbox[t]{3.2cm}{FonioAI\\\url{fonio.ai}} & Voice \gls{gls-AI} & KMU & \parbox[t]{3.2cm}{EU-Anbieter, \gls{gls-DSGVO}-konform, mit \gls{gls-AVV} absichern} \\
\bottomrule
\end{tabularx}
\begin{flushleft}\footnotesize
\textit{Quelle:} Eigene Darstellung nach \parencite{cloudonix},
\parencite{dialogshift_phone_ai}, \parencite{retellai_sip_blog}, \parencite{videosdk_phone_agent}, \parencite{botfriends_phonebot}, \parencite{synthflow}, \parencite{fonio}.
\end{flushleft}
\end{table}

\section{Kostenmodelle und Transparenz}
Neben Funktionalität und Compliance ist Transparenz bei der Preisgestaltung ein starkes Kriterium, das zur Entscheidungsfindung beiträgt. Gerade für \gls{gls-KMU}, die Skalierungskosten verlässlich planen müssen, sollten diese Aspekte bei der Einführung eines \gls{gls-KI}-basierten Telefonsystems beachtet werden. Die folgende Tabelle~\ref{tab:pricing_comparison} zeigt aktuell öffentlich verfügbare Preismodelle (Stand: Januar 2026) für realistische Szenarien.

\begin{table}[H]
\caption{Preisvergleich anhand von Bereitstellungsmodellen}\label{tab:pricing_comparison}
\small
\begin{tabularx}{\linewidth}{@{}l L L L L@{}}
\toprule
\textbf{Anbieter} & \textbf{Preismodell} & \textbf{100 Anrufe/Monat (5h)} & \textbf{1000 Anrufe/Monat (50h)} & \textbf{Transparenz} \\
\midrule
Cloudonix & Abo-Modell (0,05–0,15 EUR/min) + Trunk & 25-45 EUR & 150-460 EUR & Öffentlich \\
DialogShift & Projekt-Basis (Custom) & auf Anfrage & auf Anfrage & Eingeschränkt \\
RetellAI & Pay-as-you-go (0,10 USD/min) + APIs & 39-71 EUR\textsuperscript{1} & 393-708 EUR\textsuperscript{1} & Teilweise \\
VideoSDK & Subscription-Tier (99 USD/Monat = 1000 min) & 90 EUR\textsuperscript{2} & 90 EUR\textsuperscript{2} & Öffentlich \\
Botfriends & Full-Service (Custom) & auf Anfrage & auf Anfrage & Eingeschränkt \\
Synthflow & Fixed Subscription-Tiers (ab 99 USD/Monat) & 90 EUR\textsuperscript{3} & 90-270 EUR\textsuperscript{3} & Öffentlich \\
FonioAI & Pay-as-you-go \gls{gls-SIP}-Trunk (9,99 EUR) + Nummern (3,99 EUR) + 0,02 EUR/min & 32-64 EUR\textsuperscript{4} & 194-514 EUR\textsuperscript{4} & Öffentlich \\
\bottomrule
\end{tabularx}
\begin{flushleft}\footnotesize
\textit{Quelle:} Öffentliche Pricing-Seiten der Anbieter (abgerufen 15. Januar 2026): \parencite{cloudonix_pricing},
\parencite{retellai_pricing}, \parencite{videosdk_pricing}, \parencite{synthflow_pricing}, \parencite{botfriends_pricing}, \parencite{fonio_pricing}. \\
\textsuperscript{1} RetellAI: 0,10 USD/min (\textasciitilde 0,09 EUR/min) $\approx$ 27 EUR (100 Anrufe) bzw. 273 EUR (1000 Anrufe) + nutzungsabhängige STT/TTS/LLM-Kosten (ca. +12-44 EUR bzw. +120-435 EUR je nach API-Wahl) \\
\textsuperscript{2} VideoSDK Starter-Tier: 99 USD/Monat (\textasciitilde 90 EUR). Minutenkontingent abhängig vom gewählten Tier (1000-3000 min) \\
\textsuperscript{3} Synthflow: Starter-Tier 99 USD/Monat; Professional 299 USD/Monat; Anrufe im Tier enthalten \\
\textsuperscript{4} FonioAI: Infrastruktur (\gls{gls-SIP}-Trunk + Nummern + 0,02 EUR/min) ca. 20 EUR (100 Anrufe) bzw. 74 EUR (1000 Anrufe). \gls{gls-AI}-Layer (TTS/STT/LLM) zusätzlich ca. +12-44 EUR (100 Anrufe) bzw. +120-440 EUR (1000 Anrufe) über externe \gls{gls-API}s je nach \gls{gls-TTS}-Anbieter (günstige Alternative vs. OpenAI Standard)
\end{flushleft}
\end{table}

\section{Kostentrends und Schlussfolgerungen}
Auf Basis der aktuellen öffentlichen Preismodelle (Stand: Januar 2026, siehe Tabelle~\ref{tab:pricing_comparison} für detaillierte Kosten) zeigen sich klare Muster. Tiered-Subscription-Modelle wie zum Beispiel VideoSDK oder Synthflow können für \gls{gls-KMU} kostengünstiger sein als Pay-as-you-go-Systeme, da sie unbegrenzte Anrufe innerhalb des Minuten-Limits ermöglichen. Dies ist ein Vorteil, der jedoch erst bei höherem Anrufvolumen beziehungsweise Anrufaufkommen sichtbar wird. Pay-as-you-go-Anbieter wie RetellAI und FonioAI bieten Flexibilität, können aber bei Volumen-Schwankungen teuer werden. Ein kritischer Faktor ist die Transparenz, wobei diese nicht bei allen betrachteten Anbietern gegeben ist. VideoSDK und Synthflow haben öffentlich zugängliche Preise, die eine klare Kostenplanung ermöglichen, während projektbezogene und dadurch individuelle Kostenmodelle, die DialogShift und Botfriends verwenden, erst nach Kontaktaufnahme in den Kostenvergleich einbezogen werden können. Dieser Nachteil muss in Entscheidungsprozessen vor der Integration oder dem Kauf eines solchen Produkts bedacht werden.

Erkenntnisse zur Zuverlässigkeit virtualisierter Netzwerke treiben Hybrid-Deployments voran \parencite{Antonova2025} \parencite{ReliabilityVirtualizedNetworks}. Daher wird Unternehmen im Sektor von \gls{gls-KMU} empfohlen, Hybrid-Modelle oder vollständig selbstgehostete Systeme zu betreiben und zu nutzen. Kunden von reinen Cloud-Modellen sind oft stark an einen bestimmten Anbieter sowie dessen Dienste, Prozesse und Richtlinien gebunden und können nur mit hohem Aufwand oder hohen Kosten zu einem anderen Anbieter wechseln. FonioAI bietet sich daher als flexibler Pay-as-you-go-Anbieter der \gls{gls-DACH}-Region an und erfüllt bereits einige wichtige Grundsätze zur Datenkonformität und zum Datenschutz. Dennoch ist es wichtig, für einen produktiven Einsatz von \gls{gls-KI}-basierten Telefonsystemen individuell nach den Anforderungen des Unternehmens und den Zielen zu entscheiden.


\chapter{Implementierung}
\label{chap:implementierung}

\noindent In diesem Kapitel wird die technische Implementierung des prototypischen Systems detailliert dargelegt. Dabei werden die gewählten Technologien, die Konfigurationsdetails sowie die iterativen Entwicklungsentscheidungen beschrieben.

% Transparenz-Hinweis zur Nutzung von KI-Werkzeugen
\noindent\textit{Hinweis (Transparenz): Zur Unterstützung wurden GitHub Copilot für Implementierungszwecke und Fehlerbehebung eingesetzt. Die Nutzung erfolgte über die für Studierende der HTWK kostenlos verfügbare Lizenz.}

\section{Entwicklungs- und Testumgebung}

Die Implementierung erfolgte auf einer dedizierten Entwicklungsumgebung mit folgenden Spezifikationen:

\begin{table}[htbp]
\centering
\caption{Hardware- und Software-Spezifikationen der Entwicklungsumgebung}
\label{tab:dev-environment}
\begin{tabular}{ll}
\toprule
\textbf{Kategorie} & \textbf{Spezifikation} \\
\midrule
\multicolumn{2}{l}{\textit{Hardware}} \\
Prozessor & Intel Core i9-9900K @ 3.6 GHz (16 logische CPUs) \\
Arbeitsspeicher & 16 GB DDR4 RAM \\
\midrule
\multicolumn{2}{l}{\textit{Software}} \\
Betriebssystem & Windows 11 Pro 64-bit \\
Entwicklungsumgebung & JetBrains WebStorm 2025.3.2 \\
Container-Runtime & Docker Desktop 4.26.1 \\
Node.js & v20.11.0 LTS \\
Python & v3.12.1 \\
\midrule
\bottomrule
\end{tabular}
\end{table}
\section{Technologien und Frameworks}

Das hier entwickelte prototypische System nutzt eine spezialisierte Kombination eines modernen Frameworks namens LiveKit für \gls{gls-KI}-basierte Sprachverarbeitung, die in fünf eng integrierten Mikroservices, die mit Docker organisiert sind, implementiert ist. Diese Architektur ermöglicht Unabhängigkeit bei der Skalierung und Wartung, da jeder Service eine spezifische Aufgabe erfüllt und über standardisierte Protokolle miteinander kommuniziert. Für detaillierte Informationen zu Architekturen siehe Kapitel \ref{chap:architektur}.

\subsection{Service-Architektur und Orchestrierung}
Die Infrastruktur wird über Docker-Compose in einem isolierten Bridge-Netz orchestriert. Das ist optimal für Entwicklung und Single-Host-Deployments, weil mehrere Container mit minimaler Konfiguration verbunden werden. Fünf Dienste arbeiten zusammen, um das System zu komplettieren. Der \textbf{{LiveKit}- Server} als zentraler \gls{gls-WebRTC}-Media-Server, der Raumerstellung verwaltet und Audioströme verarbeitet. Die \textbf{{LiveKit} \gls{gls-SIP}-Bridge} dient als Gateway zwischen klassischem \gls{gls-SIP} und \gls{gls-WebRTC}. \textbf{Asterisk} fungiert als ausgereifte Open-Source-\gls{gls-PBX} und übernimmt daher klassische Funktionen der \gls{gls-SIP}-Telefonie. \textbf{Redis} ist die Instanz, die für die Session-State-Verwaltung und Anrufmetadaten zuständig ist, damit \gls{gls-SIP}-Bridge und Asterisk konsistent bleiben. Der \textbf{Python Agent Worker} nutzt das LiveKit Agents \gls{gls-SDK}, um Room-Events zu verarbeiten, eine \gls{gls-VAD} vorzuschalten und die \gls{gls-KI}-Pipeline samt Metrik-Logging zu orchestrieren. Die Architektur verteilt Aufgaben auf die einzelnen spezialisierten Dienste, anstatt sie zu zentralisieren, und erhöht so die Fehlertoleranz und Wiederverwendbarkeit.

\subsection{Paketmanagement und Paketabhängigkeiten}

Das System nutzt zwei separate Dependency-Management-Systeme für Frontend und Backend.

\textbf{Frontend-Abhängigkeiten (npm):}\\
Die \texttt{package.json} definiert alle JavaScript/TypeScript-Abhängigkeiten für das Next.js-Frontend. Die Kern-Abhängigkeiten sind: \texttt{@livekit/components-react} (v2.0.0) für vorgefertigte React-Komponenten zur \gls{gls-WebRTC}-Integration, \texttt{livekit-client} (v2.0.0) als Low-Level-\gls{gls-WebRTC}-Client-Bibliothek, \texttt{livekit-server-sdk} (v2.0.0) für serverseitige Token-Signierung, \texttt{next} (v14.0.0) als React-Framework mit Server-Side-Rendering, \texttt{react} und \texttt{react-dom} (v18.2.0) als \gls{gls-UI}-Bibliotheken. Die Installation erfolgt via \texttt{npm install}, wodurch ein \texttt{node\_modules}-Verzeichnis mit allen transitiven Abhängigkeiten erstellt wird.
\clearpage
\textbf{Backend-Abhängigkeiten (Python):}\\
Die \texttt{requirements.txt} spezifiziert Python-Pakete für den Agent Worker. Kritische Abhängigkeiten sind: \texttt{livekit} (>=0.11.0) als Python-Client für LiveKit, \texttt{livekit-agents} (>=0.8.0) als High-Level-Framework für Sprachagenten, \texttt{livekit-plugins-openai} (>=0.6.0) für OpenAI-\gls{gls-API}-Integration, \texttt{livekit-plugins-silero} (>=0.6.0) für lokale \gls{gls-VAD}, \texttt{python-dotenv} (>=1.0.0) zum Laden von Umgebungsvariablen sowie \texttt{requests} (>=2.31.0) und \texttt{PyJWT} (>=2.8.0) für \gls{gls-HTTP}-Requests und \gls{gls-JWT}-Handling. Die Installation erfolgt via \texttt{pip install -r requirements.txt}, was alle Pakete inklusive ihrer transitiven Abhängigkeiten aus dem Python Package Index bezieht.

\subsection{Frontend- und Backend-Implementierung}

Das Web-Frontend wurde mit Next.js 14.0.0 entwickelt und bildet eine React-basierte \gls{gls-SPA}. Die Architektur setzt auf zwei Haupt-Dependency-Gruppen: Erstens \texttt{@livekit/components-react} (v2.0.0), eine von LiveKit gepflegte React-Komponenten-Bibliothek, die WebRTC-Integration abstrahiert und reusable \gls{gls-UI}-Komponenten für Räume, Aufzeichnungen und Teilnehmer anbietet. Zweitens \texttt{livekit-client} (v2.0.0), die zugrunde liegende JavaScript \gls{gls-WebRTC}-Bibliothek, die die Low-Level-Connectivity-Logik und
Streaming-Protokolle handhabt.

Die Komponenten-Architektur im Frontend folgt dem \textbf{Observer-Pattern}: Komponenten nutzen Hooks wie \texttt{useConnectionState()} und \texttt{useTracks()} aus der LiveKit React-Bibliothek, um sich auf LiveKit Räume zu abonnieren. Bei Änderungen werden die Komponenten automatisch neu gerendert, ohne dass die Applikation aktiv Zustandsabfragen durchführen muss.

Der Token-Abruf implementiert ein klassisches \textbf{Server-Side Token Pattern}: Die Frontend-Komponente \texttt{VoiceAgent} ruft die Next.js \gls{gls-API}-Route \texttt{/api/livekit/token} auf, die auf der Server-Seite läuft. Diese Route nutzt das \texttt{livekit-server-sdk} um einen \gls{gls-JWT} zu signieren, der folgende Informationen enthält: Eine eindeutige Identität der Teilnehmer (z.\,B. ``user''), eine zugeordnete Raum-ID (z.\,B. ``voice-agent-room'') und Berechtigungen (\texttt{canPublish}, \texttt{canSubscribe}, \texttt{canPublishData}). Der signierte Token wird dann an das Frontend zurückgegeben. Dies verhindert, dass das \gls{gls-API}-Geheimnis, welches für die Token-Signierung notwendig ist, im Frontend-Code sichtbar ist.

Das \textbf{Reconnect-Policy-System} ist kritisch für die Stabilität des Systems: Die \texttt{LiveKitRoom}-Komponente wird mit einer benutzerdefinierten \texttt{reconnectPolicy} konfiguriert, die exponentielles Backoff implementiert. Konkret wartet das System beim ersten Verbindungsabbruch 5 Sekunden, beim zweiten 10 Sekunden und beim dritten 15 Sekunden, bevor es nach 3 Versuchen aufgibt. Diese Wartezeiten sind notwendig, weil der Agent Worker Container beim Starten Zeit benötigt, um den LiveKit Raum zu betreten und die \gls{gls-KI}-Pipeline zu initialisieren. Ohne gro\ss zügige Wartezeiten würde der Browser-Client die Verbindung zu früh als gescheitert einstufen. Dieses konkrete Warten bietet einen Kompromiss zwischen Fehlertoleranz und Benutzer-Feedback. Nach 15 Sekunden wird dem Nutzer ein Fehler mitgeteilt, falls die Verbindung nicht hergestellt werden kann.

\noindent Die praktische Implementierung des Server-Side Token Pattern funktioniert wie folgt: Die Frontend-Komponente \texttt{VoiceAgent} wird initial mit dem ``Start''-Button in einer Lade-Phase gerendert. Beim Laden wird ein \texttt{useEffect}-Hook ausgeführt, der einen asynchronen \texttt{fetch()}-Request an die Next.js \gls{gls-API}-Route \texttt{/api/livekit/token} mit \texttt{Content-Type: application/json} und den Payload-Parametern \texttt{roomName} und \texttt{participantName} durchführt. Die serverseitige Route validiert zunächst diese Parameter und extrahiert dann die Umgebungsvariablen \texttt{LIVEKIT\_API\_KEY} und \texttt{LIVEKIT\_API\_SECRET} aus der Systemkonfiguration \texttt{.env.local}. Mit Hilfe des\\ \texttt{livekit-server-sdk}-Moduls wird ein neues \texttt{AccessToken}-Objekt erzeugt mit einer \gls{gls-TTL} von 10 Stunden, und mittels \texttt{at.addGrant()} werden die Berechtigungen gesetzt: \texttt{roomJoin: true} betritt den Raum, \texttt{room: roomName}, sowie die Publikations- und Abonnement-Berechtigungen. Der signierte Token wird dann via \texttt{at.toJwt()} kodiert und zusammen mit der \texttt{LIVEKIT\_URL} \texttt{ws://localhost:7880} lokal oder einer öffentlichen URL in der Produktivumgebung, an das Frontend zurückgegeben. Das Frontend speichert beide Werte im React State und rendert dann die Schaltfläche ``Voice Agent starten''. Erst wenn der Nutzer diese Schaltfläche anklickt, wird der \texttt{connect}-Parameter auf \texttt{true} gesetzt, und die \texttt{LiveKitRoom}-Komponente initiiert die tatsächliche WebSocket-Verbindung mit dem Token und der URL. Bei der Nutzung der \gls{gls-SIP}-Integration ist lediglich das Ausführen des Agent Workers notwendig.

\noindent Die Verbindungserkennung erfolgt mittels des \texttt{useConnectionState()}-Hooks, der die folgende Zustandsmaschine überwacht: Der \texttt{ConnectionState} ist initial \texttt{Disconnected}, wechselt auf \texttt{Connecting}, wenn der Browser versucht, die Verbindung mit dem WebSocket herzustellen, und wird \texttt{Connected}, sobald der Handshake erfolgreich ist. Wenn der Verbindungsabbruch erfolgt, springt der Status zurück auf \texttt{Disconnected}, und die \texttt{reconnectPolicy}-Logik wird aktiviert.\newline Die Policy-Funktion \texttt{nextRetryDelayInMs(context)} wird mit einem Kontextobjekt aufgerufen, das \texttt{context.retryCount} enthält. Wenn \texttt{retryCount < 3}, berechnet die Funktion \texttt{5000 * (retryCount + 1)}, was 5 s (retry 0), 10 s (retry 1), bzw. 15 s (retry 2) ergibt). Wenn der Kontext erfolgreich verbunden wurde, wird die Richtlinie nicht erneut aktiviert. Wenn alle drei Versuche fehlschlagen, gibt die Funktion \texttt{null} zurück, was signalisiert, dass keine weiteren Reconnect-Versuche stattfinden sollen, und eine Fehlerkomponente wird im Frontend angezeigt.

\looseness=-1
Der Agent Worker ist in Python 3.12 implementiert und nutzt das LiveKit Agents \gls{gls-SDK} (v0.8.0) als Framework für die Sprachverarbeitung. Die \gls{gls-KI}-Pipeline orchestriert vier spezialisierte Komponenten in Echtzeit, die wie folgt zusammenwirken: (1) Es läuft \textbf{Silero \gls{gls-VAD}} lokal ohne \gls{gls-API}-Calls und fungiert als Torwächter, um zu entscheiden, welche Audio-Blöcke Sprache enthalten und daher an Whisper gesendet werden sollten. (2) Die OpenAI Whisper \gls{gls-API} wird nur bei erkannter Sprache aufgerufen, um Audio-Blöcke in Text umzuwandeln. (3)  OpenAI GPT-4o-mini verarbeitet den erkannten Text und generiert eine Antwort, wobei Tokens gestreamt statt gesammelt werden. (4) OpenAI \gls{gls-TTS}  synthetisiert die \gls{gls-LLM}-Antwort in Audio, wobei Streamingmodus verwendet wird, um Latenz zu minimieren.

\textbf{Metrik-Sammlung und Performance-Tracking:}\\
Der Agent Worker implementiert umfassendes Performance-Tracking über eine \texttt{UsageCollector}-Klasse, die relevante Metriken während eines Anrufs erfasst. Für jeden Anruf werden folgende Daten gesammelt: Anruf-ID und Raumname zur Identifikation, Start- und Endezeitstempel für die Berechnung der Antwortzeiten und der Dauer. Die Sammlung umfasst Teilehmer-IDs aller Gesprächsteilnehmer, \gls{gls-STT} und \gls{gls-STT}-Latenzen pro Äu\ss erung. Au\ss erdem wird die Anzahl generierter Audiodaten sowie ein Verbindungsabbruchsgrund und -Zeitpunkt festgehalten.

Die Metriken werden während des Anrufs in Echtzeit erfasst und bei Anrufende als \gls{gls-JSON}-Zusammenfassung in einen separaten Ordner exportiert. Dies ermöglicht nachträgliche Anrufanalysen zur Optimierung der Pipeline und zur Identifikation von Leistungsengpässen. Die \texttt{summarize()}-Methode berechnet Durchschnittswerte für alle Latenzen und aggregiert Gesamtaudiodaten, was für Kostenkalkulationen und zur Optimierung des Systems essenziell ist. Diese Werte werden in Kapitel \ref{chap:evaluation} nähergehend betrachtet.

\section{Konfiguration und Datenstrukturen}

Das System speichert und verarbeitet Daten über mehrere Schichten hinweg mit standardisierten Strukturen und ist über drei Hauptkonfigurationsdateien definiert: \texttt{docker-compose.yml} für die Orchestrierung, Asterisk-Konfigurationsdateien (\texttt{pjsip.conf}, \texttt{extensions.conf}, \texttt{modules.conf}) für das \gls{gls-SIP}-Routing und
LiveKit-Konfigurationsdateien (\texttt{livekit-config.yaml}, \texttt{sip-config.yaml}) für den Media-Server und die \gls{gls-SIP}-Bridge.

\subsection{Docker-Compose Konfiguration und Service-Vernetzung}

 Die Datei \texttt{docker-compose.yml} orchestriert die gesamte Infrastruktur mit expliziten IP-Adressen im privaten Netzwerk \texttt{172.20.0.0/16}. Dieses isolierte Netzwerk-Setup wird als Bridge-Netzwerk konfiguriert, was bedeutet, dass alle Container intern über diese vorhersehbaren IP-Adressen miteinander kommunizieren können, während der Host nur die extern weitergeleiteten Ports sieht.

\looseness=-1
Die Service-zu-Service-Topologie ist wie folgt organisiert:\newline \textbf{Redis} läuft auf \texttt{172.20.0.2:6379} und wird sowohl von der LiveKit \gls{gls-SIP}-Bridge als auch von Asterisk für die Verwaltung des Sitzungsstatus und dessen Persistierung benötigt. \textbf{LiveKit Server} läuft auf \texttt{172.20.0.3} und exponiert intern Port 7880 für WebSocket/\gls{gls-HTTP}, 7881 für \gls{gls-TCP}/\gls{gls-TURN} und 7882 für \gls{gls-UDP}. Extern werden die Ports 50000-50100 für \gls{gls-WebRTC}-Clients weitergeleitet. \textbf{Asterisk \gls{gls-PBX}} läuft auf \texttt{172.20.0.10:5060} für \gls{gls-SIP}-Signalisierung und Ports 10200-10299 werden für \gls{gls-RTP} verwendet. Die \textbf{LiveKit \gls{gls-SIP}-Bridge} läuft auf \texttt{172.20.0.5} und lauscht intern auf Port 5060, wird aber extern auf 5069 weitergeleitet, um einen Portkonflikt mit Asterisk zu vermeiden.

\looseness=-1
Ein kritischer Aspekt ist die \textbf{Planung des Portnutzungsbereichs}: Asterisk nutzt \gls{gls-RTP}-Ports 10200-10299 für bidirektionales Audio zu externen \gls{gls-SIP}-Providern. Die LiveKit \gls{gls-SIP}-Bridge nutzt Ports 10000-10100 für lokale \gls{gls-RTP}-Kommunikation mit Asterisk und zur Konvertierung nach \gls{gls-WebRTC}. Der LiveKit Server nutzt die Ports 50000-50100 für externe \gls{gls-WebRTC}-Verbindungen vom Browser. Diese sorgfältige Separation verhindert Port-Konflikte.

Die Konfiguration setzt zusätzliche Umgebungsvariablen:\newline \texttt{LIVEKIT\_ANNOUNCE\_IP} teilt dem LiveKit Server mit, unter welcher öffentlichen IP-Adresse er erreichbar ist, damit WebRTC-Clients sich korrekt mit ihm verbinden können. Aus datenschutzrechtlichen Aspekten wurde hier eine Beispiel-IP aus dem RFC 5737 Dokumentations-Bereich verwendet. Für die \gls{gls-SIP}-Bridge werden interne Docker-IP-Adressen \texttt{SIP\_INTERNAL\_IP} verwendet, die vom Bridge-Netzwerk automatisch geroutet werden.

Da das System eingehende \gls{gls-SIP}-Anrufe vom Provider empfangen muss, ist die Konfiguration von \gls{gls-NAT}-Traversierung und Portfreigaben auf dem Router kritisch. Für den prototypischen Betrieb wurden folgende Portfreigaben auf dem Router konfiguriert: Port 5060/\gls{gls-UDP} für \gls{gls-SIP}-Signalisierung wird an die interne Server-IP weitergeleitet, damit \gls{gls-SIP} INVITE-Pakete vom Provider empfangen werden können. Die \gls{gls-RTP}-Port-Range 10200-10299/\gls{gls-UDP} wird ebenfalls weitergeleitet, um bidirektionales Audio zwischen dem Provider und
Asterisk zu ermöglichen. Zusätzlich werden die Ports 7880/\gls{gls-TCP}, 7881/\gls{gls-TCP} (LiveKit WebSocket/TCP-Fallback) sowie 50000-50100/\gls{gls-UDP} \gls{gls-WebRTC} freigegeben, damit externe \gls{gls-WebRTC}-Clients sich mit dem LiveKit Server verbinden können. Ohne diese Portfreigaben würde die Firewall des Routers eingehende Verbindungen blockieren, was dazu führen würde, dass Anrufe vom Provider nicht ankommen oder das Audio aufgrund fehlender \gls{gls-RTP}-Durchleitung nicht funktioniert.

\subsection{Asterisk-Konfiguration: PJSIP und Anrufweiterleitung}

Asterisk wird über drei Konfigurationsdateien gesteuert, die im Container unter \texttt{/etc/asterisk/} bereitgestellt werden:

\textbf{Transport und Trunk-Definition:} (\texttt{pjsip.conf})\\
Die Datei definiert zunächst einen \gls{gls-UDP}-Transport (\texttt{[transport-udp]}) auf Port 5060 mit explizit konfigurierten externen IP-Adressen. Dies sind die\\ (\texttt{external\_media\_address}) und die (\texttt{external\_signaling\_address}), wobei beide auf die externe IP-Adresse des Server-Hosts zeigen müssen. Dies ist auf das Single-Host-Szenario mit nur einer öffentlichen IP-Adresse zurückzuführen.

Der \gls{gls-SIP}-Trunk ist als Registrierung (\texttt{['Provider-Name']}) konfiguriert. Es muss sich beim Provider-Server (\texttt{sip:sip.provider.de}) unter einem Benutzernamen\\ (\texttt{sipuser\_xxxxx\_00}), über die Digest-Zugriffsauthentifizierung, angemeldet werden. Diese Zugangsdaten sind je nach Provider unterschiedlich und müssen mit dem entsprechenden Konto übereinstimmen. Aus Gründen des Datenschutzes und der Anonymität sind die tatsächlichen Zugangsdaten hier nicht einsehbar.

Ein Endpunkt (\texttt{[provider-endpoint]}) definiert die Behandlung eingehender Anrufe über den Provider: Erlaubte Codecs sind \texttt{ulaw, alaw, Opus}. Der Parameter\\ \texttt{context=from-provider} leitet eingehende Anrufe in den entsprechenden Dialplan-Kontext weiter. Parameter wie \texttt{rewrite\_contact=yes}, \texttt{force\_rport=yes}, und\\ \texttt{rtp\_symmetric=yes} sind \gls{gls-NAT}-Traversierungstechniken, die notwendig sind, wenn der Provider hinter einem Netzwerk, das mit \gls{gls-NAT} konfiguriert ist, liegt. Somit werden Asymmetrien vermieden.

Ein zweiter Endpunkt (\texttt{[livekit]}) ist für die interne Docker-Kommunikation mit der LiveKit \gls{gls-SIP}-Bridge konfiguriert. Dieser Endpunkt hat bewusst keine Authentifizierung, da er nur intern läuft. Die \texttt{media\_address='IP-Adresse'} zeigt auf die Docker-IP des Asterisk-Containers, und der \texttt{contact} zeigt auf die statische IP der LiveKit \gls{gls-SIP}-Bridge (\texttt{172.20.0.5:5060}).

Ein \texttt{['provider-identifizierung']}-Block listet den IP-Adressbereich des Providers (\texttt{'Externe-IP'}) auf. Dies ermöglicht es Asterisk, eingehende \gls{gls-SIP}-INVITE-Pakete vom Provider automatisch dem richtigen Endpunkt (\texttt{'provider-endpunkt'}) zuzuordnen, ohne dass eine explizite Digest-Zugriffsauthentifizierung für jedes Paket nötig ist.

\textbf{Dialplan und Routing:} (\texttt{extensions.conf})\\
Der Dialplan ist das Herzstück des Routings. Er besteht aus Kontexten, Extensions und Befehlen:

Der Kontext \texttt{['from-provider']} definiert, wie eingehende Anrufe vom Provider behandelt werden. Wenn ein Anruf über den Provider eingeht (identifiziert durch den \texttt{['provider-identifizierung']}-Block), durchläuft er diesen Kontext. Eine spezifische Extension \texttt{+491234567890} repräsentiert die \textbf{\gls{gls-DID}}-Nummer, die dem System beim Provider zugewiesen ist. Der \texttt{Dial()}-Befehl leitet den Anruf weiter: \texttt{Dial(\gls{gls-PJSIP}/\linebreak[0]\$\{CALLERID(num)\}@livekit,120,g)}. Dies bedeutet: "Nutze den PJSIP-Kanal, sende die Anrufer-ID des ursprünglichen Anrufers unverändert, adressiere den Endpunkt 'livekit', warte bis zu 120 Sekunden auf Annahme, und 'g' bedeutet, erlauben sie dem Angerufenen, das Gespräch aufzulegen".

Ein Catch-all-Fallback (\texttt{exten => \_X.,1}) handhabt unerwartete Extensions und routet sie ebenfalls an LiveKit.


\textbf{Selective Module Loading:} (\texttt{modules.conf})\\
Diese Datei definiert, welche Asterisk-Module geladen werden. Die Module\\ (\texttt{res\_pjsip.so}, \texttt{res\_pjsip\_session.so}, \texttt{res\_pjsip\_registrar.so},\\ \texttt{res\_pjsip\_endpoint\_identifier\_ip.so}), sind notwendig, um die moderne \gls{gls-SIP}-Implementierung bereitzustellen. Das Dialplan-Engine-Modul (\texttt{pbx\_config.so}) wird geladen, um die \texttt{extensions.conf} zu parsen.

Die folgenden Codec-Module werden explizit geladen: \texttt{codec\_ulaw.so}, \texttt{codec\_alaw.so},\newline \texttt{codec\_opus.so}. Diese behandeln die Kompression und Dekompression von Audio mit verschiedenen Kodierungsschemata.

Entscheidend ist: \textbf{Audio-Hardware-Module werden explizit NICHT geladen}. \texttt{noload => chan\_alsa.so} deaktiviert die Advanced Linux Sound Architecture, \texttt{noload => chan\_console.so} deaktiviert die Konsolen-Audio. Dies ist notwendig, da das System in einem Docker-Container ohne lokale Audio-Hardware läuft. Stattdessen wird das gesamte Audio über \gls{gls-RTP}/\gls{gls-WebRTC} transportiert.

\subsection{LiveKit-Konfiguration}

Die Datei \textbf{\texttt{livekit-config.yaml}} konfiguriert den LiveKit Server:

\begin{itemize}
  \item \texttt{port: 7880}: Hauptport für \gls{gls-HTTP}/WebSocket-Verbindungen. Clients (Browser,
Agent Worker) stellen hier ihre Verbindung her.
  \item \texttt{bind\_addresses: 0.0.0.0}: Lauscht auf allen Netzwerkschnittstellen des Containers
  \item \texttt{rtc.port\_range\_start/end: 50000-50100}: \gls{gls-UDP}-Ports für \gls{gls-RTC}-Media (Audio/Video-Streams)
  \item \texttt{rtc.tcp\_port: 7881}: Fallback-Port für Clients, die über Firewalls nur \gls{gls-TCP} verwenden können
  \item \texttt{keys: devkey: secret}: Entwicklungs-API-Keys (nicht produktionsgerecht! Produktivumgebungen nutzen starke Keys)
  \item \texttt{room.auto\_create: true}: Neue Rooms werden automatisch erzeugt, wenn der erste Participant beitritt
  \item \texttt{room.empty\_timeout: 300}: Räume werden nach 5 Minuten Inaktivität automatisch gelöscht
  \item \texttt{redis.address: redis:6379}: Verbindung zu Redis für verteilte\newline Sitzungskoordination
\end{itemize}

\subsection{SIP-Bridge Konfiguration}
\label{sec:sip-bridge-config}

Die Datei \textbf{\texttt{sip-config.yaml}} konfiguriert die LiveKit \gls{gls-SIP}-Bridge:

\begin{itemize}
  \item \texttt{ws\_url: ws://livekit:7880}: WebSocket-Verbindung zum LiveKit Server (interner Docker-Name)
  \item \texttt{api\_key/api\_secret: devkey/secret}: Authentifizierung gegen den LiveKit Server (nicht produktionsgerecht! Produktivumgebungen nutzen starke Keys)
  \item \texttt{redis.address: redis:6379}: Für Trunk- und Zuteilungsregelung und Persistierung
  \item \texttt{log\_level: debug}: Ausführliche Protokollierung für Debugging von \gls{gls-SIP}-Problemen
\end{itemize}

Die LiveKit \gls{gls-SIP}-Bridge wird nicht über statische Konfigurationsdateien eingerichtet, sondern \textbf{programmgesteuert über die LiveKit \gls{gls-API}}. So können Trunks und Dispatch-Rules flexibel zur Laufzeit verwaltet werden. Die Einrichtung erfolgt über ein Python-Setup-Skript (\texttt{setup\_sip\_bridge\_v2.py}), das das \texttt{livekit-api} \gls{gls-SDK} verwendet.
\clearpage
\textbf{Inbound-Trunk-Registrierung:}\\
Ein \textbf{Inbound-Trunk} definiert, wie die \gls{gls-SIP}-Bridge eingehende Anrufe vom Provider empfängt. Das Setup-Skript erstellt einen Trunk mit folgenden Parametern: \texttt{name} (Identifikation, z.\,B. ``Plusnet Inbound''), \texttt{numbers} (Liste der \gls{gls-DID},\\ z.\,B.[\texttt{+491234567890}]), \texttt{auth\_username} und \texttt{auth\_password} (Providerzugangsdaten für die \gls{gls-SIP}-Authentication), \texttt{allowed\_addresses} (IP-Whitelist: leer = alle erlaubt) und \texttt{metadata} (JSON-String mit Provider-Informationen).

Die Implementierung nutzt die asynchrone \gls{gls-API}-Methode\\ \texttt{create\_sip\_inbound\_trunk()}, die ein \texttt{SIPInboundTrunkInfo}-Objekt als Parameter akzeptiert. Nach erfolgreicher Registrierung gibt die \gls{gls-API} eine eindeutige Trunk-ID zurück, die für spätere Referenzierungen verwendet wird.

\textbf{Erstellung der Zuteilungsregeln:}\\
Eine \textbf{Zuteilungsregel} definiert das Routing: Wohin soll ein eingehender Anruf geleitet werden? Das Setup-Skript erstellt eine Regel vom Typ \texttt{SIPDispatchRuleDirect}, die Anrufe direkt an einen LiveKit-Raum weiterleitet. Der Raumname folgt dem Pattern \texttt{sip-call-\{callID\}}, wobei \texttt{\{callID\}} zur Laufzeit durch eine eindeutige Anruferkennung ersetzt wird.

Die \texttt{create\_sip\_dispatch\_rule()}-Methode akzeptiert ein \texttt{SIPDispatchRule}-Objekt mit einem \texttt{dispatch\_rule\_direct}-Feld, das \texttt{room\_name} und optional einen \texttt{pin} für geschützte Räume enthält. Das prototypische System nutzt keine PIN, da die Authentifizierung über die Zugandsdaten des \gls{gls-SIP}-Trunk erfolgt.

\textbf{Aufräumen und Fehlerbehandlung:}\\
Das Setup-Skript implementiert eine robuste Fehlerbehandlung: Vor der Erstellung neuer Trunks werden alle existierenden Trunks und Zuteilungsregeln gelöscht\\ (\texttt{cleanup\_existing\_config()}), um Konflikte zu vermeiden. Dies ist notwendig, da die LiveKit \gls{gls-API} keine Duplikaterkennung bietet. Die Cleanup-Funktion listet alle existierenden Ressourcen via \texttt{list\_inbound\_trunk()} und \texttt{list\_dispatch\_rule()} auf und löscht sie mit den entsprechenden \texttt{delete}-Methoden.

Ein Verifikationsschritt (\texttt{verify\_configuration()}) prüft nach dem Setup, ob alle Ressourcen korrekt erstellt sind, und gibt eine Zusammenfassung der aktiven Trunks und Regeln aus. Dies ermöglicht schnelles Debugging bei Konfigurationsproblemen.

\textbf{Praktische Nutzung:}\\
Das Setup-Skript wird einmalig nach dem Deployment ausgeführt:\newline \texttt{python python/setup\_sip\_bridge\_v2.py}. Es liest Credentials aus Umgebungsvariablen (\texttt{SIP\_PROVIDER\_USERNAME}, \texttt{SIP\_PROVIDER\_PASSWORD}, \texttt{SIP\_DID\_NUMBER}) und konfiguriert die \gls{gls-SIP}-Bridge automatisch. Nach erfolgreichem Setup können eingehende Anrufe auf die konfigurierte \gls{gls-DID} direkt an den Agent Worker weitergeleitet werden, ohne manuelle Eingriffe.

\subsection{Strukturierte Protokollierung und Eventverfolgung}

Alle Dienste schreiben strukturierte Events in \textbf{\gls{gls-NDJSON}}-Format. Jedes Event ist ein separates \gls{gls-JSON}-Objekt auf einer eigenen Zeile:

\begin{center}
\texttt{\{call\_id:sip-call-1234567890,timestamp:2026-01-16T14:30:45.123Z,service:agent-worker,event-type:stt-transcription,transcript:"Hallo",latency-ms:523}
\end{center}

Die standardisierten Felder sind:

\begin{itemize}
  \item \textbf{call\_id}: Eindeutige Identifier für dienstübergreifendes Tracking (z.\,B. \texttt{sip-call-incoming-main-1234567890})
  \item \textbf{timestamp}: ISO 8601 Format mit Millisekunden-Auflösung
  \item \textbf{service}: Asterisk, Livekit-sip, Livekit-server, Agent Worker, Frontend
  \item \textbf{log\_level}: INFO, WARN, ERROR
  \item \textbf{event\_type}: Klassifizierung: \texttt{room\_joined},
\texttt{stt\_transcription}, \texttt{llm\_response}, \texttt{tts\_synthesis}, \texttt{call\_ended}, \texttt{api\_error}
  \item \textbf{metadata}: Event-spezifische Daten
\end{itemize}

Diese Struktur ermöglicht eine automatisierte Extraktion von wichtigen Metriken. Nachbearbeitungsanalysen der Anrufe können alle Events mit einer bestimmten \texttt{call\_id} sammeln und die gesamte Anrufhistorie mit Zeitstempeln rekonstruieren.

\subsection{Audio-Datenformat und Root Mean Square}

Das System nutzt durchgehend \gls{gls-PCM16} in 16-Bit-Auflösung mit 16 kHz Abtastrate \parencite{rfc3551}. Die Audio-Chunks sind standardisiert auf \textbf{160 Samples}, was exakt 20 Millisekunden Audio pro Chunk entspricht (Berechnung: 160 Samples / 16.000 Hz = 0,01 s).

Für jeden Chunk wird das \textbf{\gls{gls-RMS}}-Level berechnet. Der \gls{gls-RMS}-Wert ist ein Signalkennwert für die Signalstärke und ist mathematisch definiert als $\sqrt{\frac{1}{N}\sum_{i=1}^{N} x_i^2}$, wobei $x_i$ die Abtastwerte des Audiosignals sind. Der \gls{gls-RMS}-Wert wird gegen einen empirisch ermittelten Schwellwert (typischerweise 0,15 auf einer normalisierten Skala von 0,0 bis 1,0) verglichen, um zwischen Sprache und Stille zu unterscheiden \parencite{Czarnecki1992RMS}.

\subsection{Sitzungsverwaltung und Raumlebenszyklus}

LiveKit Räume werden als \gls{gls-JSON}-Objekte mit eindeutigen Identifikatoren verwaltet. Das System erstellt Raäume mit einem vorhersehbaren Muster:\newline \texttt{sip-call-\{dispatchRuleID\}-\{timestamp\}}, wobei \texttt{dispatchRuleID} der Kennung der Zuteilungsregel entspricht (z.\,B. incoming-main) und
\texttt{timestamp} der Unix-Zeit beim Anrufeingang entspricht.

Jeder Raum enthält genau zwei Teilnehmer: Erstens den Anrufer, identifiziert über seine Anrufer-ID (z.\,B. +491234567890), und zweitens den \textbf{Agent Worker}, identifiziert durch seinen Container-Hostname (z.\,B. livekit-agent-1).

Metadaten werden im Raum als \gls{gls-JSON}-Objekt gespeichert, um späteren Abfragen Kontext zu geben:
\begin{itemize}
  \item \textbf{source}: plusnet (extern \gls{gls-SIP}) oder web (intern Browser)
  \item \textbf{sip\_trunk}: Name der \gls{gls-SIP}-Trunk-Konfiguration
  \item \textbf{dispatch\_rule\_match}: Welche Zuteilungsregel diesen Anruf weiterleitet (z.\,B. incoming-main)
  \item \textbf{caller\_id}: Für spätere Ablaufverfolgung
  \item \textbf{did\_dialed}: Die \gls{gls-DID}-Nummer, die der Anrufer gewählt hat
  \item \textbf{join\_timestamp}: Zeit nach ISO 8601 des Raumbeitritts
\end{itemize}

Diese Metadaten sind essenziell für Compliance (Rechnung, Auditierung), Debugging (Rekonstruktion von Problemen) und
Reporting (Anrufhistorie).

\section{Zusammenfassung und Deployment}

Die Implementierung des prototypischen Systems umfasst fünf spezialisierte Komponenten (LiveKit Server, LiveKit \gls{gls-SIP}-Bridge, Asterisk, Redis, Agent Worker). Sie stellt ein Next.js-Frontend mit \gls{gls-WebRTC}-Integration bereit und enthält umfassende Konfigurationsdateien für \gls{gls-SIP}-Routing, Medienserver und Sitzungsverwaltung. Die Abhängigkeiten werden über npm (\texttt{package.json}) für das Frontend und pip (\texttt{requirements.txt}) für den Backend-Agent verwaltet. Die \gls{gls-SIP}-Bridge-Konfiguration erfolgt über ein Python-Skript, das die LiveKit \gls{gls-API} für die Registrierung von Trunks und die Einrichtung von Zuteilungsregeln verwendet. Die Metrikverfolgung über die \texttt{UsageCollector}-Klasse ermöglicht Metrikanalysen zur Performance-Optimierung des Systems.

Detaillierte Anweisungen zur Installation, Konfiguration und Nutzung des Systems befinden sich im Anhang~\ref{ch:anhang-deployment}. Dort sind vollständige Befehle für die Installation der Abhängigkeiten, das Starten der Docker Container, das \gls{gls-SIP}-Bridge-Setup sowie Hinweise für auftretende Probleme dokumentiert.

% Kapitel 6: Evaluation
% Ergebnisse der SIP-Integration.
\chapter{Evaluation}
\label{chap:evaluation}

In dieser Evaluation werden lediglich Tests aus der \gls{gls-SIP}-Integration behandelt. Eine quantitative Auswertung der Web\gls{gls-UI} selbst ist nicht das Ziel der Evaluation.

\section[SIP-Evaluation]{\gls{gls-SIP}-Evaluation}
Die \gls{gls-SIP}-Anbindung wurde im Gegensatz zur Web\gls{gls-UI} detailliert verfolgt und ausgewertet. Die Signalisierungsebene und der Medienpfad funktionieren stabil: \gls{gls-SIP}-Trunk-Registration beim Provider, Zuteilungsregeln und die automatische Raumerstellung greifen zuverlässig. Eingehende Anrufe werden korrekt erkannt und an den \gls{gls-KI}-Agenten weitergeleitet.

Nach den durchgeführten \gls{gls-SIP}-Testanrufen zeigt sich, dass:

\begin{itemize}
  \item Bidirektionale Audioübertragung ohne \gls{gls-RTP}-Timeout erfolgt
  \item Mehrfache Gesprächsumläufe stabil verarbeitet (keine Abbrüche oder Timeouts)
  \item Anrufer Begrü\ss ung und Sprachausgabe des Agents deutlich empfangen
  \item Audio-Qualität und Latenz im erwarteten Bereich liegen
\end{itemize}

\subsection{Integration, Durchführung und Testszenarien}

Die \gls{gls-SIP}-Integration wurde in einer kontrollierten \gls{gls-LAN}-Umgebung mit Docker Desktop betrieben, während die Testanrufe selbst über das öffentliche Mobilfunknetz (Telekom) durchgeführt wurden. Aus der Gesamtheit der Testanrufe wurden exemplarisch 5 Anrufe für eine detaillierte Analyse ausgewählt. Die folgenden Beispiele zeigen die Vielfalt der getesteten Szenarien:

\begin{itemize}
  \item \textbf{Anruf 1:} Fragen zu physikalischen Dimensionen (Fu\ss ballfeld, Handball-Umfang): 3 Turns, 54,97s Dauer
  \item \textbf{Anruf 2:} Geographie und Zeitberechnungen (tiefste Meeresstelle): 3 Turns, 56,75s Dauer
  \item \textbf{Anruf 3:} Infrastruktur-Fragen (Deutsche Autobahnen): 3 Turns, 49,0s Dauer
  \item \textbf{Anruf 4:} Tier-Rekorde (schnellstes/langsamstes Tier): 4 Turns, 46,96s Dauer
  \item \textbf{Anruf 5:} Technologie-Abkürzungen (IT-Bereiche): 3 Turns, 37,96s Dauer
\end{itemize}

Die Testanrufe wurden bewusst mit unterschiedlicher Komplexität und Anzahl an Turns gestaltet , um eine Varianz in der Ergebnisdarstellung zu erreichen. Alle Anrufe wurden regulär durch Auflegen des Anrufers beendet (CLIENT\_INITIATED disconnect).

\subsection{Metriken und Analyse}

Aus den durchgeführten Testanrufen wurden fünf Anrufe für eine detaillierte Metrikanalyse ausgewählt. Sitzungen ohne Turns (z.\,B. kurz abgebrochene Rufaufbauten durch die Testperson) wurden von der Aggregation ausgeschlossen.

Zur qualitativen Einordnung der \gls{gls-STT}-Robustheit werden exemplarisch die in den Logs erkannten Nutzerfragen der fünf Testanrufe aufgeführt; dabei ist eine einzelne Fehlinterpretation sichtbar, die die Grenzen der Erkennung in Telefonszenarien verdeutlicht:
\begin{itemize}
  \item \textbf{Anruf 1 (Physikalische Dimensionen):} „Hallo, wie groß ist ein Fu\ss ballfeld?“, „Gib mir den Umfang für einen Handball.“, „Danke.“
  \item \textbf{Anruf 2 (Geographie/Zeiten):} „Gib mir die tiefste Stelle im Meer.“, „Wie lange würde es dauern, dort hinunterzutauchen?“, „Danke.“
  \item \textbf{Anruf 3 (Infrastruktur):} „Wie lang ist die längste Autobahn in Deutschland?“, „Welche ist die kürzeste?“, „Danke.“
  \item \textbf{Anruf 4 (Tier-Rekorde):} „Welches ist das schnellste Tier der Welt?“, \textit{„Welche Systemstelle?“} (Fehlinterpretation), „Das Langsamste.“, „Danke.“
  \item \textbf{Anruf 5 (Technologie-Abkürzungen):} „Was ist die Abkürzung IT?“, „Gib mir einen Bereich davon.“, „Danke.“
\end{itemize}

Die folgende Tabelle~\ref{tab:sip-metriken-aggregiert} zeigt die aggregierten Messwerte über diese fünf exemplarischen Testanrufe hinweg:

\begin{table}[H]
  \centering
  \begin{tabular}{|l|c|c|c|c|}
    \hline
    \textbf{Metrik} & \textbf{Mittelwert} & \textbf{Median} & \textbf{Min} & \textbf{Max} \\ \hline
    Anruf-Dauer (in s) & 49,13 & 49,00 & 37,96 & 56,75 \\ \hline
    Turns pro Anruf & 3,2 & 3 & 3 & 4 \\ \hline
    STT-Latenz (Ø in s) & 0,86 & 0,78 & 0,52 & 1,19 \\ \hline
    Transcript-Delay (Ø in s) & 1,30 & 1,31 & 1,05 & 1,74 \\ \hline
    TTS-Latenz (Ø in s) & 1,45 & 1,41 & 1,15 & 1,83 \\ \hline
    Gesamt Audio (s) & 21,15 & 22,35 & 10,73 & 26,78 \\ \hline
    Audio-Bytes (Ø in KiB) & 495,6 & 523,7 & 251,6 & 627,6 \\ \hline
  \end{tabular}
  \caption{Aggregierte Metriken und Latenzen der \gls{gls-SIP}-Tests}
  \label{tab:sip-metriken-aggregiert}
\end{table}

% --- SIP-Aggregat: Kurzverweis auf Statistik, Fokus auf Netzwerk- und Protokollaspekte ---
Die Analyse der aggregierten Daten zeigt konsistente Leistungswerte über die fünf detailliert ausgewerteten Testanrufe hinweg. Die \gls{gls-STT}-Latenz variiert zwischen 0,52 s im besten Fall und 1,19 s bei längeren Prompts. Im \gls{gls-SIP}-Kontext treten zusätzliche Latenzen durch Netzwerk- und Protokollaspekte auf, insbesondere bei der Transkriptverzögerung, die mit durchschnittlich 1,30 s deutlich über der reinen \gls{gls-STT}-Latenz liegt \parencite{ITU-T-G114}.

Die \gls{gls-TTS}-Latenz bleibt mit durchschnittlich 1,45 s über alle Anrufe hinweg konsistent \parencite{openai_whisper_docs}. Die durchschnittlich verarbeiteten Audiodaten pro Anruf liegen mit rund 496 KiB im niedrigen Kilobyte-Bereich. Dies entspricht den Empfehlungen für effiziente Sprachübertragung mit modernen Codecs wie Opus \parencite{rfc6716}.

\subsection{Detailanalyse von Testanrufen}

Zur tieferen Einordnung werden im Folgenden die Kennzahlen der fünf exemplarisch ausgewählten Testanrufe dargestellt (Tabelle~\ref{tab:sip-calls-detail}):

\begin{table}[H]
  \centering
  \small
  \begin{tabular}{|l|c|c|c|c|c|}
    \hline
    \textbf{Metrik} & \textbf{Anruf 1} & \textbf{Anruf 2} & \textbf{Anruf 3} & \textbf{Anruf 4} & \textbf{Anruf 5} \\ \hline
    Dauer (in s) & 54,97 & 56,75 & 49,00 & 46,96 & 37,96 \\ \hline
    Turns (Anzahl) & 3 & 3 & 3 & 4 & 3 \\ \hline
    STT-Durchschnitt (in s) & 1,19 & 0,77 & 0,52 & 1,07 & 0,78 \\ \hline
    TTS-Durchschnitt (in s) & 1,83 & 1,58 & 1,29 & 1,15 & 1,41 \\ \hline
    Audio (in s) & 25,07 & 26,78 & 22,35 & 20,81 & 10,73 \\ \hline
    Bytes (in KiB) & 587,6 & 627,6 & 523,7 & 487,7 & 251,6 \\ \hline
    STT-Fehlinterpretationen (Anzahl) & 0 & 0 & 0 & 1 & 0 \\ \hline
  \end{tabular}
  \caption{Detaillierte Metriken der \gls{gls-SIP}-Tests}
  \label{tab:sip-calls-detail}
\end{table}

% --- Detailanalyse: Fokus auf Gesprächsdynamik und Nutzerverhalten ---
Die Detailanalyse der einzelnen Anrufe zeigt Unterschiede in den gemessenen Leistungsmetriken. In Anruf 4 trat eine einzelne Fehlinterpretation des Transkripts auf, während in den übrigen Anrufen keine weiteren Auffälligkeiten dokumentiert wurden. Anruf 1 und Anruf 4 weisen höhere \gls{gls-STT}-Latenzen auf (1,19 s bzw. 1,07 s), während Anruf 3 mit 0,52 s die geringste \gls{gls-STT}-Latenz erreicht. Die \gls{gls-TTS}-Latenz variiert zwischen 1,15 s und 1,83 s.

Anruf 5 weist mit einer Gesamtdauer von 37,96 s sowie der geringsten Audiodatenmenge (10,73 s und 257,6 KiB) die niedrigsten Werte auf. Die beobachteten Unterschiede verdeutlichen die Variabilität der Systemmetriken zwischen einzelnen Gesprächen, deren Länge und der Anzahl Turns \parencite{Data2025_AIDataQuality,ssrn5348245}.

\subsection{Dialogbasierte Latenzanalyse}

Um die Konsistenz der Leistung innerhalb einzelner Gespräche zu bewerten, wurde eine dialogbasierte Analyse durchgeführt (Tabelle~\ref{tab:turn-latencies}):

\begin{table}[H]
  \centering
  \begin{tabular}{|l|c|c|c|c|}
    \hline
    \textbf{Anruf} & \textbf{Turn 1 STT (s)} & \textbf{Turn 2 STT (s)} & \textbf{Turn 3 STT (s)} \\ \hline
    Anruf 1 & 1,02 & 2,10 & 0,46  \\ \hline
    Anruf 2 & 0,84 & 0,91 & 0,55  \\ \hline
    Anruf 3 & 0,82 & 0,38 & 0,35  \\ \hline
    Anruf 4 & 1,07 & 0,81 & 1,80  \\ \hline
    Anruf 5 & 0,68 & 1,10 & 0,54  \\ \hline
  \end{tabular}
  \caption{Dialogbasierte Latenzanalyse einzelner \gls{gls-SIP}-Tests}
  \label{tab:turn-latencies}
\end{table}

% --- dialogbasiert: Fokus auf Schwankungen und Robustheit, keine Wiederholung der Statistik-Argumentation ---
Die dialogbasierte Analyse offenbart eine erhebliche Varianz der \gls{gls-STT}-Latenz innerhalb einzelner Anrufe. In Anruf 1 steigt die Latenz im zweiten Turn auf 2,10 s an, sinkt jedoch im dritten Turn wieder auf 0,46 s ab. Anruf 4 weist mit 1,80 s im dritten Turn eine besonders hohe Varianz auf. Im Gegensatz dazu bleibt Anruf 3 durchgehend konsistent und unterschreitet in allen Turns die 0,85 s. Die Ratio zwischen dem schnellsten und dem langsamsten Turn liegt bei 6,1:1 (0,35 s vs 2,10 s), was im erwartbaren Bereich für cloudbasierte \gls{gls-KI}-Dienste liegt \parencite{whisper_quant_2025,ITU-T-G114}.

Alle Turns wurden erfolgreich verarbeitet, ohne dass Timeouts auftraten. Eine Fehlinterpretation wurde in Anruf 4 dokumentiert.

\subsection{Audioqualität und Medienpfad-Stabilität}

Die Stabilität des Medienpfads wurde durch kontinuierliche bidirektionale Audioübertragung über die gesamte Anrufdauer nachgewiesen. Folgende qualitative Beobachtungen wurden während aller \gls{gls-SIP}-Testanrufe dokumentiert:

\begin{itemize}
  \item \textbf{Codec:} Opus 48 kHz, 24 kbps \parencite{rfc6716}
  \item \textbf{Paketverlust:} In \gls{gls-RTP}-Statistiken/LiveKit-Logs keine Paketverluste beobachtet
  \item \textbf{Jitter:} Stabile Audioausgabe ohne starke Verzögerungen
  \item \textbf{Echo:} Keine Echo-Bildung beobachtet (subjektive Wahrnehmung)
  \item \textbf{Verständlichkeit:} Subjektiv geschätzte \gls{gls-MOS} Werte von 4,0--4,5 / 5,0 (gut bis sehr gut)
  \item \textbf{\gls{gls-RTP}-Timeout:} Kein Auftreten bei allen Testanrufe
\end{itemize}

% --- Audioqualität: Fokus auf subjektive und objektive Bewertung, Codec, Netzwerkstabilität ---
Die konsistente Audioqualität über alle Testanrufe hinweg bestätigt die Stabilität des Medienpfads. Die \gls{gls-MOS}-Einschätzung ist subjektiv und basiert nicht auf standardisierten ITU-T-Hörversuchen. Neben der technischen Stabilität ist insbesondere die subjektive Verständlichkeit ein zentrales Qualitätsmerkmal für den Praxiseinsatz \parencite{ITU-T-G114}. Die Nutzung des Opus-Codecs und das Fehlen von Paketverlusten oder Jitter unterstreichen die objektive und subjektive Güte der Audioübertragung \parencite{rfc6716}.

\section{Hypothesenbezogene Evaluationsergebnisse}
\label{sec:hypothesen-resultate}

Die erhobenen Daten ermöglichen eine direkte Prüfung der in Kapitel~\ref{chap:einleitung} formulierten Hypothesen. Im Folgenden werden die Ergebnisse kurz dokumentiert, ohne inhaltliche Interpretation:

\paragraph{H1 (Funktionalität):} Gestützt durch Daten.
\begin{itemize}
  \item Keine funktionellen Fehler oder \gls{gls-RTP}-Timeouts über alle Testanrufe
  \item Eine \gls{gls-STT}-Fehlinterpretation wurde dokumentiert (Anruf 4)
  \item Bidirektionale Audioübertragung ohne Medienpfad-Fehler
  \item Alle Verbdindungsabbrüche regulär durch Auflegen des Anrufers (CLIENT\_INITIATED)
\end{itemize}

\paragraph{H2 (Latenz):} Gestützt durch Daten.
\begin{itemize}
  \item \gls{gls-STT}-Latenz: Ø 0,86\,s (Median: 0,78\,s, Spanne: 0,52--1,19\,s)
  \item \gls{gls-TTS}-Latenz: Ø 1,45\,s (Median 1,41\,s, Spanne: 1,15--1,83\,s)
  \item Transkriptverzögerung: Ø 1,30\,s
\end{itemize}

\paragraph{H3 (Ressourceneffizienz):} Gestützt durch Daten.
\begin{itemize}
  \item \gls{gls-CPU}-Auslastung: Ø 11{,}2\,\% (Max. 26\,\%)
  \item\gls{gls-RAM}-Nutzung: 11{,}2--12{,}1\,GB bei verfügbaren 16\,GB
\end{itemize}

\paragraph{H4 (Skalierbarkeit):} Gestützt durch Architekturdesign.
\begin{itemize}
  \item Modulare Komponentenstruktur
  \item Flexible Ressourcenzuteilung durch Containerbasierung
  \item Reserve in \gls{gls-CPU} und \gls{gls-RAM} für Erweiterbarkeit
\end{itemize}

\paragraph{H5 (STT-Robustheit):} Qualitativ gestützt durch Daten.
\begin{itemize}
  \item Die Mehrheit der deutschsprachigen Testprompts wurde korrekt erkannt
  \item Eine einzelne Fehlinterpretation wurde dokumentiert (Anruf 4)
  \item Es wurden keine systematischen Transkriptionsfehler dokumentiert
  \item Eine quantitative Messung der Word-Error-Rate wurde nicht durchgeführt
\end{itemize}

\section{Bedrohungen der Validität}
\label{sec:threats-to-validity}

Die Validität der Evaluationsergebnisse unterliegt mehreren Einschränkungen, die bei der Interpretation der Daten berücksichtigt werden müssen.

\subsection{Interne Validität}
\textbf{Stichprobengrö\ss e:}\\
Die Evaluation basiert auf \gls{gls-SIP}-Testanrufen, wovon fünf Anrufe detailliert analysiert wurden. Für produktive Systeme wären grö\ss ere Testszenarien erforderlich, um statistische Signifikanz zu gewährleisten. Dies ist für das hier getestete prototypische System nicht umgesetzt.

\textbf{Testszenario-Diversität:}\\
Die Tests wurden in einem mit Docker betriebenen System in einer kontrollierten \gls{gls-LAN}-Umgebung durchgeführt, wobei die \gls{gls-SIP}-Anrufe über eine Mobilfunk-Verbindung initiiert wurden. Andere Architekturen können abweichende Ergebnisse liefern.

\textbf{Prompt-Komplexität:}\\
Die verwendeten Testprompts decken Themen des Allgemeinwissen ab (Geographie,
Physik, Technologie). Domänenspezifische Fachsprache oder stark akzentuierte Aussprache könnte die \gls{gls-STT}-Genauigkeit beeinflussen.

\subsection{Externe Validität}
\textbf{Hardware-Abhängigkeit:}\\
Die Messungen wurden auf einem spezifischen Testsystem (16\,GB \gls{gls-RAM},
8 Kerne) durchgeführt. Die Übertragbarkeit auf andere Hardware-Konfigurationen (z.\,B. ARM-Architekturen, Cloud-Instanzen) ist nicht geprüft.

\textbf{Cloud-\gls{gls-API}-Variabilität:}\\
Die gemessenen Latenzen für \gls{gls-STT} und
\gls{gls-TTS} hängen von der Auslastung der OpenAI-Server und der geografischen Distanz ab. Wiederholte Messungen zu anderen Zeitpunkten oder von anderen Standorten könnten abweichende Werte liefern.

\textbf{Telefonanbieter-Spezifika:}\\
Die \gls{gls-SIP}-Tests erfolgten mit einem spezifischen Provider und einem dazugehörigen \gls{gls-SIP}-Trunk. Andere Provider könnten abweichende Codec-Unterstützung oder Netzwerkcharakteristika aufweisen.

\subsection{Konstrukt-Validität}
\textbf{Subjektive \gls{gls-MOS}-Bewertung:}\\
Die \gls{gls-MOS}-Werte basieren auf subjektiver Einschätzung ohne standardisiertes \gls{gls-ITU-T}-Testprotokoll. Eine objektive Messung mittels \gls{gls-PESQ} oder \gls{gls-POLQA} ist nicht Bestandteil der Hypothesen dieser Arbeit.

\textbf{Word-Error-Rate-Messung:}\\
Die qualitative Beobachtung, dass alle Prompts korrekt erkannt wurden, ersetzt keine quantitative Messung der Word-Error-Rate mit Referenztranskripten. Dies ist kein Bestandteil der Hypothesen dieser Arbeit.

\subsection{Reliabilität}
\textbf{Reproduzierbarkeit:}\\
Alle Tests sind durch die bereitgestellten Komponenten und deren Konfigurationen reproduzierbar. Die Cloud-\gls{gls-API}-Abhängigkeit führt jedoch zu nicht-deterministischen Latenzen bei Wiederholungen.

\textbf{Logging-Präzision:}\\
Die Messungen der Latenzen basieren auf Log-Zeitstempeln mit Millisekunden-Präzision. Systematische Logging-Verzögerungen, die durch Region oder Systemabhängigkeiten entstehen, könnten die Messungen verfälschen.
\chapter{Diskussion, Limitationen und Fazit}\label{chap:diskussion}

Dieses Kapitel erörtert die wesentlichsten Ergebnisse aus der Evaluation und ordnet sie in den bestehenden Kontext der technischen Literatur ein. Der Fokus liegt auf der Bewertung, inwiefern die entwickelte Architektur und der Prototyp die anfangs formulierten Hypothesen erfüllen, sowie auf der Identifikation von Grenzen und Optimierungsmöglichkeiten. Die Ergebnisse werden mit Blick auf praktische Implikationen für produktive Deployments diskutiert.

Die Evaluationsergebnisse bestätigen, dass die entwickelte Architektur eine stabile bidirektionale Audioverbindung zwischen \gls{gls-SIP}-Telefonie und \gls{gls-WebRTC} ermöglicht. Dies zeigt sich in der konsistenten Codec-Verhandlung, der erfolgreichen Überwindung von Netzwerkhindernissen durch \gls{gls-NAT}-Traversierungsmechanismen und der stabilen Verschlüsselung von Medienströmen. Die Anwendung von Opus als Codec und die strikte Einhaltung der \gls{gls-ICE}-, \gls{gls-STUN}- und \gls{gls-TURN}-Standards entsprechen dabei anerkannten Best-Practices für echtzeitfähige Multimediakommunikation \parencite{rfc5245,rfc5766,rfc5764,rfc6716}.

Die gemessenen Systemlatenzen liegen über denen der klassischen IP-Telefonie, was durch die zusätzlichen Verarbeitungsschritte der \gls{gls-KI}-Pipeline (\gls{gls-STT}, \gls{gls-LLM}-Inferenz, \gls{gls-TTS}) zu erklären ist. Ein direkter Vergleich mit klassischen Telefoniestandards ist daher nur eingeschränkt aussagekräftig. Stattdessen sollten die Ergebnisse im Kontext asynchroner, dialogorientierter Anwendungen interpretiert werden.

Zusammengefasst belegen die Evaluationsergebnisse, dass die gewählte Architektur sowohl den aktuellen Stand der Forschung und praktischen Szenarien abbildet. Die implementierten Komponenten folgen den Best-Practices für Echtzeitkommunikation und containerbasierten Deployments. Allerdings identifiziert die Analyse auch deutliche Limitationen und Ansatzpunkte zur Verbesserung, die in den folgenden Abschnitten adressiert werden.

\section{Bewertung der Hypothesenvalidierung}

Die entwickelte Architektur adressiert die in der Einleitung formulierten Hypothesen systematisch. Im Folgenden wird dokumentiert, wie die Evaluation diese Annahmen validiert oder modifiziert hat:

\textbf{H1 (Funktionalität):} Validiert.\\ Die Tests zeigen keine Verarbeitungsfehler oder \gls{gls-RTP}-Timeouts; eine einzelne \gls{gls-STT}-Fehlinterpretation wurde dokumentiert. Die stabile Medienpfad-Integration (\gls{gls-ICE}, \gls{gls-STUN}, \gls{gls-TURN}) bestätigt die bidirektionale Audioübertragung über Systemgrenzen hinweg \parencite{rfc5245,rfc5389,rfc5766,rfc3711}.

\textbf{H2 (Latenz):} Partiell validiert.\\ Die gemessenen Latenzbeiträge der Sprachverarbeitung (\gls{gls-STT}, Transkriptverzögerung, \gls{gls-TTS}) liegen im erwartbaren Bereich für \gls{gls-KI}‑Dialoge und ermöglichen praxistaugliche Interaktionen. Eine vollständige Einordnung nach klassischen \gls{gls-VoIP}‑Standards ist nur eingeschränkt möglich, da externe Netzwerk‑/\gls{gls-SIP}‑Anteile nicht gemessen wurden \parencite{ITU-T-G114}.

\textbf{H3 (Ressourceneffizienz):} Validiert.\\ Die \gls{gls-CPU}-Auslastung von 11,2\,\% und das stabile Verhalten unter Testlast demonstrieren effiziente Ressourcennutzung auf Standard-Hardware im Sinne des \gls{gls-ETSI}-\gls{gls-NFV}-Paradigmas \parencite{ETSI-NFV,ETSI-REL}.

\textbf{H4 (Skalierbarkeit):} Validiert.\\ Die modulare Komponentenstruktur ermöglicht horizontale Skalierung durch Replikation, wie in produktiven Architekturen, gestützt durch aktuelle Literatur, üblich ist \parencite{ETSI-REL,Antonova2025}.

\textbf{H5 (STT-Robustheit):} Qualitativ validiert.\\ Die Transkriptionen zeigten in den Testszenarien überwiegend korrekte Erkennung. Eine Fehlinterpretation wurde dokumentiert. Für produktive Systeme wird eine erweiterte Analyse empfohlen, die in dieser Arbeit nicht durchgeführt wurde und kein Bestandteil der Hypothesen darstellt.

\section{Iterative Entwicklung und Erkenntnisse}

Die iterative Entwicklung zeigte, dass die Menge an Integrationspunktne und ereignisbasierte Architekturprinzipien die Systemstabilität und Ressourcennutzung beeinflussen. In einer frühen Phase wurde Baresip als leichtgewichtiger \gls{gls-SIP} User-Agent genutzt, um die Grundlagen von \gls{gls-SIP}-Registrierung, Authentifizierung und Medienpfad-Steuerung zu verstehen. Baresip ermöglichte es, die Funktionsmechaniken eines \gls{gls-SIP}-Trunks zu lernen, ohne dabei von umfangreicher Dialplan-Logik belastet zu sein. Die gewonnenen Erkenntnisse führten zur Auswahl von Asterisk als stabilem \gls{gls-SIP}-Backend mit reproduzierbarer Dialplan-Konfiguration und deterministischem Routing. Aufbauend darauf wurde LiveKit als Media-Server und Integrationsschicht für \gls{gls-WebRTC}-Funktionalität und \gls{gls-KI}-Sprachverarbeitung ausgewählt. Damit konnte die Architektur schrittweise von einem experimentellen Setup zu einem belastbaren Prototypen, der mit einem modernen Framework implementiert ist, überführt werden.

\section{Limitationen der Arbeit}

Die Generalisierbarkeit der Ergebnisse wird durch Faktoren eingeschränkt: Die Tests wurden in einer kontrollierten Laborumgebung durchgeführt, die nicht alle Aspekte von Produktivumgebungen im Unternehmenskontext abdeckt. Die Hardware-Ausstattung des Testfeldes ist spezifisch und kann sich von Deployment-Szenarien in anderen Infrastrukturen unterscheiden. Die Abhängigkeit von Cloud-basierten \gls{gls-KI}-\gls{gls-API}s bedeutet, dass die Ergebnisse an die Performance und Verfügbarkeit dieser externen Services gebunden sind \parencite{ReliabilityVirtualizedNetworks,ITU-T-G114}.

\section{Praktische Implikationen und Ausblick}

Für die Praxis zeigen die Ergebnisse, dass containerbasierte Architekturen ein vielversprechender Weg sind, um klassische Telefonsysteme mit modernen \gls{gls-KI}-Komponenten zu verbinden. Die Orientierung an etablierten \gls{gls-VoIP}-Standards und die Nutzung bewährter Open-Source-Komponenten bieten ein solides Fundament für produktive Deployments \parencite{Antonova2025}. Allerdings erfordern produktive Systeme zusätzliche Investitionen in Orchestrierungsmechanismen, Monitoring sowie Sicherheitsma\ss nahmen. Zukünftige Arbeiten könnten sich auf systematische Evaluationen von Edge-Deployments, multimodalen Erweiterungen und domänenspezifischen Testsets konzentrieren. Die Berücksichtigung der objektiven Nutzerwahrnehmung wird als zukünftige Herausforderung identifiziert \parencite{ITU-T-G114}.

\section{Schlussbemerkung}

Diese Arbeit hat die zentrale Forschungsfrage beantwortet: Die stabile Integration eines \gls{gls-LLM} in ein \gls{gls-SIP}-basiertes Telefonsystem, das mit einer containerbasierten Architektur bereitgestellt wird, ist technisch realisierbar und praktisch sinnvoll. Die prototypische Implementierung demonstriert, dass Open-Source-Ansätze nicht nur mit proprietären Plattformen konkurrieren können, sondern sogar strategische Vorteile bieten: Vermeidung von Vendor-Lock-in, datenschutzgerechte On-Premise-Varianten und domänenspezifische Anpassbarkeit.

Die gemessenen Leistungsmetriken aus der \gls{gls-SIP}-Evaluation liegen im akzeptablen Bereich für konversative Telefonieanwendungen. Entscheidend ist jedoch die Interpretation: Diese Werte reflektieren nicht technische Grenzen, sondern fundamentale Trade-offs zwischen Modellkomplexität, Latenz und Kosten.

Die identifizierten Limitationen sind transparent dokumentiert und skizzieren einen klaren Pfad für zukünftige Arbeiten. Die prioritären Forschungsrichtungen, wie zum Beispiel die lokale \gls{gls-LLM}-Inferenz, Edge-Deployment oder multimodale Erweiterung, können gezieltere Einblicke in die vielfältigen Themen geben, die sich mit künstlicher Intelligenz und Telefonsystemen beschäftigen.

\clearpage
\pagenumbering{roman}
\pagestyle{scrheadings}

\appendix
% ============================================================================
% ANHANG - Zusätzliche Dokumentation und Deployment-Anleitung
% ============================================================================
% Die \appendix Direktive in main.tex führt dazu, dass alle \chapter{} hier
% automatisch mit Buchstaben (A, B, C, ...) statt Nummern nummeriert werden.

\chapter{Use-Case: Hotline}
\thispagestyle{scrheadings}
\markboth{ANHANG}{ANHANG}
\label{ch:anhang-usecase}

Dieses Kapitel illustriert beispielhaft eine praktische Anwendung anhand eines praxisrelevanten Use-Cases.

\paragraph{Bezeichnung:}
Automatisierte Hotline mit \gls{gls-KI}-basierter Vorbearbeitung.

\paragraph{Beschreibung:}
Eingehende Anrufe werden automatisiert angenommen, transkribiert und durch ein \gls{gls-LLM} vorbearbeitet. Bei Bedarf erfolgt eine Eskalation an menschliche Agenten.

\paragraph{Akteure:}
An diesem Use-Case sind drei Parteien beteiligt: Der \textbf{Anrufer}, der sich extern über ein \gls{gls-SIP}-Client oder das Public Switched Telephone Network, das klassische öffentliche Telefonnetz, verbindet. Das \textbf{System}, das aus einem \gls{gls-SIP}-Gateway, eine Schnittstelle zur Medienübertragung, einer Speech-To-Text Komponente und dem \gls{gls-LLM} und Ticket-System besteht, sowie ein menschlicher \textbf{Support-Agent}, der bei komplexen Anfragen zur Eskalation bereitsteht.

\paragraph{Vorbedingungen:}
Um einen erfolgreichen Ablauf zu gewährleisten müssen mehrere Voraussetzungen erfüllt sein. Die \gls{gls-SIP}-Infrastruktur muss erreichbar und nach RFC 3261 konform konfiguriert sein. Der \gls{gls-LLM}-Service sollte verfügbar sein und
über konfigurierte Richtlinien zur Erkennung von Benutzerabsichten, verfügen. Au\ss erdem muss die Ticket-Datenbank erreichbar sein und die Agent-Authentifizierung über einen Identity-Service möglich sein.

\paragraph{Hauptablauf:}
Der typische Ablauf beginnt mit der \textbf{Anruf-Annahme}, bei der sich der Anrufer mit dem System verbindet. Das \gls{gls-SIP}-Gateway empfängt das INVITE und erstellt eine neue Room-Session (\textbf{Session-Erstellung}). Das System sendet anschlie\ss end eine automatische \textbf{Begrü\ss ungs-Ansage} via \gls{gls-TTS}. Während des Gesprächs transkribiert \gls{gls-ASR} die Sprache des Anrufers in Echtzeit (\textbf{Transkription}), woraufhin das \gls{gls-LLM} das Transkript analysiert und die Benutzerabsicht erkennt (\textbf{Intent-Erkennung}, z.\,B. "Technische Unterstützung"). Wichtig ist ebenfalls eine Aufforderung des \gls{gls-KI}-Agenten, sodass der verbundene Anrufer zulassen kann, dass personenbezogene Daten erhoben oder benötigt werden.

In der Phase der \textbf{intelligenten Verarbeitung} wird je nach erkanntem \textbf{Intent} eine von drei Aktionen eingeleitet: Bei einer einfachen Lösung, die mithilfe von Informationsbereitstellung erbracht werden kann, antwortet das \gls{gls-LLM} direkt mit relevanter Information via \gls{gls-TTS}. Alternativ wird das System dazu angeleitet, ein Support-Ticket zu erstellen und speichert das Transkript (Ticket-Erstellung). Bei komplexeren Anfragen wird ein Agent benachrichtigt und der Anrufer zur Warteschlange zur \textbf{Eskalation} weitergeleitet. Zum \textbf{Abschluss} wird das Gespräch entweder beendet oder ein Agent übernimmt, während das Ticket aktualisiert wird.

\paragraph{Alternativer Ablauf:}
Falls die \gls{gls-ASR} keinen Treffer liefert oder ein Verständnisproblem auftritt, wird ein zweiter Durchgang zur Support-Anfrage eingeleitet. Sollte erneut kein Treffer geliefert werden, wird der Anruf direkt an einen menschlichen Agenten weitergeleitet. Bei einem \gls{gls-LLM}-Timeout von mehr als 10\,s erfolgt eine automatische Eskalation mit entsprechender Ankündigung an den Anrufer. Aus Datenschutzgründen ist bei \gls{gls-DSGVO}-relevanten Verarbeitungen ein Aufzeichnungs-Opt-in vor der eigentlichen Verarbeitung erforderlich.

\paragraph{Erfolgskriterien:}
Der Ablauf gilt als erfolgreich, wenn der Anrufer eine automatische Begrü\ss ung innerhalb der von \gls{gls-ITU-T} G.114 empfohlenen Latenz empfängt. Das erfasste Transkript muss korrekt dem entsprechenden Ticket zugeordnet werden, während der Intent entweder erfolgreich erkannt wird oder eine Eskalation eingeleitet wird. Im Fall einer Eskalation sollte der menschliche Agent vollständigen Kontext erhalten, inklusive Transkript, erkanntem Intent und
Ticket-ID.

\begin{figure}[H]
\centering
\includegraphics[width=1.0\textwidth,height=0.63\textheight,keepaspectratio]{images/use_case_architektur.pdf}
\caption{Sequenzdiagramm UC-01: Automatisierte Hotline mit \gls{gls-KI}-gestützter Vorbearbeitung (drei Pfade: A) Automatische Lösung,
B) Ticket-Erstellung, C) Eskalation zu Agent).}
\label{fig:usecase_sequence}
\end{figure}


\chapter{Deployment und Inbetriebnahme}
\markboth{ANHANG}{ANHANG}
\label{ch:anhang-deployment}

Dieses Kapitel beschreibt die vollständige Installation und Inbetriebnahme des Systems zur Reproduzierbarkeit der Implementierung.

\section{Systemvoraussetzungen}
\label{sec:anhang-systemvoraussetzungen}

Folgende Software muss auf dem Host-System installiert sein:

\begin{itemize}
  \item \textbf{Node.js:} v20.11.0 LTS oder höher (für Next.js Frontend)
  \item \textbf{Python:} v3.12.1 oder höher (für Agent Worker)
   \item \textbf{Docker Desktop:} v4.26.1 oder höher (für Container-Orchestrierung, Entwicklung unter Windows/macOS)
  \item \textbf{Docker Engine:} 23.x oder höher (für produktives Hosting unter Linux, alternativ zu Docker Desktop)
  \item \textbf{npm:} v10.0.0 oder höher (Paketmanager für Node.js)
  \item \textbf{pip:} v23.0.0 oder höher (Paketmanager für Python)
\end{itemize}

\noindent Hardwareanforderungen:
\begin{itemize}
  \item Mindestens 16 GB RAM (empfohlen für parallele Container)
  \item Mindestens 4 CPU-Kerne (für \gls{gls-VAD} und Audio-Processing)
  \item Mindestens 10 GB freier Festplattenspeicher
  \item Stabile Internetverbindung (für OpenAI \gls{gls-API}-Zugriff)
\end{itemize}

\section{Installation}
\label{sec:anhang-installation}

\subsection{Schritt 1: Dependencies installieren}

\subsubsection{Frontend-Dependencies (npm):}

\begin{lstlisting}[language=bash,caption={Node.js-Pakete installieren}]
npm install
\end{lstlisting}

\noindent Dies installiert alle in \texttt{package.json} definierten Pakete inklusive LiveKit-Komponenten, Next.js, React und TypeScript.

\subsubsection{Backend-Dependencies (Python):}

\begin{lstlisting}[language=bash,caption={Python-Pakete installieren}]
python -m pip install -r requirements.txt
\end{lstlisting}

\noindent Dies installiert alle in \texttt{requirements.txt} definierten Pakete inklusive LiveKit Agents \gls{gls-SDK}, OpenAI-Plugins und Silero \gls{gls-VAD}.

\subsection{Schritt 2: Umgebungsvariablen konfigurieren}

Erstellen Sie eine \texttt{.env.local}-Datei im Projekt-Root mit folgenden Variablen:

\begin{lstlisting}[language=bash,caption={Umgebungsvariablen (.env.local)}]
# LiveKit Server
LIVEKIT_URL=ws://localhost:7880
LIVEKIT_API_KEY=<your-api-key>
LIVEKIT_API_SECRET=<your-api-secret>

# OpenAI API
OPENAI_API_KEY=sk-<your-openai-api-key>

# SIP Provider Credentials
SIP_PROVIDER_DOMAIN=<sip-provider-domain>
SIP_PROVIDER_USERNAME=<sip-username>
SIP_PROVIDER_PASSWORD=<sip-password>
SIP_DID_NUMBER=<your-did-number>
\end{lstlisting}

\textit{Hinweis: Ersetzen Sie die Platzhalter in spitzen Klammern mit Ihren tatsächlichen Werten. Der OpenAI \gls{gls-API}-Key ist unter \url{https://platform.openai.com/api-keys} erhältlich.}

\subsection{Schritt 3: Docker-Container starten}

\begin{lstlisting}[language=bash,caption={Docker-Container im Hintergrund starten}]
docker-compose up -d
\end{lstlisting}

\noindent Dies startet fünf Container: LiveKit Server (Port 7880), LiveKit \gls{gls-SIP}-Bridge (Port 5069), Asterisk (Port 5060), Redis (Port 6379) und den internen Docker-Bridge-Netzwerk-Manager. Die Container laufen dauerhaft im Hintergrund (\texttt{-d} Flag).

\subsection{Schritt 4: SIP-Bridge konfigurieren}

Nach dem Start der Container muss die \gls{gls-SIP}-Bridge mit dem Provider-Trunk konfiguriert werden:

\begin{lstlisting}[language=bash,caption={SIP-Trunk und Dispatch-Rules registrieren}]
python python/setup_sip_bridge_v2.py
\end{lstlisting}

\noindent Dieses Skript erstellt einen Inbound-Trunk für eingehende Anrufe vom Provider und eine Dispatch-Rule, die Anrufe an den LiveKit-Room \texttt{sip-call-\{callID\}} routet.

\subsection{Schritt 5: Agent Worker starten}

Der Agent Worker muss gestartet werden, um eingehende Anrufe zu verarbeiten:

\begin{lstlisting}[language=bash,caption={Voice Agent Worker starten}]
python python/agent_worker.py
\end{lstlisting}

\noindent Der Agent Worker verbindet sich mit dem LiveKit Server und wartet auf eingehende Room-Events. Alle Events und Metriken werden in \texttt{agent\_worker\_debug.log} geschrieben. Lassen Sie dieses Terminal-Fenster offen, solange das System laufen soll.

\subsection{Schritt 6: Frontend starten (optional)}

Für die Web-Benutzeroberfläche starten Sie den Next.js Development Server in einem separaten Terminal:
\clearpage
\begin{lstlisting}[language=bash,caption={Next.js Development Server starten}]
npm run dev
\end{lstlisting}

\noindent Das Frontend ist danach unter \url{http://localhost:3000} erreichbar und ermöglicht die Verwaltung und Überwachung des Voice Agent Systems.

\section{Tests und Vorbereitung}

\subsection{Container-Status prüfen}

\begin{lstlisting}[language=bash,caption={Status aller Docker-Container anzeigen}]
docker-compose ps
\end{lstlisting}

\noindent Alle fünf Container sollten im Status ``Up'' sein.

\subsection{LiveKit-Konnektivität testen}

Öffnen Sie im Browser \url{http://localhost:3000} und klicken Sie auf ``Voice Agent starten''. Die Verbindungsstatusanzeige sollte von ``Getrennt'' über ``Verbindet...'' auf ``Verbunden'' wechseln (grüner Indikator). Dies ist nur notwendig, wenn das System über den Browser verwendet wird.

\subsection{SIP-Telefonie testen}

Rufen Sie die konfigurierte DID-Nummer von einem externen Telefon an. Der Anruf sollte automatisch angenommen werden und der Agent Worker sollte in den Logs ein neues Room-Event anzeigen. Dafür muss der Agent Worker gestartet sein.
\clearpage
\section{Logs und Debugging}
\label{sec:anhang-logs}

\subsection{Container-Logs anzeigen}

\begin{lstlisting}[language=bash,caption={Logs für einzelne Container anzeigen}]
# LiveKit Server
docker-compose logs -f livekit

# LiveKit SIP-Bridge
docker-compose logs -f livekit-sip

# Asterisk
docker-compose logs -f asterisk
\end{lstlisting}

\subsection{Agent Worker Logs}

Der Agent Worker schreibt alle Events in \texttt{agent\_worker\_debug.log}:

\begin{lstlisting}[language=bash,caption={Agent Worker Logs in Echtzeit anzeigen}]
tail -f agent_worker_debug.log
\end{lstlisting}

\subsection{Mögliche Probleme durch fehlende Konfigurationswerte}

\subsubsection{Problem: \texttt{LIVEKIT\_API\_KEY} ist nicht definiert}

\begin{itemize}
  \item \textbf{Lösung:} Prüfen Sie, ob \texttt{.env.local} existiert und korrekt formatiert ist (keine Leerzeichen um das ``=''-Zeichen).
\end{itemize}

\subsubsection{Problem: Eingehende Anrufe kommen nicht an}

\begin{itemize}
\item \textbf{Lösung:} Prüfen Sie die Router-Portfreigaben (Port 5060/\gls{gls-UDP}, 10200\textendash10299/\gls{gls-UDP})
  \item Prüfen Sie die \gls{gls-SIP}-Registrierung: \texttt{docker exec asterisk asterisk -rx "pjsip show registrations"}
\end{itemize}

\subsubsection{Problem: Agent antwortet nicht}

\begin{itemize}
  \item \textbf{Lösung:} Prüfen Sie den OpenAI \gls{gls-API}-Key in \texttt{.env.local}.
  \item Prüfen Sie Agent Worker Logs auf Fehler: \texttt{grep ERROR agent\_worker\_debug.log}
\end{itemize}

\section{System beenden}
\label{sec:anhang-shutdown}

\begin{lstlisting}[language=bash,caption={Alle Komponenten beenden}]
# Next.js beenden (Strg+C im Terminal)
# Agent Worker beenden (Strg+C im Terminal)
# Docker-Container stoppen
docker-compose down
\end{lstlisting}

\chapter{Konfigurationsdateien}
\markboth{ANHANG}{ANHANG}
\label{ch:anhang-config}

Dieses Kapitel enthält die vollständigen Konfigurationsdateien des implementierten Systems zur Nachvollziehbarkeit und Reproduzierbarkeit. Standort- und personenbezogene Daten wurden maskiert oder durch Platzhalter ersetzt.

\section{Docker-Compose Orchestrierung}
\label{sec:anhang-docker-compose}

Die vollständige \texttt{docker-compose.yml} orchestriert alle fünf Mikroservices mit expliziten IP-Adressen und Port-Mappings:

\begin{lstlisting}[language=yaml,caption={docker-compose.yml - Service-Orchestrierung}]
services:
  # Asterisk SIP Bridge
  asterisk:
    image: andrius/asterisk:latest
    restart: unless-stopped
    ports:
      - "5060:5060/udp"
      - "10200-10299:10200-10299/udp"
    volumes:
      - ./asterisk/pjsip.conf:/etc/asterisk/pjsip.conf:ro
      - ./asterisk/extensions.conf:/etc/asterisk/extensions.conf:ro
      - ./asterisk/modules.conf:/etc/asterisk/modules.conf:ro
    networks:
      livekit-network:
        ipv4_address: 172.20.0.10

  # Redis - Session State
  redis:
    image: redis:7-alpine
    restart: unless-stopped
    ports:
      - "6379:6379"
    networks:
      livekit-network:
        ipv4_address: 172.20.0.2

  # LiveKit Server
  livekit:
    image: livekit/livekit-server:latest
    command: --config /etc/livekit-config.yaml
    restart: unless-stopped
    ports:
      - "7880:7880"
      - "7881:7881"
      - "7882:7882/udp"
      - "50000-50100:50000-50100/udp"
    environment:
      LIVEKIT_KEYS: "<api-key>: <api-secret>"
    volumes:
      - ./livekit-config.yaml:/etc/livekit-config.yaml:ro
    networks:
      livekit-network:
        ipv4_address: 172.20.0.3

  # LiveKit SIP Bridge
  livekit-sip:
    image: livekit/sip:latest
    restart: unless-stopped
    ports:
      - "5069:5060/udp"
      - "10000-10100:10000-10100/udp"
    environment:
      - LIVEKIT_URL=ws://livekit:7880
      - LIVEKIT_API_KEY=<api-key>
      - LIVEKIT_API_SECRET=<api-secret>
      - REDIS_HOST=redis
      - REDIS_PORT=6379
    volumes:
      - ./sip-config.yaml:/sip/config.yaml:ro
    networks:
      livekit-network:
        ipv4_address: 172.20.0.5

networks:
  livekit-network:
    driver: bridge
    ipam:
      config:
        - subnet: 172.20.0.0/16
\end{lstlisting}

\section{Asterisk-PBX Konfiguration}
\label{sec:anhang-asterisk}
Im folgenden werden die Konfigurationsdateien für die Asterisk-\gls{gls-PBX} vorgestellt.
\subsection{PJSIP Transport und Trunks}

Die \texttt{pjsip.conf} definiert \gls{gls-SIP}-Transports, Provider-Registrierung und Endpoints:

\begin{lstlisting}[language=yaml,caption={pjsip.conf - SIP-Konfiguration}]
[transport-udp]
type=transport
protocol=udp
bind=0.0.0.0:5060
external_media_address=<external-ip>
external_signaling_address=<external-ip>

[provider]
type=registration
outbound_auth=provider-auth
server_uri=sip:<sip-provider-domain>
client_uri=sip:<sip-username>@<sip-provider-domain>
retry_interval=60

[provider-auth]
type=auth
auth_type=userpass
username=<sip-username>
password=<sip-password>

[provider-endpoint]
type=endpoint
context=from-provider
disallow=all
allow=ulaw,alaw
rtp_symmetric=yes
force_rport=yes
rewrite_contact=yes
from_user=<sip-username>
outbound_auth=provider-auth

[provider-endpoint]
type=aor
contact=sip:<sip-provider-domain>

[provider-endpoint]
type=identify
endpoint=provider-endpoint
match=<sip-provider-domain>
\end{lstlisting}

\subsection{Dialplan und Call-Routing}

Die \texttt{extensions.conf} enthält die Routing-Logik für eingehende Anrufe:

\begin{lstlisting}[language=ini,caption={extensions.conf - Dialplan}]
[from-provider]
exten => <your-did-number>,1,NoOp(Incoming call from SIP Provider)
 same => n,Answer()
 same => n,Stasis(livekit-bridge)
 same => n,Hangup()
\end{lstlisting}

\subsection{Modul-Loading}

Die \texttt{modules.conf} definiert, welche Asterisk-Module geladen werden:

\begin{lstlisting}[language=ini,caption={modules.conf - Modulkonfiguration}]
[modules]
autoload=yes

noload => chan_sip.so
load => res_pjsip.so
load => res_pjsip_session.so
load => res_pjsip_transport_websocket.so
load => res_ari.so
load => res_stasis.so
\end{lstlisting}

\section{LiveKit Server Konfiguration}
\label{sec:anhang-livekit}
Im folgenden werden die Konfigurationsdateien für die LiveKit Komponenten vorgestellt.
\subsection{Media Server}

Die \texttt{livekit-config.yaml} konfiguriert den zentralen \gls{gls-WebRTC} Media Server:

\begin{lstlisting}[language=yaml,caption={livekit-config.yaml - Media Server}]
port: 7880
rtc:
  port_range_start: 50000
  port_range_end: 50100
  use_external_ip: true
keys:
  <api-key>: <api-secret>
room:
  auto_create: true
  empty_timeout: 300
  max_participants: 10
\end{lstlisting}

\subsection{SIP Bridge}

Die \texttt{sip-config.yaml} konfiguriert das \gls{gls-SIP}-zu-\gls{gls-WebRTC} Gateway:

\begin{lstlisting}[language=yaml,caption={sip-config.yaml - SIP Bridge}]
logging:
  level: debug
sip:
  port: 5060
  local_net:
    - 172.20.0.0/16
redis:
  address: redis:6379
livekit:
  url: ws://livekit:7880
  api_key: <api-key>
  api_secret: <api-secret>
\end{lstlisting}




\printbibliography
% Eidesstattliche Versicherung / Selbständigkeitserklärung
% Nach HTWK Leipzig Vorgaben

\chapter*{Eidesstattliche Versicherung}
\addcontentsline{toc}{chapter}{Eidesstattliche Versicherung}

Ich erkläre hiermit, dass ich die vorliegende Arbeit selbstständig ohne Hilfe Dritter und ohne Benutzung anderer als der angegebenen Quellen und Hilfsmittel verfasst habe. Alle den benutzten Quellen wörtlich oder sinngemä\ss{} entnommenen Stellen sind als solche einzeln kenntlich gemacht.

Diese Arbeit ist bislang keiner anderen Prüfungsbehörde vorgelegt und auch nicht veröffentlicht worden.

Ich bin mir bewusst, dass eine falsche Erklärung rechtliche Folgen haben wird.

\vspace{2cm}

\noindent
\begin{tabular}{@{}l@{\hspace{4cm}}l@{}}
Ort, Datum & Unterschrift \\
\\
\line(1,0){60} & \line(1,0){60}
\end{tabular}

\end{document}